% El consumo responsable y la ética profesional (continuación) : expresión escrita — Espagnol 1ère A — Cours signé OnBuch+
% Programme officiel MINESEC. Compiler avec : tectonic EA16-consumo-responsable-etica-continuacion.tex
\newcommand{\DOCMATIERE}{Espagnol}
\newcommand{\DOCNIVEAU}{1ère A}
\newcommand{\DOCMODULE}{Módulo 4 : Vie économique}
\newcommand{\DOCLECON}{EA16}
\newcommand{\DOCTITRE}{El consumo responsable y la ética profesional (continuación) : expresión escrita}
% OnBuch+ — préambule « cours signés » ENRICHI (compilé avec tectonic / XeLaTeX)
% Système visuel « Pop » d'OnBuch+ : fond crème (jamais blanc), bordures encre
% épaisses, ombres « dures » décalées, titres Archivo Black en capitales.
% Version MATHÉMATIQUES : graphes (pgfplots/tikz), boîtes pédagogiques
% (définition, propriété, méthode, attention, le savais-tu, exercice résolu,
% exercice, TP, bilan), page de garde et signature OnBuch+ ancrée en bas.
%
% Macros attendues dans main.tex AVANT \input{preamble} :
%   \DOCMATIERE  \DOCNIVEAU  \DOCMODULE  \DOCLECON (numéro)  \DOCTITRE
\documentclass[11pt]{article}

\usepackage{fontspec}
\usepackage[spanish,french]{babel} % français = langue principale ; l'espagnol via \esp{..} / espagnol
\usepackage[a4paper,margin=2cm,top=3.4cm,bottom=2.8cm,headheight=1.4cm,footskip=1.3cm]{geometry}
\usepackage{amsmath,amssymb,amsfonts}
\usepackage{mathtools} % \xmapsto, \coloneqq... souvent utilisés par les modèles
\usepackage{cancel} % simplifications barrées (bilans de forces)
% \abs{z} (module, valeur absolue) et \norm{u} (norme), utilisés par les modèles
\providecommand{\abs}{}\let\abs\relax\DeclarePairedDelimiter{\abs}{\lvert}{\rvert}
\providecommand{\norm}{}\let\norm\relax\DeclarePairedDelimiter{\norm}{\lVert}{\rVert}
\usepackage[table]{xcolor} % [table] : fournit aussi \rowcolors (alternance de lignes)
\usepackage{pagecolor}
\usepackage{tikz}
\usetikzlibrary{arrows.meta,positioning,calc,shapes.geometric,shapes.misc,shapes.arrows,decorations.pathmorphing,decorations.pathreplacing,decorations.markings,patterns,shadows,mindmap,trees,fit,backgrounds,matrix,intersections,angles,quotes}
\usepackage{pgfplots}
\usepackage[version=4]{mhchem} % équations nucléaires \ce{^{235}_{92}U}
\usepackage[european,siunitx]{circuitikz} % schémas électriques
\pgfplotsset{compat=1.18}
\usepgfplotslibrary{fillbetween,groupplots} % aires entre courbes (calcul intégral)
\usepackage{enumitem}
\usepackage{fancyhdr}
% [most] charge toutes les bibliothèques tcolorbox (listings, minted, raster,
% documentation...) et gonfle inutilement la mémoire TeX sur un document long
% (~300 boîtes) : on ne charge que ce qui est réellement utilisé.
\usepackage[skins,breakable]{tcolorbox}
\usepackage{graphicx}
\usepackage{colortbl}
\usepackage{booktabs}
\usepackage{tabularx}
\usepackage{diagbox} % case d'en-tête diagonale des tableaux à double entrée
% Colonnes extensibles centrée / gauche / droite (C, L, R), souvent utilisées par les modèles
\newcolumntype{C}{>{\centering\arraybackslash}X}
\newcolumntype{L}{>{\raggedright\arraybackslash}X}
\newcolumntype{R}{>{\raggedleft\arraybackslash}X}
\usepackage{array}
\usepackage{multicol}
\usepackage{siunitx}
% parse-numbers=false : accepte aussi \SI{2,5 \times 10^{-3}}{...}
\sisetup{locale=FR,output-decimal-marker={,},group-minimum-digits=4,parse-numbers=false}
\DeclareSIUnit\ohm{\ensuremath{\symup{\Omega}}} % U+2126 absent de la police
\DeclareSIUnit\division{div} % réglages d'oscilloscope (ms/div, V/div)
\DeclareSIUnit\tr{tr} % tours (vitesse de rotation, ex. \SI{1500}{\tr\per\minute})
\DeclareSIUnit\molar{M} % concentration molaire (ex. \SI{3}{\molar}, \SI{1}{\micro\molar})
\DeclareSIUnit\an{an} % année (le modèle confond parfois avec \year, un compteur TeX) — ex. \SI{5}{\kilo\meter\per\an}
\DeclareSIUnit\FCFA{FCFA} % franc CFA (ex. \SI{500}{\FCFA\per\kilo\gram})
\DeclareSIUnit\degreeBrix{\textdegree Brix} % teneur en sucre d'un sirop/jus (réfractométrie)
\DeclareSIUnit\month{mois} % le modèle utilise parfois \month, un compteur TeX, comme unité de durée
\usepackage{qrcode}
% \degree : commande fréquemment utilisée par les modèles pour un angle ou une
% température en dehors de \SI{}{} (non fournie par les paquets chargés).
\providecommand{\degree}{^{\circ}}
% \qed (fin de démonstration), utilisé par les modèles sans amsthm
\providecommand{\qed}{\hfill$\square$}
% \barre : double barre des valeurs interdites dans un tableau de variations (tkz-tab)
\providecommand{\barre}{\ensuremath{\|}}
% environnement proof (amsthm non chargé) utilisé par les modèles
\ifdefined\proof\else\newenvironment{proof}[1][Démonstration]{\par\noindent\textit{#1.}\ }{\qed\par}\fi
\usepackage{lastpage}
\usepackage{hyperref}
\hypersetup{hidelinks}

% ============================================================
% Palette « Pop » OnBuch+ — mêmes hex que la classe P (plus_ui.dart)
% ============================================================
\definecolor{poporange}{HTML}{F59321}
\definecolor{popdark}{HTML}{E07A0C}
\definecolor{popcream}{HTML}{FFF6E8}
\definecolor{popink}{HTML}{1C1714}
\definecolor{poppurple}{HTML}{7A5AE0}
\definecolor{popgreen}{HTML}{2E9E6A}
\definecolor{poppink}{HTML}{E0457B}
\definecolor{popblue}{HTML}{2F7DE1}
\definecolor{popgold}{HTML}{C98A12}
\definecolor{popmuted}{HTML}{8A7F76}
\definecolor{popline}{HTML}{EADFD2}
\definecolor{popgray}{HTML}{8A7F76}
\colorlet{popgrey}{popgray}
% Alias tolérants (le modèle nomme parfois les couleurs différemment)
\colorlet{popwhite}{white}
\colorlet{popblack}{popink}
\colorlet{popred}{poppink}
\colorlet{popyellow}{popgold}
% Teintes claires (fonds de tableaux/encadrés) — le modèle nomme parfois
% les couleurs avec un suffixe L (ex. popblueL) sans les avoir définies.
\colorlet{poporangeL}{poporange!20}
\colorlet{popdarkL}{popdark!20}
\colorlet{poppurpleL}{poppurple!20}
\colorlet{popgreenL}{popgreen!20}
\colorlet{poppinkL}{poppink!20}
\colorlet{popblueL}{popblue!20}
\colorlet{popgoldL}{popgold!20}
\colorlet{popredL}{popred!20}
\colorlet{popgrayL}{popgray!20}
% Teintes claires pour fonds de tableaux / zones
\colorlet{popblueL}{popblue!10!white}
\colorlet{poporangeL}{poporange!14!white}
\colorlet{popgreenL}{popgreen!12!white}
\colorlet{poppurpleL}{poppurple!12!white}
\colorlet{poppinkL}{poppink!12!white}
\colorlet{popgoldL}{popgold!14!white}

\pagecolor{popcream}

% ============================================================
% Typographie
% ============================================================
\defaultfontfeatures{Ligatures=TeX}
\setmainfont{PlusJakartaSans-Medium.ttf}[Path=../../../fonts/,BoldFont=PlusJakartaSans-Bold.ttf]
% Maths en sans-serif (Fira Math) avec chiffres et lettres droites de Plus
% Jakarta, pour que les formules se fondent dans le texte.
\usepackage[math-style=ISO,bold-style=ISO]{unicode-math}
\setmathfont{FiraMath-Regular.otf}[Path=../../../fonts/]
\setmathfont{PlusJakartaSans-Medium.ttf}[Path=../../../fonts/,range={up/num,up/{Latin,latin},\mathup}]
\usepackage{icomma} % virgule décimale sans espace en mode math
% Caractères mathématiques Unicode absents des polices de texte (Plus Jakarta,
% Archivo Black) : rendus par leur équivalent mathématique, en texte comme en
% formule — sinon ils s'affichent en carré vide (ex. « Arithmétique dans ℤ »).
\usepackage{newunicodechar}
\newunicodechar{ℝ}{\ensuremath{\mathbb{R}}}
\newunicodechar{ℤ}{\ensuremath{\mathbb{Z}}}
\newunicodechar{ℂ}{\ensuremath{\mathbb{C}}}
\newunicodechar{ℕ}{\ensuremath{\mathbb{N}}}
\newunicodechar{ℚ}{\ensuremath{\mathbb{Q}}}
\newunicodechar{𝔻}{\ensuremath{\mathbb{D}}}
\newunicodechar{α}{\ensuremath{\alpha}}
\newunicodechar{β}{\ensuremath{\beta}}
\newunicodechar{γ}{\ensuremath{\gamma}}
\newunicodechar{δ}{\ensuremath{\delta}}
\newunicodechar{ε}{\ensuremath{\varepsilon}}
\newunicodechar{θ}{\ensuremath{\theta}}
\newunicodechar{λ}{\ensuremath{\lambda}}
\newunicodechar{μ}{\ensuremath{\mu}}
\newunicodechar{σ}{\ensuremath{\sigma}}
\newunicodechar{τ}{\ensuremath{\tau}}
\newunicodechar{φ}{\ensuremath{\varphi}}
\newunicodechar{ψ}{\ensuremath{\psi}}
\newunicodechar{ω}{\ensuremath{\omega}}
\newunicodechar{Σ}{\ensuremath{\Sigma}}
% majuscules produites par \MakeUppercase dans les titres (α-aminés -> Α)
\newunicodechar{Α}{\ensuremath{\alpha}}
\newunicodechar{Β}{\ensuremath{\beta}}
\newunicodechar{Γ}{\ensuremath{\Gamma}}
\newunicodechar{Δ}{\ensuremath{\Delta}}
\newunicodechar{Ε}{\ensuremath{\varepsilon}}
\newunicodechar{Θ}{\ensuremath{\Theta}}
\newunicodechar{Λ}{\ensuremath{\Lambda}}
\newunicodechar{Μ}{\ensuremath{\mu}}
\newunicodechar{Τ}{\ensuremath{\tau}}
\newunicodechar{Φ}{\ensuremath{\Phi}}
\newunicodechar{Ψ}{\ensuremath{\Psi}}
\newunicodechar{Ω}{\ensuremath{\Omega}}
\newunicodechar{↦}{\ensuremath{\mapsto}}
\newunicodechar{∈}{\ensuremath{\in}}
\newunicodechar{∉}{\ensuremath{\notin}}
\newunicodechar{∧}{\ensuremath{\wedge}}
\newunicodechar{⇔}{\ensuremath{\Leftrightarrow}}
\newunicodechar{⟺}{\ensuremath{\Longleftrightarrow}}
\newunicodechar{⇒}{\ensuremath{\Rightarrow}}
\newunicodechar{⟹}{\ensuremath{\Longrightarrow}}
\newunicodechar{∘}{\ensuremath{\circ}}
\newunicodechar{∪}{\ensuremath{\cup}}
\newunicodechar{∩}{\ensuremath{\cap}}
\newunicodechar{≡}{\ensuremath{\equiv}}
\newunicodechar{⊂}{\ensuremath{\subset}}
\newunicodechar{⊥}{\ensuremath{\perp}}
\newunicodechar{∃}{\ensuremath{\exists}}
\newunicodechar{∀}{\ensuremath{\forall}}
\newunicodechar{∛}{\ensuremath{\sqrt[3]{\,}}}
\newunicodechar{∜}{\ensuremath{\sqrt[4]{\,}}}
\newunicodechar{⃗}{\ensuremath{{}^{\to}}}
\newunicodechar{ⁿ}{\ifmmode^{n}\else\textsuperscript{\ensuremath{n}}\fi}
\newunicodechar{ˣ}{\ifmmode^{x}\else\textsuperscript{\ensuremath{x}}\fi}
\newunicodechar{ᵇ}{\ifmmode^{b}\else\textsuperscript{\ensuremath{b}}\fi}
\newunicodechar{ʳ}{\ifmmode^{r}\else\textsuperscript{\ensuremath{r}}\fi}
\newunicodechar{ᵘ}{\ifmmode^{u}\else\textsuperscript{\ensuremath{u}}\fi}
\newunicodechar{ʸ}{\ifmmode^{y}\else\textsuperscript{\ensuremath{y}}\fi}
\newunicodechar{ᵃ}{\ifmmode^{a}\else\textsuperscript{\ensuremath{a}}\fi}
\newunicodechar{ᵉ}{\ifmmode^{e}\else\textsuperscript{\ensuremath{e}}\fi}
\newunicodechar{ᵏ}{\ifmmode^{k}\else\textsuperscript{\ensuremath{k}}\fi}
\newunicodechar{ᵖ}{\ifmmode^{p}\else\textsuperscript{\ensuremath{p}}\fi}
\newunicodechar{ᴺ}{\ifmmode^{N}\else\textsuperscript{\ensuremath{N}}\fi}
\newunicodechar{ᶜ}{\ifmmode^{c}\else\textsuperscript{\ensuremath{c}}\fi}
\newunicodechar{ʰ}{\ifmmode^{h}\else\textsuperscript{\ensuremath{h}}\fi}
\newunicodechar{ᵠ}{\ifmmode^{q}\else\textsuperscript{\ensuremath{q}}\fi}
\newunicodechar{ₙ}{\ifmmode_{n}\else\textsubscript{\ensuremath{n}}\fi}
\newunicodechar{ₐ}{\ifmmode_{a}\else\textsubscript{\ensuremath{a}}\fi}
\newunicodechar{ᵢ}{\ifmmode_{i}\else\textsubscript{\ensuremath{i}}\fi}
\newunicodechar{ᵣ}{\ifmmode_{r}\else\textsubscript{\ensuremath{r}}\fi}
\newunicodechar{ₖ}{\ifmmode_{k}\else\textsubscript{\ensuremath{k}}\fi}
\newunicodechar{ᵦ}{\ifmmode_{\beta}\else\textsubscript{\ensuremath{\beta}}\fi}
\newunicodechar{ₓ}{\ifmmode_{x}\else\textsubscript{\ensuremath{x}}\fi}
\newfontfamily\popdisplay{ArchivoBlack-Regular.ttf}[Path=../../../fonts/]
\newfontfamily\poplogo{SpaceGrotesk-Bold.ttf}[Path=../../../fonts/]
\newfontfamily\popsemi{PlusJakartaSans-SemiBold.ttf}[Path=../../../fonts/]
\newfontfamily\popxbold{PlusJakartaSans-ExtraBold.ttf}[Path=../../../fonts/]
\newfontfamily\popxboldls{PlusJakartaSans-ExtraBold.ttf}[Path=../../../fonts/,LetterSpace=3.0]

\setlist[enumerate]{leftmargin=1.7em,itemsep=5pt,topsep=4pt}
\setlist[itemize]{leftmargin=1.7em,itemsep=4pt,topsep=4pt,label={\color{poporange}\textbullet}}
\setlength{\parindent}{0pt}
\setlength{\parskip}{5pt}
\renewcommand{\arraystretch}{1.25}

% pgfplots : style Pop par défaut
\pgfplotsset{
  every axis/.append style={axis line style={popink,line width=1pt},tick style={popink},
    grid style={popline,line width=0.5pt},label style={font=\small},tick label style={font=\footnotesize}},
  popcurve/.style={poporange,line width=1.8pt,no markers,smooth},
}
% Même style disponible dans un \draw TikZ simple (hors pgfplots)
\tikzset{popcurve/.style={poporange,line width=1.8pt}}
\tikzset{popgrid/.style={popline,very thin}}
\colorlet{popcurve}{poporange} % le nom du style est parfois utilisé comme couleur

% ============================================================
% Pill / squiggle
% ============================================================
\newcommand{\poppill}[3]{%
  \tikz[baseline=-0.55ex]\node[draw=popink,line width=1.2pt,rounded corners=2.6pt,%
    fill=#2,inner xsep=6.5pt,inner ysep=3pt]{%
    \popxboldls\fontsize{7}{7}\selectfont\color{#3}\MakeUppercase{#1}};%
}
\newcommand{\popsquiggle}[1]{%
  \tikz[baseline=-0.4ex]{
    \draw[#1,line width=2.6pt,line cap=round]
      plot [smooth] coordinates {(0,0.09) (0.34,0.01) (0.68,0.21) (1.02,0.01) (1.36,0.10)};
  }%
}
% Mot-clé surligné (marqueur orange sous le texte)
\newcommand{\cle}[1]{{\popxbold\color{popdark}#1}}

% ============================================================
% En-tête / pied
% ============================================================
\pagestyle{fancy}
\fancyhf{}
\renewcommand{\headrulewidth}{0pt}
\renewcommand{\footrulewidth}{0pt}
\lhead{\raisebox{-0.15ex}{\poplogo\fontsize{13}{13}\selectfont\color{popink}OnBuch\color{poporange}+}}
\rhead{\poppill{\DOCMATIERE\ \textbullet\ \DOCNIVEAU\ \textbullet\ Leçon \DOCLECON}{white}{popink}}
\lfoot{\popxbold\fontsize{7.6}{7.6}\selectfont\color{popmuted}\MakeUppercase{Cours signé OnBuch+}\ \textbullet\ {\popsemi\DOCTITRE}}
\rfoot{\tikz[baseline=-0.6ex]\node[rounded corners=8.5pt,fill=popink,minimum height=17pt,minimum width=17pt,inner xsep=5pt,inner sep=0]{\popxbold\fontsize{8}{8}\selectfont\color{popcream}\thepage\,/\,\pageref*{LastPage}};}

% ============================================================
% Titres
% ============================================================
\newcommand{\chaptitre}[2]{%
  \begin{tcolorbox}[enhanced,colback=poporange,colframe=popink,boxrule=2.6pt,arc=11pt,
      drop shadow=popink,left=15pt,right=15pt,top=12pt,bottom=11pt,before skip=3pt,after skip=18pt]
    \poppill{#2}{popink}{popcream}\par\vspace{9pt}
    {\raggedright\hyphenpenalty=10000\exhyphenpenalty=10000\popdisplay\fontsize{22}{24}\selectfont\color{popcream}\MakeUppercase{#1}\par}
  \end{tcolorbox}
}
% Section numérotée : pastille orange avec le numéro + titre Archivo
\newcounter{popsec}
\newcommand{\coursec}[1]{%
  \stepcounter{popsec}\setcounter{popsub}{0}%
  \par\vspace{16pt}\noindent
  \begin{minipage}{\linewidth}
  \tikz[baseline=-0.7ex]\node[draw=popink,line width=1.6pt,fill=poporange,rounded corners=4pt,minimum size=22pt,inner sep=0]{\popdisplay\fontsize{12}{12}\selectfont\color{popcream}\thepopsec};%
  \hspace{8pt}{\popdisplay\fontsize{15}{17}\selectfont\color{popink}\MakeUppercase{#1}}\par
  \vspace{4pt}\noindent{\color{poporange}\rule{36pt}{3.6pt}}
  \end{minipage}\par\nopagebreak\vspace{8pt}
}
\newcounter{popsub}
\newcommand{\courssub}[1]{%
  \stepcounter{popsub}%
  \par\vspace{9pt}\noindent{\popxbold\fontsize{11.5}{13}\selectfont\color{popink}{\color{poporange}\thepopsec.\thepopsub}\hspace{6pt}#1}\par\nopagebreak\vspace{4pt}
}

% ============================================================
% Boîtes pédagogiques — même recette : fond blanc, bordure épaisse colorée,
% ombre dure encre, badge plein. Titre optionnel [#1].
% ============================================================
\tcbset{popbase/.style={breakable,enhanced,colback=white,boxrule=2.3pt,arc=9pt,
  shadow={3pt}{-3pt}{0pt}{popink},left=14pt,right=14pt,top=11pt,bottom=11pt,before skip=11pt,after skip=15pt,
  fontupper=\popsemi\fontsize{10.6}{14.5}\selectfont\color{popink},
  fontlower=\popsemi\fontsize{10.6}{14.5}\selectfont\color{popink}}}
% Coche dessinée (le glyphe ✓ n'existe pas dans les polices du document)
\newcommand{\popcheck}[1]{\tikz[baseline=-0.5ex]\draw[#1,line width=1.6pt,line cap=round,line join=round](-0.13,0.0)--(-0.04,-0.09)--(0.13,0.1);}
\providecommand{\coche}{\popcheck{popgreen}}
\providecommand{\resultat}[1]{\par\smallskip\noindent\fcolorbox{popgreen}{popgreenL}{\parbox{\dimexpr\linewidth-2\fboxsep-2\fboxrule\relax}{\textbf{Résultat.} #1}}\par\smallskip}
\newcommand{\popboxhead}[3]{\poppill{#1}{#2}{white}\ifx\relax#3\relax\else\hspace{6pt}{\popxbold\fontsize{10.6}{12}\selectfont\color{popink}#3}\fi\par\vspace{7pt}}

\newtcolorbox{aretenir}[1][]{popbase,colframe=popgreen,before upper={\popboxhead{À retenir}{popgreen}{#1}}}
% Texte en espagnol (document de lecture, dialogue, poème, extrait) : cadre
% d'étude. Un titre optionnel (source, auteur) s'affiche dans l'en-tête.
\newtcolorbox{textebox}[1][]{popbase,colframe=popblue,before upper={\popboxhead{Texto}{popblue}{#1}\begin{otherlanguage}{spanish}},after upper={\end{otherlanguage}}}
% Marques ✓ / ✗ (forme correcte / fautive) souvent utilisées par les modèles
\providecommand{\cross}{\textcolor{poppink}{\ensuremath{\times}}}
\providecommand{\xmark}{\cross}\providecommand{\crossmark}{\cross}\providecommand{\cmark}{\ensuremath{\checkmark}}
% Ligne de réponse / justification à compléter (inventée par les modèles)
\providecommand{\justification}[1]{\par\noindent\textit{Justification~:} \dotfill\par}
% \textipa (package tipa non chargé) : la police ne couvre pas l'API -> simple passage
\providecommand{\textipa}[1]{#1}
% Soulignement (ulem non chargé) : \uline, \ul
\providecommand{\uline}[1]{\underline{#1}}\providecommand{\ul}[1]{\underline{#1}}
% Ordinaux français écrits en macro par les modèles : 1\ière, 3\ième
\newcommand{\ière}{\textsuperscript{re}}\newcommand{\ième}{\textsuperscript{e}}
% \exemple (inventée par les modèles) : étiquette « Exemple : »
\providecommand{\exemple}{\textbf{Exemple~:}\ }
% Mot ou phrase en espagnol à l'intérieur d'un paragraphe français
\newcommand{\esp}[1]{\ifmmode\text{\textit{#1}}\else\begingroup\selectlanguage{spanish}\itshape #1\endgroup\fi}
\newtcolorbox{exemplebox}[1][]{popbase,colframe=poporange,before upper={\popboxhead{Exemple}{poporange}{#1}}}
\newtcolorbox{definition}[1][]{popbase,colframe=popblue,before upper={\popboxhead{Définition}{popblue}{#1}}}
\newtcolorbox{propriete}[1][]{popbase,colframe=poppurple,before upper={\popboxhead{Propriété}{poppurple}{#1}}}
\newtcolorbox{methode}[1][]{popbase,colframe=popgold,before upper={\popboxhead{Méthode}{popgold}{#1}}}
\newtcolorbox{attention}[1][]{popbase,colframe=poppink,before upper={\popboxhead{Attention}{poppink}{#1}}}
\newtcolorbox{experience}[1][]{popbase,colframe=popblue,colback=popblueL,before upper={\popboxhead{Expérience}{popblue}{#1}}}
% « Le savais-tu ? » : carte violette pleine (contexte, culture, Cameroun)
\newtcolorbox{savaistu}[1][]{popbase,colback=poppurple,colframe=popink,
  fontupper=\popsemi\fontsize{10.4}{14.2}\selectfont\color{white},
  before upper={\poppill{Le savais-tu\,?}{white}{poppurple}\ifx\relax#1\relax\else\hspace{6pt}{\popxbold\color{white}#1}\fi\par\vspace{7pt}}}
% Exercice résolu : énoncé en haut, \tcblower puis la solution sur fond crème
\newcounter{popexo}\newcounter{popexr}
\newtcolorbox{exoresolu}[1][]{popbase,colframe=poporange,colbacklower=popcream,
  segmentation style={popink,line width=1.2pt,dashed},
  before upper={\stepcounter{popexr}\popboxhead{Exercice résolu \thepopexr}{poporange}{#1}},
  before lower={\poppill{Solution}{popink}{popcream}\par\vspace{6pt}}}
% Exercice (non résolu en place) — la correction est dans la partie Corrigés
\newtcolorbox{exercice}[1][]{popbase,colframe=popink,
  before upper={\stepcounter{popexo}\popboxhead{Exercice \thepopexo}{popink}{#1}}}
% Corrigé
\newtcolorbox{corrige}[1][]{popbase,colframe=popgreen,colback=popgreenL,
  before upper={\popboxhead{Corrigé}{popgreen}{#1}}}
% Objectifs / prérequis / bilan
\newtcolorbox{objectifs}{popbase,colframe=popink,colback=poporangeL,
  before upper={\popboxhead{Objectifs de la leçon}{popink}{}}}
\newtcolorbox{prerequis}{popbase,colframe=popink,colback=popblueL,
  before upper={\popboxhead{Prérequis}{popblue}{}}}
\newtcolorbox{bilan}{popbase,colframe=popink,colback=white,boxrule=2.8pt,
  before upper={\popboxhead{Fiche bilan}{poporange}{L'essentiel à connaître pour l'examen}}}
% Figure encadrée avec légende
\newtcolorbox{popfigure}[1][]{enhanced,breakable=false,colback=white,colframe=popline,boxrule=1.4pt,arc=8pt,
  left=10pt,right=10pt,top=10pt,bottom=8pt,before skip=10pt,after skip=12pt,halign=center,
  lower separated=false,halign lower=center,
  fontlower=\popsemi\fontsize{9}{11.5}\selectfont\color{popmuted},
  #1}
\newcommand{\legende}[1]{\par\vspace{6pt}{\popsemi\fontsize{9}{11.5}\selectfont\color{popmuted}#1}}

% ============================================================
% Signature OnBuch+
% ============================================================
\newcommand{\poplogoinline}[1]{{\poplogo\fontsize{#1}{#1}\selectfont\color{popink}OnBuch\color{poporange}+}}
\newcommand{\signatureonbuch}{%
  \begin{tcolorbox}[enhanced,colback=popink,colframe=popink,boxrule=2.2pt,arc=10pt,
      drop shadow=poporange,left=14pt,right=14pt,top=10pt,bottom=10pt,before skip=0pt,after skip=0pt]
    \tikz[baseline=-0.7ex]\node[circle,fill=popgreen,draw=popcream,line width=1.3pt,minimum size=22pt,inner sep=0]{\popcheck{white}};%
    \hspace{9pt}\begin{minipage}[c]{0.62\linewidth}
      {\popxbold\fontsize{12}{13}\selectfont\color{popcream}Signé\ \textbullet\ {\poplogo OnBuch\color{poporange}+}}\par\vspace{2pt}
      {\raggedright\popsemi\fontsize{8}{10.5}\selectfont\color{popline}\MakeUppercase{Équipe pédagogique OnBuch+}\\\MakeUppercase{Conforme au programme officiel MINESEC}\par}
    \end{minipage}\hfill
    {\poplogo\fontsize{22}{22}\selectfont\color{popcream}OnBuch\color{poporange}+}
  \end{tcolorbox}%
}

% ============================================================
% Page de garde
% #1 titre · #2 matière · #3 niveau · #4 module · #5 n° leçon · #6 durée ·
% #7 description / objectifs (texte ou itemize)
% ============================================================
\newcommand{\pagegarde}[7]{%
  \thispagestyle{empty}
  \newgeometry{a4paper,margin=1.7cm,top=1.6cm,bottom=1.4cm}%
  \begin{tikzpicture}[remember picture,overlay]
    \fill[poporange!18] ([xshift=-3.2cm,yshift=-3.6cm]current page.north east) circle (4.6cm);
    \fill[poppurple!14] ([xshift=2.4cm,yshift=7.6cm]current page.south west) circle (3.8cm);
  \end{tikzpicture}%
  {\poplogo\fontsize{30}{30}\selectfont\color{popink}OnBuch\color{poporange}+}\hfill
  \raisebox{0.6ex}{\poppill{Cours signé \textbullet\ Premium}{popink}{popcream}}\par
  \vspace{1pt}\popsquiggle{poppurple}\par
  \vspace{8pt}
  \begin{tcolorbox}[enhanced,colback=poporange,colframe=popink,boxrule=2.8pt,arc=13pt,
      drop shadow=popink,left=18pt,right=18pt,top=12pt,bottom=13pt,after skip=10pt]
    \poppill{#2 \textbullet\ #3}{popink}{popcream}\hspace{5pt}\poppill{#4}{white}{popink}\par\vspace{8pt}
    {\popdisplay\fontsize{14}{14}\selectfont\color{popink}LEÇON #5}\par\vspace{4pt}
    {\raggedright\hyphenpenalty=10000\exhyphenpenalty=10000\popdisplay\fontsize{25}{27}\selectfont\color{popcream}\MakeUppercase{#1}\par}
    \vspace{8pt}
    {\popxbold\fontsize{9.5}{9.5}\selectfont\color{popink}\MakeUppercase{Durée conseillée : #6}}
  \end{tcolorbox}
  \begin{tcolorbox}[enhanced,colback=white,colframe=popink,boxrule=2.2pt,arc=10pt,
      drop shadow=popink,left=16pt,right=16pt,top=10pt,bottom=10pt,after skip=8pt]
    \poppill{Dans cette leçon}{poppurple}{white}\par\vspace{5pt}
    \popsemi\fontsize{9.3}{12.6}\selectfont\color{popink}#7
  \end{tcolorbox}
  \noindent\begin{minipage}[t]{0.487\linewidth}
    \begin{tcolorbox}[enhanced,colback=popgreen,colframe=popink,boxrule=2.2pt,arc=10pt,
        drop shadow=popink,left=11pt,right=11pt,top=10pt,bottom=10pt,height=3.1cm,valign=center]
      \tcbox[colback=white,colframe=popink,boxrule=1.3pt,arc=5pt,boxsep=0pt,nobeforeafter,
        left=3pt,right=3pt,top=3pt,bottom=3pt,valign=center]{\qrcode[height=1.9cm]{https://play.google.com/store/apps/details?id=com.luvvix.onbuch&hl=fr}}%
      \hspace{8pt}\begin{minipage}[c]{0.5\linewidth}
        {\popdisplay\fontsize{11}{12.5}\selectfont\color{white}\MakeUppercase{Télécharge OnBuch}}\par\vspace{4pt}
        {\raggedright\popsemi\fontsize{8.4}{11}\selectfont\color{white}Gratuit \textbullet\ Google Play\\Tuteur IA, annales, concours, résultats\par}
      \end{minipage}
    \end{tcolorbox}
  \end{minipage}\hfill
  \begin{minipage}[t]{0.487\linewidth}
    \begin{tcolorbox}[enhanced,colback=poppurple,colframe=popink,boxrule=2.2pt,arc=10pt,
        drop shadow=popink,left=11pt,right=11pt,top=10pt,bottom=10pt,height=3.1cm,valign=center]
      \tikz[baseline=-0.7ex]\node[circle,fill=white,draw=popink,line width=1.3pt,minimum size=1.6cm,inner sep=0]{\popdisplay\fontsize{24}{24}\selectfont\color{poppurple}+};%
      \hspace{8pt}\begin{minipage}[c]{0.6\linewidth}
        {\popdisplay\fontsize{11}{12.5}\selectfont\color{white}\MakeUppercase{Passe à OnBuch+}}\par\vspace{4pt}
        {\raggedright\popsemi\fontsize{8.4}{11}\selectfont\color{white}Tous les cours signés, Tuteur IA illimité, fiches PDF — onglet OnBuch+ de l'app\par}
      \end{minipage}
    \end{tcolorbox}
  \end{minipage}
  \vspace{8pt}
  \popcontenu
  \vfill
  \signatureonbuch
  \restoregeometry
  \newpage
}

% Rangée « Ce cours contient » (page de garde)
\newcommand{\poptile}[3]{% #1 couleur #2 gros texte #3 légende
  \begin{tcolorbox}[enhanced,colback=#1,colframe=popink,boxrule=2pt,arc=9pt,shadow={2.5pt}{-2.5pt}{0pt}{popink},
      left=6pt,right=6pt,top=9pt,bottom=9pt,halign=center,valign=center,height=1.75cm,width=\linewidth,nobeforeafter]
    {\popdisplay\fontsize{12.5}{13}\selectfont\color{white}\MakeUppercase{#2}}\par\vspace{3pt}
    {\centering\popsemi\fontsize{7.8}{9.5}\selectfont\color{white}#3\par}
  \end{tcolorbox}}
\newcommand{\popcontenu}{%
  \par\noindent{\popxboldls\fontsize{7.5}{7.5}\selectfont\color{popmuted}CE COURS SIGNÉ CONTIENT}\par\vspace{7pt}
  \noindent\begin{minipage}[t]{0.235\linewidth}\poptile{popblue}{Cours}{complet, illustré, pas à pas}\end{minipage}\hfill
  \begin{minipage}[t]{0.235\linewidth}\poptile{poporange}{Méthodes}{et exercices résolus}\end{minipage}\hfill
  \begin{minipage}[t]{0.235\linewidth}\poptile{poppink}{Exercices}{type examen + corrigés}\end{minipage}\hfill
  \begin{minipage}[t]{0.235\linewidth}\poptile{popgreen}{Fiche bilan}{l'essentiel à retenir}\end{minipage}\par}

% Sommaire
\newtcolorbox{sommaire}{enhanced,colback=white,colframe=popink,boxrule=2.2pt,arc=10pt,
  drop shadow=popink,left=18pt,right=18pt,top=13pt,bottom=13pt,before skip=6pt,after skip=20pt,
  before upper={\poppill{Sommaire}{poppurple}{white}\par\vspace{9pt}\popsemi\fontsize{10.6}{14.5}\selectfont\color{popink}}}

% Signature de fin de cours — poussée tout en bas de la dernière page
\newcommand{\findecours}{%
  \par\vfill\nopagebreak
  \begin{center}\popsquiggle{poporange}\end{center}\vspace{4pt}
  \signatureonbuch
}
\AtBeginDocument{\def\llbracket{\mathopen{[\mkern-3.5mu[}}\def\rrbracket{\mathclose{]\mkern-3.5mu]}}} % intervalles d'entiers (glyphe ⟦ absent de la police)
% Glyphes absents de Fira Math : redéfinis à partir de glyphes disponibles
\AtBeginDocument{%
  \def\setminus{\mathbin{\backslash}}\let\smallsetminus\setminus
  \def\vdots{\mathord{\vbox{\baselineskip3.2pt\lineskiplimit0pt\kern4pt\hbox{.}\hbox{.}\hbox{.}}}}%
  \def\ddots{\mathinner{\mkern1mu\raise6.5pt\hbox{.}\mkern2mu\raise3.6pt\hbox{.}\mkern2mu\raise0.7pt\hbox{.}\mkern1mu}}%
  \def\triangle{\mathord{\scalebox{0.9}{$\symup{\Delta}$}}}%
  \def\nabla{\mathord{\raisebox{\depth}{\scalebox{1}[-1]{$\symup{\Delta}$}}}}%
  \def\checkmark{\popcheck{popdark}}%
  \def\star{\mathbin{\ast}}\def\bigstar{\mathord{\ast}}%
  \def\diamond{\mathbin{\scalebox{0.6}{$\Diamond$}}}%
  \def\bigcup{\mathop{\mathchoice{\raisebox{-0.3em}{\scalebox{1.7}{$\cup$}}}{\scalebox{1.25}{$\cup$}}{\cup}{\cup}}\displaylimits}%
  \def\bigcap{\mathop{\mathchoice{\raisebox{-0.3em}{\scalebox{1.7}{$\cap$}}}{\scalebox{1.25}{$\cap$}}{\cap}{\cap}}\displaylimits}%
  \def\lesseqqgtr{\mathrel{\lessgtr}}\def\bigoplus{\mathop{\oplus}\displaylimits}%
}
\providecommand{\ssi}{si et seulement si} % abréviation fréquente
\AtBeginDocument{\providecommand{\wideparen}{\overparen}} % arc de cercle

%% FIGFIX-BEGIN v5 — correctifs globaux des figures (docs/onbuch-generation/tools/figfix)
\usepackage{adjustbox}\usepackage{etoolbox}\tcbuselibrary{raster}
% toute figure TikZ/pgfplots plus large que la ligne est réduite à la largeur disponible
\BeforeBeginEnvironment{tikzpicture}{\begin{adjustbox}{max width=\linewidth}}
\AfterEndEnvironment{tikzpicture}{\end{adjustbox}}
% légendes pgfplots lisibles par-dessus les courbes
\pgfplotsset{every axis/.append style={legend style={fill=white,fill opacity=0.92,text opacity=1,draw=black!15,inner sep=3pt}}}
% étiquettes d'arêtes (syntaxe edge["..."]) avec fond blanc
\tikzset{every edge quotes/.append style={fill=white,inner sep=1.2pt}}
%% FIGFIX-END
\begin{document}

\pagegarde{El consumo responsable y la ética profesional (continuación) : expresión escrita}{Espagnol}{1ère A}{Módulo 4 : Vie économique}{EA16}{2 h}{Cette leçon mixte poursuit l'étude de la consommation responsable et de l'éthique professionnelle à travers un dialogue client-vendeur, l'expression de la cause, du but et de la condition (porque, para, si), la rédaction d'un court texte sur ses habitudes de consommation et la traduction de phrases sur le commerce et l'éthique.
\vspace{4pt}
\begin{multicols}{2}\small
\begin{itemize}
\item À la fin de cette leçon, je sais comprendre et jouer un dialogue entre un client et un vendeur en espagnol.
\item À la fin de cette leçon, je sais utiliser le lexique du commerce, de la consommation et de l'éthique professionnelle.
\item À la fin de cette leçon, je sais exprimer la cause avec porque, pues, como, a causa de, gracias a, debido a.
\item À la fin de cette leçon, je sais exprimer le but avec para + infinitif et para que + subjonctif.
\item À la fin de cette leçon, je sais exprimer la condition avec si + indicatif et je connais l'interdiction si + présent + futur.
\item À la fin de cette leçon, je sais rédiger un texte de 80 mots sur mes habitudes de consommation.
\end{itemize}\end{multicols}}

\begin{sommaire}
\begin{multicols}{2}
\begin{enumerate}
\item Texte support : dialogue au magasin
\item Lexique : le commerce et l'éthique
\item Grammaire 1 : la cause (porque et ses équivalents)
\item Grammaire 2 : le but (para) et la condition (si)
\item Traduction : thème et version sur le commerce
\item Production écrite guidée : mes habitudes de consommation
\item Activité d'intégration
\item Méthodes et exercices résolus
\item Exercices
\item Corrigés des exercices
\item Fiche bilan
\end{enumerate}
\end{multicols}
\end{sommaire}

\begin{objectifs}
\begin{itemize}
\item À la fin de cette leçon, je sais comprendre et jouer un dialogue entre un client et un vendeur en espagnol.
\item À la fin de cette leçon, je sais utiliser le lexique du commerce, de la consommation et de l'éthique professionnelle.
\item À la fin de cette leçon, je sais exprimer la cause avec porque, pues, como, a causa de, gracias a, debido a.
\item À la fin de cette leçon, je sais exprimer le but avec para + infinitif et para que + subjonctif.
\item À la fin de cette leçon, je sais exprimer la condition avec si + indicatif et je connais l'interdiction si + présent + futur.
\item À la fin de cette leçon, je sais rédiger un texte de 80 mots sur mes habitudes de consommation.
\item À la fin de cette leçon, je sais traduire des phrases sur le commerce et l'éthique du français vers l'espagnol et inversement.
\item À la fin de cette leçon, je sais éviter les pièges de traduction liés à porque, para et si.
\end{itemize}
\end{objectifs}

\begin{prerequis}
\begin{itemize}
\item Le présent de l'indicatif des verbes réguliers et irréguliers fréquents (Seconde)
\item Le lexique de base des achats et des prix (Seconde)
\item Les connecteurs simples : y, pero, también, además (Seconde)
\item La notion de subjonctif présent après querer que et es importante que (début d'année de Première)
\item Le vocabulaire de la consommation responsable vu à la leçon EA15 (première partie)
\end{itemize}
\end{prerequis}

\newpage

\coursec{Texte support : dialogue au magasin}

\courssub{Situation d'accroche et problématique}

C'est samedi au marché Mokolo de Yaoundé : la foule se presse entre les étals, les vendeuses appellent les clients, l'odeur des épices flotte dans l'air. Amina veut acheter un sac, mais elle hésite entre le sac importé bon marché et le sac en cuir fabriqué par un artisan camerounais. \esp{Porque quiero apoyar a los artesanos locales}, dit-elle à sa sœur. \esp{Si compras este, durará más}, répond celle-ci. Finalement, Amina choisit le sac de l'artisan : \esp{Compro este bolso para que mi dinero sirva a mi comunidad}.

Dans cette petite scène de la vie quotidienne, Amina a utilisé, sans même y penser, trois outils grammaticaux essentiels : la \cle{cause} (\esp{porque}), la \cle{condition} (\esp{si}) et le \cle{but} (\esp{para}). Voici notre problématique : \textbf{comment dialoguer en espagnol avec un vendeur, exprimer la cause de son choix, le but de son achat et la condition de sa décision ?} Pour y répondre, nous allons d'abord lire un dialogue authentique entre une cliente et un vendeur à Douala, observer le vocabulaire du commerce et de l'éthique, puis repérer les connecteurs qui structurent la conversation.

\courssub{Lecture du dialogue : « En la tienda de artesanía »}

Lisons le dialogue suivant. Il se déroule dans une boutique d'artisanat du quartier Akwa, à Douala. María, une jeune cliente, cherche un cadeau pour sa mère. Don Pedro, le vendeur, est un commerçant honnête qui défend le commerce équitable.

\begin{textebox}[En la tienda de artesanía — Diálogo en Douala]
\textbf{Vendedor:} Buenos días, señorita, bienvenida. ¿En qué puedo ayudarle?\par
\textbf{Cliente:} Buenos días. Busco un regalo para mi madre porque mañana es su cumpleaños.\par
\textbf{Vendedor:} ¡Qué bien! Tenemos bolsos hechos a mano por artesanos de la región. Mire este de cuero: es muy resistente.\par
\textbf{Cliente:} Es bonito. ¿Cuánto cuesta este?\par
\textbf{Vendedor:} Veinte mil francos. Es un poco caro, pero es de comercio justo: el artesano recibe un precio digno por su trabajo.\par
\textbf{Cliente:} Entiendo, pero es mucho dinero para mí. ¿Tiene algo más barato?\par
\textbf{Vendedor:} Sí, hay bolsos de segunda mano a diez mil francos, pero no tienen garantía. Este, en cambio, tiene una garantía de dos años.\par
\textbf{Cliente:} Hmm... Si me hace un descuento, lo compro.\par
\textbf{Vendedor:} De acuerdo. Como es para su madre, se lo dejo en dieciocho mil francos.\par
\textbf{Cliente:} ¡Perfecto! Me lo llevo. Compro aquí para apoyar a los artesanos locales.\par
\textbf{Vendedor:} Muchas gracias, señorita. Aquí tiene su recibo. Si hay algún problema, puede volver cuando quiera.\par
\textbf{Cliente:} Gracias, don Pedro. ¡Hasta pronto!\par
\textbf{Vendedor:} ¡Hasta luego! Y felicite a su madre de mi parte.
\end{textebox}

\textbf{Traduction du dialogue (pour vérifier votre compréhension) :}

« Vendeur : Bonjour mademoiselle, bienvenue. En quoi puis-je vous aider ? — Cliente : Bonjour. Je cherche un cadeau pour ma mère parce que demain c'est son anniversaire. — Vendeur : C'est bien ! Nous avons des sacs faits main par des artisans de la région. Regardez celui-ci en cuir : il est très résistant. — Cliente : Il est joli. Combien coûte celui-ci ? — Vendeur : Vingt mille francs. C'est un peu cher, mais c'est du commerce équitable : l'artisan reçoit un prix digne pour son travail. — Cliente : Je comprends, mais c'est beaucoup d'argent pour moi. Avez-vous quelque chose de moins cher ? — Vendeur : Oui, il y a des sacs d'occasion à dix mille francs, mais ils n'ont pas de garantie. Celui-ci, en revanche, a une garantie de deux ans. — Cliente : Hmm... Si vous me faites une réduction, je l'achète. — Vendeur : D'accord. Comme c'est pour votre mère, je vous le laisse à dix-huit mille francs. — Cliente : Parfait ! Je le prends. J'achète ici pour soutenir les artisans locaux. — Vendeur : Merci beaucoup, mademoiselle. Voici votre reçu. S'il y a un problème, vous pouvez revenir quand vous voulez. — Cliente : Merci, don Pedro. À bientôt ! — Vendeur : Au revoir ! Et félicitez votre mère de ma part. »

\begin{definition}[Glossaire : les mots nouveaux du dialogue]
\begin{itemize}
\item \esp{regatear} : marchander (verbe régulier en -ar) ; \esp{¿Puedo regatear el precio?} « Puis-je marchander le prix ? »
\item \esp{la garantía} : la garantie ; \esp{tiene una garantía de dos años} « il a une garantie de deux ans »
\item \esp{de segunda mano} : d'occasion ; \esp{un bolso de segunda mano} « un sac d'occasion »
\item \esp{el recibo} : le reçu, la quittance ; \esp{aquí tiene su recibo} « voici votre reçu »
\item \esp{hecho a mano} : fait main ; \esp{bolsos hechos a mano} « des sacs faits main »
\item \esp{el descuento} : la réduction, la remise ; \esp{hacer un descuento} « faire une réduction »
\item \esp{resistente} : résistant, solide ; \esp{el cuero} : le cuir
\end{itemize}
\end{definition}

\begin{definition}[El comercio justo]
\esp{El comercio justo} (le commerce équitable) est un système d'échange commercial qui garantit un \cle{prix équitable} au producteur ou à l'artisan, ainsi que de bonnes conditions de travail. Quand María achète un sac \esp{de comercio justo}, elle sait que l'artisan camerounais reçoit une part juste du prix final. C'est le contraire du système où les intermédiaires gardent presque tout l'argent.
\end{definition}

\begin{savaistu}[Le marchandage, une culture partagée]
En Espagne, on marchande (\esp{regatear}) surtout dans les marchés et les petites boutiques, jamais dans les grands magasins où les prix sont fixes (\esp{precios fijos}). En Amérique latine, au Cameroun et en Guinée équatoriale, le marchandage fait partie de la vie quotidienne au marché : c'est presque un art social ! En espagnol, on dit \esp{¿Me lo deja más barato?} (« Vous me le laissez moins cher ? ») pour commencer à négocier poliment.
\end{savaistu}

\courssub{Compréhension globale du dialogue}

Avant d'entrer dans les détails, vérifions la compréhension générale du texte.

\begin{exoresolu}[Compréhension globale : qui, où, quoi, pourquoi ?]
Répondez en français aux questions suivantes : 1) Qui sont les personnages du dialogue ? 2) Où se déroule la scène ? 3) Que cherche la cliente et pourquoi ? 4) Quel produit achète-t-elle finalement et à quel prix ?
\tcblower
\textbf{Solution détaillée :}
\begin{enumerate}
\item Les personnages sont \textbf{María}, une cliente, et \textbf{Don Pedro}, le vendeur de la boutique d'artisanat.
\item La scène se déroule dans une \textbf{boutique d'artisanat à Douala} (quartier Akwa).
\item María cherche \textbf{un cadeau pour sa mère}, parce que \textbf{demain c'est l'anniversaire de celle-ci} : \esp{Busco un regalo para mi madre porque mañana es su cumpleaños}.
\item Elle achète finalement \textbf{le sac en cuir fait main}, après négociation, à \textbf{dix-huit mille francs} au lieu de vingt mille : \esp{se lo dejo en dieciocho mil francos}.
\end{enumerate}
\end{exoresolu}

\begin{exercice}[Compréhension détaillée]
Répondez en espagnol, avec des phrases complètes :
\begin{enumerate}
\item ¿Qué regalo busca María?
\item ¿Por qué es caro el bolso de cuero?
\item ¿Qué diferencia hay entre los bolsos de segunda mano y el bolso de cuero?
\item ¿Qué descuento hace el vendedor?
\item ¿Para qué compra María en esta tienda?
\end{enumerate}
\end{exercice}

\begin{corrige}[Exercice : compréhension détaillée]
\begin{enumerate}
\item \esp{María busca un regalo para su madre, porque mañana es su cumpleaños.}
\item \esp{Es caro porque es de comercio justo y está hecho a mano por artesanos de la región.}
\item \esp{Los bolsos de segunda mano no tienen garantía; el bolso de cuero tiene una garantía de dos años.}
\item \esp{El vendedor le hace un descuento de dos mil francos: se lo deja en dieciocho mil.}
\item \esp{María compra allí para apoyar a los artesanos locales.}
\end{enumerate}
Remarquez que les questions 2 et 5 font appel à \esp{porque} (cause) et \esp{para} (but) : ces deux connecteurs seront étudiés en détail dans les sections de grammaire.
\end{corrige}

\courssub{Repérage des formules de politesse commerciale}

Un dialogue commercial en espagnol suit des étapes précises, avec des \cle{formules-clés} qu'il faut connaître par cœur. Observons-les dans le texte.

\begin{center}
\begin{tabularx}{\linewidth}{|X|X|X|}\hline
\rowcolor{popblueL} \textbf{Étape du dialogue} & \textbf{Formule espagnole} & \textbf{Traduction française} \\ \hline
Accueil du client & \esp{¿En qué puedo ayudarle?} & En quoi puis-je vous aider ? \\ \hline
Demande du client & \esp{Busco un regalo...} & Je cherche un cadeau... \\ \hline
Question sur le prix & \esp{¿Cuánto cuesta este?} & Combien coûte celui-ci ? \\ \hline
Négociation & \esp{Si me hace un descuento, lo compro.} & Si vous me faites une réduction, je l'achète. \\ \hline
Décision d'achat & \esp{Me lo llevo.} & Je le prends. \\ \hline
Conclusion & \esp{Aquí tiene su recibo.} & Voici votre reçu. \\ \hline
\end{tabularx}
\end{center}

\begin{attention}[Me lo llevo ou Me voy ?]
Ne confondez pas \esp{Me lo llevo} et \esp{Me voy} !
\begin{itemize}
\item \esp{Me lo llevo} = « \textbf{je le prends} » : c'est la formule de \textbf{décision d'achat} (du verbe \esp{llevarse} : emporter). \esp{lo} désigne l'objet acheté (le sac).
\item \esp{Me voy} = « \textbf{je m'en vais} » : cela signifie qu'on quitte le magasin (du verbe \esp{irse} : partir).
\end{itemize}
Erreur fréquente des francophones : dire \esp{me voy} pour annoncer qu'on achète le produit. Le vendeur comprendrait alors que vous partez sans rien acheter !
\end{attention}

\begin{exemplebox}[Exemples traduits à retenir]
\begin{itemize}
\item \esp{Si me hace un descuento, lo compro.} = « Si vous me faites une réduction, je l'achète. » (condition avec \esp{si} + présent, futur de décision au présent)
\item \esp{Compro aquí para apoyar a los artesanos.} = « J'achète ici pour soutenir les artisans. » (but avec \esp{para} + infinitif)
\item \esp{Busco un regalo porque mañana es su cumpleaños.} = « Je cherche un cadeau parce que demain c'est son anniversaire. » (cause avec \esp{porque})
\item \esp{¿Cuánto cuesta este bolso?} = « Combien coûte ce sac ? »
\item \esp{Es un poco caro, pero es de comercio justo.} = « C'est un peu cher, mais c'est du commerce équitable. »
\end{itemize}
\end{exemplebox}

\courssub{Les arguments du vendeur éthique}

Don Pedro n'est pas un vendeur ordinaire : il défend une consommation responsable. Relevons dans le texte ses trois grands arguments de vente.

\begin{center}
\begin{tabularx}{\linewidth}{|X|X|X|}\hline
\rowcolor{popgreenL} \textbf{Argument} & \textbf{Phrase du texte} & \textbf{Traduction} \\ \hline
La qualité & \esp{es muy resistente} / \esp{hecho a mano} & il est très résistant / fait main \\ \hline
Le commerce équitable & \esp{es de comercio justo: el artesano recibe un precio digno} & c'est du commerce équitable : l'artisan reçoit un prix digne \\ \hline
La garantie & \esp{tiene una garantía de dos años} & il a une garantie de deux ans \\ \hline
\end{tabularx}
\end{center}

Ces arguments sont typiques du \cle{consommateur responsable} : on ne regarde pas seulement le prix, mais aussi la qualité du produit, l'origine (\esp{artesanos de la región}) et l'impact social de l'achat. María, de son côté, montre qu'elle est une cliente responsable : \esp{Compro aquí para apoyar a los artesanos locales}.

\courssub{La structure d'un dialogue commercial}

Le dialogue suit toujours la même progression logique, en six étapes. Mémoriser cette structure permet de rejouer le dialogue sans apprendre chaque mot par cœur.

\begin{popfigure}
\begin{center}
\begin{tikzpicture}[
  etape/.style={rectangle, rounded corners=4pt, draw=popblue, fill=popblueL, text width=2.6cm, align=center, minimum height=1cm, font=\small},
  fleche/.style={-{Stealth[length=2.5mm]}, thick, poporange}
]
\node[etape, align=center] (e1) at (0,0){\textbf{1. Salutation}\\ \esp{Buenos días}\\ \esp{¿En qué puedo ayudarle?}};
\node[etape, right=0.5cm of e1, align=center] (e2){\textbf{2. Demande}\\ \esp{Busco un regalo...}};
\node[etape, right=0.5cm of e2, align=center] (e3){\textbf{3. Argumentation}\\ \esp{hecho a mano, comercio justo}};
\node[etape, below=0.6cm of e3, align=center] (e4){\textbf{4. Négociation}\\ \esp{Si me hace un descuento...}};
\node[etape, left=0.5cm of e4, align=center] (e5){\textbf{5. Décision}\\ \esp{Me lo llevo}};
\node[etape, left=0.5cm of e5, align=center] (e6){\textbf{6. Conclusion}\\ \esp{Aquí tiene su recibo}};
\draw[fleche] (e1) -- (e2);
\draw[fleche] (e2) -- (e3);
\draw[fleche] (e3) -- (e4);
\draw[fleche] (e4) -- (e5);
\draw[fleche] (e5) -- (e6);
\end{tikzpicture}
\end{center}
\legende{Organigramme des six étapes du dialogue commercial : de la salutation à la conclusion.}
\end{popfigure}

\courssub{Jouer le dialogue en binômes}

\begin{methode}[Comment jouer un dialogue en trois temps]
\begin{enumerate}
\item \textbf{Lire d'abord silencieusement} le dialogue deux fois : une fois pour le sens général, une fois en soulignant les formules-clés (\esp{¿En qué puedo ayudarle?}, \esp{¿Cuánto cuesta?}, \esp{Me lo llevo}).
\item \textbf{Mémoriser la structure, pas les mots} : retenez les six étapes de l'organigramme (salutation, demande, argumentation, négociation, décision, conclusion). Vous pourrez ainsi improviser avec vos propres produits et prix.
\item \textbf{Jouer avec un partenaire} : l'un est le vendeur, l'autre le client. Répétez deux fois, puis échangez les rôles. Parlez à voix haute, en articulant, et regardez votre partenaire.
\end{enumerate}
\end{methode}

\begin{experience}[Jeu de rôle : au marché de Douala]
Par binômes, rejouez le dialogue en \textbf{substituant le produit et le prix}. Choisissez un produit dans la liste : \esp{una camiseta de fútbol} (un maillot de football), \esp{un collar de perlas} (un collier de perles), \esp{una escultura de madera} (une sculpture en bois), \esp{un kilo de cacao} (un kilo de cacao). Choisissez un prix de départ (par exemple \esp{quince mil francos}) et négociez une réduction. Contraintes obligatoires :
\begin{itemize}
\item utiliser au moins une fois \esp{porque} (la cause), \esp{para} (le but) et \esp{si} (la condition) ;
\item utiliser les formules \esp{¿En qué puedo ayudarle?}, \esp{¿Cuánto cuesta?} et \esp{Me lo llevo} ;
\item le vendeur doit donner au moins deux arguments éthiques (qualité, commerce équitable, garantie, produit local).
\end{itemize}
Exemple de substitution : \esp{Busco una camiseta de fútbol porque mañana juega mi equipo. Si me la deja en doce mil, la compro. La quiero para regalársela a mi hermano.}
\end{experience}

\begin{aretenir}[L'essentiel de cette section]
\begin{itemize}
\item Un dialogue commercial suit six étapes : salutation, demande, argumentation, négociation, décision, conclusion.
\item Formules-clés : \esp{¿En qué puedo ayudarle?} (accueil), \esp{¿Cuánto cuesta?} (prix), \esp{Me lo llevo} (décision d'achat), \esp{aquí tiene su recibo} (conclusion).
\item Trois connecteurs structurent le dialogue : \esp{porque} (cause : \esp{porque mañana es su cumpleaños}), \esp{para} + infinitif (but : \esp{para apoyar a los artesanos}), \esp{si} + présent (condition : \esp{si me hace un descuento, lo compro}).
\item Vocabulaire nouveau : \esp{regatear}, \esp{la garantía}, \esp{de segunda mano}, \esp{el recibo}, \esp{hecho a mano}, \esp{el descuento}, \esp{el comercio justo}.
\item Piège : \esp{Me lo llevo} (je le prends) n'est pas \esp{Me voy} (je m'en vais).
\end{itemize}
\end{aretenir}

\coursec{Lexique : le commerce et l'éthique}

Dans la section précédente, nous avons lu et commenté un dialogue entre un client et une vendeuse dans un magasin de Yaoundé. Ce dialogue regorgeait de mots liés à l'achat, à la vente et au comportement honnête du commerçant. Pour pouvoir jouer nous-mêmes un tel dialogue, rédiger un texte sur nos habitudes de consommation et traduire des phrases sur le commerce, il faut d'abord maîtriser ce vocabulaire de manière précise : c'est l'objectif de cette section. Nous allons organiser les mots en quatre champs lexicaux, les observer en contexte, les mémoriser dans des tableaux espagnol-français, puis les réutiliser dans un exercice.

\courssub{Observer le vocabulaire en contexte}

Avant d'apprendre des listes de mots, lisons ce petit texte original qui les utilise naturellement. Soulignez mentalement tous les mots liés au commerce.

\begin{textebox}[Un sábado en el mercado de Mokolo]
El sábado por la mañana, Amina va al mercado de Mokolo, en Yaoundé, para comprar fruta y tela. En el primer puesto, pregunta al vendedor : « ¿Cuánto vale este tejido? » El vendedor responde : « Cuesta ocho mil francos. » Amina piensa que es demasiado caro y empieza a regatear. Después de cinco minutos, el vendedor le hace un descuento y la tela cuesta seis mil quinientos francos. Amina paga en efectivo y pide el recibo.

En el segundo puesto, una vendedora honrada le aconseja : « Esta tela es de mejor calidad y está hecha a mano. Además, es un producto local. » Amina está contenta porque le gusta comprar productos locales para ayudar a los artesanos de su país. Ella es una consumidora responsable : prefiere los productos ecológicos y reciclables, y a veces compra ropa de segunda mano para gastar menos y proteger el medio ambiente.

Antes de irse, Amina ve a un cliente furioso : un dependiente le ha engañado con una balanza falsa. El cliente grita : « ¡Esto es una estafa! Quiero devolver el producto y que me devuelvan mi dinero. » El jefe del puesto, que es un hombre honrado, respeta al cliente, le devuelve el dinero y le pide perdón. Él siempre cumple su palabra : por eso sus clientes confían en él y vuelven cada semana.
\end{textebox}

\textbf{Glossaire :} \esp{el puesto} = l'étal, le stand ; \esp{el tejido} = le tissu ; \esp{regatear} = marchander ; \esp{en efectivo} = en espèces ; \esp{la balanza} = la balance ; \esp{confiar en} = faire confiance à ; \esp{pedir perdón} = demander pardon.

\textbf{Questions de compréhension :}
\begin{enumerate}
\item ¿Cuánto cuesta la tela después del descuento?
\item ¿Por qué compra Amina productos locales?
\item ¿Qué hace el jefe del puesto cuando el cliente protesta?
\end{enumerate}

\textbf{Réponses attendues :} 1. \esp{Después del descuento, la tela cuesta seis mil quinientos francos.} 2. \esp{Compra productos locales para ayudar a los artesanos de su país.} 3. \esp{Respeta al cliente, le devuelve el dinero y le pide perdón.}

\courssub{Le champ lexical de l'achat}

Observons d'abord deux phrases du texte : \esp{Amina paga en efectivo y pide el recibo.} « Amina paie en espèces et demande le reçu. » ; \esp{El vendedor le hace un descuento.} « Le vendeur lui fait une remise. » Tous ces mots tournent autour de l'acte d'acheter.

\begin{definition}[Le vocabulaire de l'achat]
\begin{center}
\begin{tabularx}{\linewidth}{|l|X|}\hline
\rowcolor{popblueL} \textbf{Español} & \textbf{Français} \\ \hline
\esp{comprar} & acheter \\ \hline
\esp{vender} & vendre \\ \hline
\esp{pagar} & payer \\ \hline
\esp{costar (ue)} & coûter \\ \hline
\esp{valer} & valoir, coûter (\esp{¿Cuánto vale?}) \\ \hline
\esp{regatear} & marchander \\ \hline
\esp{el precio} & le prix \\ \hline
\esp{el descuento} & la remise, le rabais \\ \hline
\esp{la rebaja} & la réduction ; \esp{las rebajas} = les soldes \\ \hline
\esp{el recibo} & le reçu \\ \hline
\esp{la factura} & la facture \\ \hline
\esp{el dinero / en efectivo} & l'argent / en espèces \\ \hline
\esp{barato / caro} & bon marché / cher \\ \hline
\end{tabularx}
\end{center}
\end{definition}

\begin{propriete}[Le verbe \esp{costar} : un verbe à changement de voyelle]
\esp{Costar} est un verbe en \textbf{-ar} avec changement \textbf{o $\rightarrow$ ue} aux trois personnes du singulier et à la troisième du pluriel (comme \esp{poder}, \esp{volver}) :
\begin{center}
\begin{tabular}{|l|l|l|}\hline
\rowcolor{popblueL} yo & \esp{cuesto} & je coûte \\ \hline
tú & \esp{cuestas} & tu coûtes \\ \hline
él / ella / usted & \esp{cuesta} & il / elle coûte \\ \hline
nosotros & \esp{costamos} & nous coûtons \\ \hline
vosotros & \esp{costáis} & vous coûtez \\ \hline
ellos / ustedes & \esp{cuestan} & ils coûtent \\ \hline
\end{tabular}
\end{center}
\esp{¿Cuánto cuesta esta camisa? Cuesta cinco mil francos.} « Combien coûte cette chemise ? Elle coûte cinq mille francs. »

Remarque : le verbe \esp{valer} (« valoir ») est irrégulier seulement à la première personne du singulier : \esp{yo valgo}, mais \esp{¿cuánto vale?}, \esp{valen}. À l'oral, \esp{¿Cuánto vale?} et \esp{¿Cuánto cuesta?} sont interchangeables.
\end{propriete}

\begin{attention}[\esp{La factura} ou \esp{la cuenta} ?]
\esp{La factura} signifie « la facture » (document d'achat, au magasin ou pour un service). Mais au restaurant, « l'addition » se dit \esp{la cuenta} : \esp{Camarero, la cuenta, por favor.} « Garçon, l'addition, s'il vous plaît. » Ne dites jamais \esp{la factura} au restaurant, et ne traduisez pas \esp{la cuenta} par « le compte » dans ce contexte.
\end{attention}

\courssub{Le champ lexical du vendeur et du client}

\begin{definition}[Les personnes et le service]
\begin{center}
\begin{tabularx}{\linewidth}{|l|X|}\hline
\rowcolor{popgreenL} \textbf{Español} & \textbf{Français} \\ \hline
\esp{el cliente / la clienta} & le client / la cliente \\ \hline
\esp{el vendedor / la vendedora} & le vendeur / la vendeuse \\ \hline
\esp{el dependiente / la dependienta} & l'employé(e) de magasin, le/la vendeur(se) \\ \hline
\esp{atender (ie)} & servir, s'occuper de (un client) \\ \hline
\esp{aconsejar} & conseiller \\ \hline
\esp{la garantía} & la garantie \\ \hline
\esp{devolver (ue)} & rendre, rapporter ; rembourser \\ \hline
\esp{el puesto} & l'étal (au marché) \\ \hline
\esp{la tienda} & le magasin, la boutique \\ \hline
\esp{el mercado / el supermercado} & le marché / le supermarché \\ \hline
\end{tabularx}
\end{center}
\end{definition}

\begin{exemplebox}[Exemples traduits]
\esp{La dependienta atiende a los clientes con amabilidad.} « L'employée sert les clients avec amabilité. » (Attention : \esp{atender} exige la préposition \esp{a} devant la personne : \esp{atender \textbf{a} los clientes}.)

\esp{Este teléfono tiene una garantía de dos años.} « Ce téléphone a une garantie de deux ans. »

\esp{Quiero devolver estos zapatos porque son demasiado pequeños.} « Je veux rendre ces chaussures parce qu'elles sont trop petites. »
\end{exemplebox}

\begin{attention}[Faux ami : \esp{devolver} ne veut pas dire « devenir »]
\esp{Devolver} (o $\rightarrow$ ue) signifie « rendre (un objet), rapporter, rembourser » : \esp{El vendedor me devolvió el dinero.} « Le vendeur m'a remboursé (littéralement : m'a rendu l'argent). » Le verbe français « devenir » se traduit par \esp{hacerse} ou \esp{volverse} : \esp{Se hizo médico.} « Il est devenu médecin. » Beaucoup de francophones traduisent « devenir » par \esp{devolver} : c'est une erreur grave de sens.
\end{attention}

\courssub{Le champ lexical de la consommation responsable}

Ce champ lexical est au cœur du module « Vie économique » : il permet de parler de ses choix de consommateur.

\begin{definition}[Consommer autrement]
\begin{center}
\begin{tabularx}{\linewidth}{|l|X|}\hline
\rowcolor{popgoldL} \textbf{Español} & \textbf{Français} \\ \hline
\esp{el comercio justo} & le commerce équitable \\ \hline
\esp{el producto local} & le produit local \\ \hline
\esp{reciclable} & recyclable \\ \hline
\esp{reciclar} & recycler \\ \hline
\esp{ecológico, -a} & écologique \\ \hline
\esp{de segunda mano} & d'occasion, de seconde main \\ \hline
\esp{hecho a mano} & fait main \\ \hline
\esp{el consumidor / la consumidora} & le consommateur / la consommatrice \\ \hline
\esp{el consumidor responsable} & le consommateur responsable \\ \hline
\esp{el medio ambiente} & l'environnement \\ \hline
\esp{proteger} & protéger \\ \hline
\esp{gastar / ahorrar} & dépenser / économiser \\ \hline
\end{tabularx}
\end{center}
\end{definition}

\begin{exemplebox}[Exemples traduits]
\esp{Este producto es reciclable.} « Ce produit est recyclable. »

\esp{Compro café de comercio justo para ayudar a los campesinos.} « J'achète du café issu du commerce équitable pour aider les paysans. »

\esp{Las bolsas de plástico no son ecológicas.} « Les sacs en plastique ne sont pas écologiques. »

\esp{Esta mesa está hecha a mano por un artesano de Foumban.} « Cette table est faite main par un artisan de Foumban. »
\end{exemplebox}

\begin{propriete}[Accord et place des adjectifs]
Les adjectifs s'accordent en genre et en nombre avec le nom : \esp{un producto ecológico} / \esp{una tela ecológica} / \esp{productos reciclables}. En espagnol, l'adjectif se place en général \textbf{après} le nom : \esp{ropa de segunda mano} « des vêtements d'occasion ». Les locutions comme \esp{de segunda mano} restent invariables ; dans \esp{hecho a mano}, seul le participe \esp{hecho} s'accorde : \esp{una mesa hecha a mano}, \esp{unos productos hechos a mano}.
\end{propriete}

\courssub{Le champ lexical de l'éthique professionnelle}

\begin{textebox}[Mini-lexique en contexte]
Este vendedor es honrado : nunca engaña a sus clientes y siempre cumple su palabra.
\end{textebox}
Traduction : « Ce vendeur est honnête : il ne trompe jamais ses clients et tient toujours sa parole. »

\begin{definition}[L'honnêteté et ses contraires]
\begin{center}
\begin{tabularx}{\linewidth}{|l|X|}\hline
\rowcolor{poppinkL} \textbf{Español} & \textbf{Français} \\ \hline
\esp{honesto, -a / honrado, -a} & honnête \\ \hline
\esp{engañar} & tromper \\ \hline
\esp{la estafa} & l'arnaque, l'escroquerie \\ \hline
\esp{estafar} & arnaquer, escroquer \\ \hline
\esp{respetar al cliente} & respecter le client \\ \hline
\esp{la calidad} & la qualité \\ \hline
\esp{cumplir (su) palabra} & tenir (sa) parole \\ \hline
\esp{mentir (ie)} & mentir \\ \hline
\esp{la mentira / la verdad} & le mensonge / la vérité \\ \hline
\esp{confiar en} & faire confiance à \\ \hline
\end{tabularx}
\end{center}
\end{definition}

\begin{exemplebox}[Exemples traduits]
\esp{Un comerciante honrado no engaña a sus clientes.} « Un commerçant honnête ne trompe pas ses clients. »

\esp{La balanza falsa es una estafa.} « La balance truquée est une arnaque. »

\esp{Siempre cumplo mi palabra.} « Je tiens toujours ma parole. »
\end{exemplebox}

\begin{attention}[Orthographe : \esp{honesto} et \esp{honrado}]
Le \textbf{h} initial de \esp{honesto} et \esp{honrado} est muet : on prononce « o-nès-to », « on-ra-do ». Ne l'oubliez pas à l'écrit : on écrit \esp{honrado} avec un \textbf{h}, comme « honnête » en français, mais sans jamais prononcer cette lettre. Même chose pour \esp{honradez} « l'honnêteté ».
\end{attention}

\courssub{Les expressions utiles du dialogue client-vendeur}

Pour jouer un dialogue au magasin, mémorisez ces formules toutes faites :

\begin{center}
\begin{tabularx}{\linewidth}{|X|X|}\hline
\rowcolor{popblueL} \textbf{Español} & \textbf{Français} \\ \hline
\esp{¿En qué puedo ayudarle?} & En quoi puis-je vous aider ? \\ \hline
\esp{¿Cuánto vale? / ¿Cuánto cuesta?} & Combien ça coûte ? \\ \hline
\esp{Es demasiado caro.} & C'est trop cher. \\ \hline
\esp{¿Tiene algo más barato?} & Avez-vous quelque chose de moins cher ? \\ \hline
\esp{¿Me hace un descuento?} & Vous me faites une remise ? \\ \hline
\esp{Me lo llevo.} & Je le prends. \\ \hline
\esp{¿Puedo pagar en efectivo?} & Puis-je payer en espèces ? \\ \hline
\esp{¿Tiene garantía?} & Est-ce qu'il y a une garantie ? \\ \hline
\end{tabularx}
\end{center}

\begin{exemplebox}[Exemple clé]
\esp{¿Tiene algo más barato?} « Avez-vous quelque chose de moins cher ? » Remarquez : \esp{algo} = « quelque chose », et le comparatif \esp{más barato} = « moins cher » (littéralement « plus bon marché »).
\end{exemplebox}

\begin{attention}[\esp{Me lo llevo} : « je le prends »]
Pour dire « je le prends » au moment de l'achat, l'espagnol utilise \esp{llevarse} : \esp{Me lo llevo.} Ne traduisez pas par \esp{Lo tomo} (qui signifie « je le bois / je le prends en main »). \esp{Llevarse algo} = « emporter, prendre quelque chose (en l'achetant) ».
\end{attention}

\courssub{Carte mentale : organiser le lexique}

La carte mentale ci-dessous relie les quatre familles de mots autour de l'idée centrale de la leçon. Recopiez-la dans votre cahier et ajoutez deux mots par branche.

\begin{popfigure}
\begin{center}
\begin{center}\tcbox[colback=popblue,colframe=popblue,coltext=white,arc=9pt,boxsep=4pt,left=6pt,right=6pt]{\bfseries\small el consumo responsable}\end{center}
\vspace{-2pt}
\begin{tcbraster}[raster columns=2,raster equal height=rows,raster column skip=6pt,raster row skip=6pt,enhanced,blankest]
\begin{tcolorbox}[enhanced,colback=poporange,colframe=poporange,coltext=white,boxrule=0.9pt,arc=3pt,left=4pt,right=4pt,top=3pt,bottom=3pt,boxsep=1pt,halign=left]\bfseries\scriptsize comprar\end{tcolorbox}
\begin{tcolorbox}[enhanced,colback=popgreen,colframe=popgreen,coltext=white,boxrule=0.9pt,arc=3pt,left=4pt,right=4pt,top=3pt,bottom=3pt,boxsep=1pt,halign=left]\bfseries\scriptsize vender\end{tcolorbox}
\begin{tcolorbox}[enhanced,colback=poppink,colframe=poppink,coltext=white,boxrule=0.9pt,arc=3pt,left=4pt,right=4pt,top=3pt,bottom=3pt,boxsep=1pt,halign=left]\bfseries\scriptsize ética\end{tcolorbox}
\begin{tcolorbox}[enhanced,colback=popgold,colframe=popgold,coltext=white,boxrule=0.9pt,arc=3pt,left=4pt,right=4pt,top=3pt,bottom=3pt,boxsep=1pt,halign=left]\bfseries\scriptsize medio ambiente\end{tcolorbox}
\end{tcbraster}

\end{center}
\legende{Carte mentale du lexique : quatre branches (acheter, vendre, éthique, environnement) autour de la consommation responsable.}
\end{popfigure}

\begin{savaistu}[Le marchandage : \esp{el regateo}]
Au Cameroun comme dans toute l'Amérique latine, le marchandage (\esp{el regateo}) est une pratique culturelle courante au marché : le prix annoncé est un point de départ, et discuter fait partie du plaisir de l'échange. En revanche, dans les grandes surfaces (\esp{los supermercados}) et les magasins modernes, les prix sont fixes et marchander serait mal vu. En Espagne, on attend plutôt les \esp{rebajas} (les soldes), en janvier et en juillet, pour obtenir de grandes réductions.
\end{savaistu}

\courssub{Exercice résolu et entraînement}

\begin{exoresolu}[Complétez avec le mot qui convient]
Complétez les phrases avec : \esp{recibo}, \esp{devolver}, \esp{regatear}, \esp{garantía}, \esp{estafa}.

1. En el mercado de Mokolo, me gusta \ldots\ para pagar menos.
2. Este vendedor es honrado : nunca hace una \ldots\ a sus clientes.
3. Cuando pago, siempre pido el \ldots.
4. Quiero \ldots\ esta camisa porque es demasiado grande.
5. Este ordenador tiene una \ldots\ de un año.
\tcblower
\textbf{Solution détaillée :}

1. \esp{regatear} : « au marché de Mokolo, j'aime marchander pour payer moins ». Après \esp{gustar}, on met l'infinitif.

2. \esp{estafa} : « ce vendeur est honnête : il ne fait jamais d'arnaque à ses clients ». Le mot doit être un nom féminin précédé de \esp{una}.

3. \esp{recibo} : « quand je paie, je demande toujours le reçu ». Ne pas confondre avec \esp{factura} : au marché, on reçoit un simple reçu.

4. \esp{devolver} : « je veux rendre cette chemise parce qu'elle est trop grande ». Après \esp{querer}, on met l'infinitif ; \esp{devolver} = rendre un objet, pas « devenir ».

5. \esp{garantía} : « cet ordinateur a une garantie d'un an ».
\end{exoresolu}

\begin{exercice}[À vous de jouer]
1. Traduisez en espagnol : a) « Avez-vous quelque chose de moins cher ? » b) « Ce produit est recyclable et fait main. » c) « Je le prends, voici l'argent. » d) « Un vendeur honnête tient toujours sa parole. »

2. Classez ces mots dans les quatre branches de la carte mentale : \esp{el precio}, \esp{honrado}, \esp{reciclable}, \esp{la clienta}, \esp{engañar}, \esp{el descuento}, \esp{el medio ambiente}, \esp{atender}.

3. Rédigez quatre phrases sur vos propres habitudes d'achat en utilisant : \esp{comprar}, \esp{regatear}, \esp{de segunda mano}, \esp{producto local}.
\end{exercice}

\begin{corrige}[Exercice : le commerce et l'éthique]
1. a) \esp{¿Tiene algo más barato?} b) \esp{Este producto es reciclable y está hecho a mano.} c) \esp{Me lo llevo, aquí está el dinero.} d) \esp{Un vendedor honrado siempre cumple su palabra.}

2. \textbf{comprar} : \esp{el precio}, \esp{el descuento} ; \textbf{vender} : \esp{la clienta}, \esp{atender} ; \textbf{ética} : \esp{honrado}, \esp{engañar} ; \textbf{medio ambiente} : \esp{reciclable}, \esp{el medio ambiente}.

3. Exemple de production : \esp{Compro fruta en el mercado. Me gusta regatear con los vendedores. A veces compro ropa de segunda mano. Prefiero los productos locales porque ayudan a mi país.}
\end{corrige}

\begin{aretenir}[L'essentiel du lexique]
\begin{itemize}
\item Achat : \esp{comprar}, \esp{pagar}, \esp{costar (ue)}, \esp{valer}, \esp{regatear}, \esp{el precio}, \esp{el descuento}, \esp{el recibo}.
\item Personnes : \esp{el cliente}, \esp{el vendedor}, \esp{el dependiente} ; \esp{atender a}, \esp{aconsejar}, \esp{devolver (ue)} = rendre/rembourser (jamais « devenir »).
\item Consommation responsable : \esp{el comercio justo}, \esp{producto local}, \esp{reciclable}, \esp{ecológico}, \esp{de segunda mano}, \esp{hecho a mano}.
\item Éthique : \esp{honrado}, \esp{engañar}, \esp{la estafa}, \esp{cumplir su palabra}, \esp{respetar al cliente}.
\item Formules : \esp{¿Cuánto vale?}, \esp{¿Tiene algo más barato?}, \esp{Me lo llevo.}
\end{itemize}
\end{aretenir}

Maintenant que ce vocabulaire est bien installé, nous pouvons passer à la grammaire : dans la section suivante, nous apprendrons à expliquer \emph{pourquoi} nous achetons tel ou tel produit grâce aux mots de la cause, \esp{porque} et ses équivalents.

\coursec{Grammaire 1 : la cause (porque et ses équivalents)}

Dans la section précédente, nous avons découvert le lexique du commerce et de l'éthique : \esp{el comercio justo}, \esp{la rebaja}, \esp{el cliente}, \esp{el vendedor}, \esp{la calidad}, \esp{el precio}. Ce vocabulaire prend tout son sens lorsqu'il est employé dans des phrases complètes. Or, dans le dialogue entre le client et la vendeuse, un mot revient sans cesse : \esp{porque}. Le client explique \emph{pourquoi} il cherche un cadeau, la vendeuse justifie les prix, chacun donne des raisons. Exprimer la cause, c'est savoir dire « parce que », « car », « comme », « à cause de », « grâce à ». C'est l'objet de cette leçon de grammaire.

\courssub{Observation : la cause dans le dialogue}

\begin{textebox}[En la tienda de comercio justo — extrait]
— Buenos días. Busco un regalo porque mañana es el cumpleaños de mi madre.\par
— Muy bien. Le recomiendo este bolso: es caro, pero es de comercio justo.\par
— ¿Por qué cuesta tanto?\par
— Porque los artesanos reciben un salario digno. Como es un producto artesanal, el precio es más alto.\par
— Entiendo. No compré nada en el supermercado ayer, pues todo era de mala calidad.\par
— Aquí, gracias a la calidad, nuestros productos duran muchos años.\par
— Perfecto. Compro el bolso por ser bonito y honesto.
\end{textebox}

\textbf{Traduction.}\par
« — Bonjour. Je cherche un cadeau parce que demain c'est l'anniversaire de ma mère. — Très bien. Je vous recommande ce sac : il est cher, mais c'est du commerce équitable. — Pourquoi coûte-t-il si cher ? — Parce que les artisans reçoivent un salaire digne. Comme c'est un produit artisanal, le prix est plus élevé. — Je comprends. Je n'ai rien acheté au supermarché hier, car tout était de mauvaise qualité. — Ici, grâce à la qualité, nos produits durent de nombreuses années. — Parfait. J'achète le sac parce qu'il est beau et honnête. »

\textbf{Glossaire.}\par
\esp{el bolso} : le sac ; \esp{artesanal} : artisanal ; \esp{un salario digno} : un salaire digne ; \esp{durar} : durer ; \esp{ayer} : hier.

\textbf{Questions d'observation.}\par
\begin{enumerate}
\item Souligne tous les mots qui introduisent une cause dans ce dialogue.
\item Quel mot espagnol correspond à « pourquoi » dans la question ? Et à « parce que » dans la réponse ?
\item Quel connecteur se trouve en tête de phrase ? Lequel se trouve en milieu de phrase ?
\end{enumerate}

On observe cinq procédés principaux pour exprimer la cause : \esp{porque}, \esp{pues}, \esp{como}, \esp{gracias a} (et ses équivalents \esp{a causa de}, \esp{debido a}) et \esp{por}. Étudions-les une à une.

\courssub{Porque + indicatif : la cause certaine}

\begin{definition}[Porque]
\esp{Porque} est la conjonction de cause la plus fréquente. Elle introduit une \cle{proposition subordonnée} qui donne la raison de l'action principale. Elle répond à la question \esp{¿Por qué?} (pourquoi ?). Le verbe qui suit est à l'\cle{indicatif}, car la cause est présentée comme un fait certain.
\end{definition}

\esp{Busco un regalo porque mañana es su cumpleaños.} « Je cherche un cadeau parce que demain c'est son anniversaire. »

\esp{No compro en esa tienda porque es muy cara.} « Je n'achète pas dans ce magasin parce qu'il est très cher. »

\begin{propriete}[Emploi de porque]
\begin{enumerate}
\item \esp{Porque} + verbe à l'indicatif exprime une cause réelle, certaine.
\item Dans une réponse courte à \esp{¿Por qué?}, la réponse commence par \esp{Porque} : \esp{—¿Por qué no vienes? —Porque estoy ocupado.} « — Pourquoi ne viens-tu pas ? — Parce que je suis occupé. »
\item En revanche, dans une phrase complète, on évite de placer \esp{porque} en tête de phrase : on lui préfère \esp{como} (voir plus bas).
\end{enumerate}
\end{propriete}

\courssub{Pues : la cause explicative faible}

\begin{definition}[Pues]
\esp{Pues} correspond au français « car ». Il introduit une explication, une justification atténuée, et se place toujours \cle{après} la proposition principale, jamais en tête de phrase. Il est un peu plus soutenu ou plus écrit que \esp{porque}.
\end{definition}

\esp{No compré nada, pues era muy caro.} « Je n'ai rien acheté, car c'était très cher. »

\esp{Cierra la puerta, pues hace frío.} « Ferme la porte, car il fait froid. »

\begin{attention}[Pues n'est pas « puis »]
Faux ami classique : \esp{pues} ne signifie presque jamais « puis » (qui se dit \esp{luego} ou \esp{después}). \esp{Pues} vaut « car », « eh bien », « donc » selon le contexte. De même, \esp{entonces} = « alors », pas « en temps ».
\end{attention}

\courssub{Como en tête de phrase : « comme »}

\begin{propriete}[Como causal]
Quand la proposition de cause est placée \cle{en tête de phrase}, l'espagnol emploie \esp{como} (= « comme », « puisque »), suivi de l'indicatif. C'est la construction la plus élégante pour commencer une phrase par une cause.
\end{propriete}

\esp{Como es de comercio justo, cuesta más.} « Comme c'est du commerce équitable, cela coûte plus cher. »

\esp{Como no tengo dinero, no compro nada.} « Comme je n'ai pas d'argent, je n'achète rien. »

Compare les deux constructions :

\begin{center}
\begin{tabularx}{\linewidth}{|X|X|}
\hline
\rowcolor{popblueL} \textbf{Cause après (porque)} & \textbf{Cause avant (como)} \\
\hline
\esp{No compro nada porque no tengo dinero.} & \esp{Como no tengo dinero, no compro nada.} \\
\hline
\esp{Cuesta más porque es artesanal.} & \esp{Como es artesanal, cuesta más.} \\
\hline
\end{tabularx}
\end{center}

\begin{attention}[Como placé ailleurs]
Attention : \esp{como} placé au milieu de la phrase signifie « comme » au sens de comparaison (\esp{Trabaja como un artesano} : « Il travaille comme un artisan ») ou « je mange » (verbe \esp{comer}, 1\up{re} personne du singulier). Le \esp{como} causal se reconnaît à sa place : toujours en tête de phrase, suivi d'une virgule.
\end{attention}

\courssub{Les causes nominales : a causa de, gracias a, debido a, por}

Quand la cause est un \cle{nom} (et non une phrase complète avec verbe conjugué), l'espagnol utilise des prépositions de cause. Le choix dépend de la valeur de la cause : positive, négative ou neutre.

\begin{definition}[Les trois prépositions de cause]
\begin{itemize}
\item \esp{a causa de} + nom : cause \cle{négative}, fâcheuse. « à cause de ».
\item \esp{gracias a} + nom : cause \cle{positive}, bénéfique. « grâce à ».
\item \esp{debido a} + nom : cause \cle{neutre}, factuelle. « en raison de », « dû à ».
\end{itemize}
\end{definition}

\begin{exemplebox}[Exemples traduits]
\esp{Gracias a la rebaja, compré dos camisas.} « Grâce à la réduction, j'ai acheté deux chemises. »\par
\esp{La tienda cerró a causa de la crisis.} « Le magasin a fermé à cause de la crise. »\par
\esp{El mercado está cerrado debido a la fiesta nacional.} « Le marché est fermé en raison de la fête nationale. »\par
\esp{Gracias a su honestidad, el vendedor tiene muchos clientes.} « Grâce à son honnêteté, le vendeur a beaucoup de clients. »
\end{exemplebox}

\begin{propriete}[Por + nom ou infinitif]
\esp{Por} exprime aussi la cause, devant un nom ou un infinitif : \esp{por falta de dinero} « faute d'argent », \esp{por ser honesto} « parce qu'il est honnête / pour être honnête (cause) ». Devant un infinitif, le sujet de l'infinitif est le même que celui du verbe principal.
\end{propriete}

\esp{Compro el bolso por ser bonito y honesto.} « J'achète le sac parce qu'il est beau et honnête. »

\esp{No viajamos por falta de dinero.} « Nous ne voyageons pas faute d'argent. »

\esp{Lo castigaron por llegar tarde.} « On l'a puni parce qu'il est arrivé en retard. »

\courssub{Les quatre formes : por qué, porque, porqué, por que}

C'est un point d'orthographe capital, très souvent évalué. L'espagnol distingue \emph{quatre} graphies différentes, que les francophones confondent facilement.

\begin{center}
\begin{tabularx}{\linewidth}{|l|X|X|}
\hline
\rowcolor{popblueL} \textbf{Forme} & \textbf{Sens} & \textbf{Exemple} \\
\hline
\esp{¿por qué?} (2 mots, accent) & question : « pourquoi ? » & \esp{¿Por qué compras eso?} « Pourquoi achètes-tu cela ? » \\
\hline
\esp{porque} (1 mot, sans accent) & réponse : « parce que » & \esp{Porque me gusta.} « Parce que cela me plaît. » \\
\hline
\esp{el porqué} (1 mot, accent, nom) & nom masculin : « la raison, le pourquoi » & \esp{No entiendo el porqué de su decisión.} « Je ne comprends pas la raison de sa décision. » \\
\hline
\esp{por que} (2 mots, sans accent) & rare : \esp{por} + pronom relatif \esp{que} & \esp{la razón por que luchamos} « la raison pour laquelle nous luttons » \\
\hline
\end{tabularx}
\end{center}

\begin{attention}[Erreurs fréquentes des francophones]
\begin{enumerate}
\item Ne jamais écrire \esp{por que} (deux mots, sans accent) pour dire « parce que » : c'est la faute la plus fréquente. « Parce que » = \esp{porque}, en un seul mot.
\item \esp{Porque} ne se traduit jamais par « pour » : « pour » exprime le but (\esp{para}), pas la cause. \esp{Lo hago porque quiero} = « Je le fais parce que je veux » ; \esp{Lo hago para aprender} = « Je le fais pour apprendre ».
\item Dans la question, n'oublie ni l'accent ni les deux points d'interrogation : \esp{¿Por qué?}
\end{enumerate}
\end{attention}

\begin{popfigure}
\begin{center}
\begin{tikzpicture}[
  box/.style={draw=popblue, fill=popblueL, rounded corners=3pt, align=center, minimum height=1.1cm, text width=3.1cm, font=\small},
  lab/.style={font=\small\itshape, text=popdark}
]
\node[box, align=center] (q) at (0,2.6){\textbf{¿Por qué?}\\ « pourquoi ? »};
\node[box, align=center] (p) at (4.6,2.6){\textbf{porque}\\ « parce que »};
\node[box, align=center] (ac) at (0,0){\textbf{a causa de}\\ « à cause de » (négatif)};
\node[box, align=center] (ga) at (4.6,0){\textbf{gracias a}\\ « grâce à » (positif)};
\node[box, align=center] (co) at (2.3,-2.6){\textbf{como} + phrase\\ « comme » (en tête)};
\draw[-{Stealth}, thick, poporange] (q) -- (p) node[midway, above, lab] {question $\rightarrow$ réponse};
\draw[-{Stealth}, thick, popgreen] (ac) -- (ga) node[midway, above, lab] {cause nominale};
\end{tikzpicture}
\end{center}
\legende{Les principaux connecteurs de cause : question, réponse et causes nominales.}
\end{popfigure}

\courssub{Comparaison avec le français}

\begin{center}
\begin{tabularx}{\linewidth}{|X|X|X|}
\hline
\rowcolor{popblueL} \textbf{Français} & \textbf{Español} & \textbf{Remarque} \\
\hline
parce que & \esp{porque} & un seul mot, sans accent \\
\hline
pourquoi ? & \esp{¿por qué?} & deux mots, accent sur \esp{qué} \\
\hline
car & \esp{pues} & jamais en tête de phrase \\
\hline
comme (cause, en tête) & \esp{como} & suivi d'une virgule \\
\hline
à cause de & \esp{a causa de} & cause négative \\
\hline
grâce à & \esp{gracias a} & cause positive \\
\hline
en raison de, dû à & \esp{debido a} & cause neutre \\
\hline
faute de, par (cause) & \esp{por} + nom/infinitif & \esp{por falta de} \\
\hline
\end{tabularx}
\end{center}

Remarque essentielle : le français utilise « parce que » partout, en tête ou au milieu de la phrase. L'espagnol, lui, réserve la tête de phrase à \esp{como} : « Parce qu'il pleut, je reste » se dira naturellement \esp{Como llueve, me quedo}, et non \esp{Porque llueve...} en début de phrase soutenue.

\begin{methode}[Choisir le bon connecteur de cause]
\begin{enumerate}
\item La cause est-elle un \textbf{nom} ? → cause positive : \esp{gracias a} ; cause négative : \esp{a causa de} ; cause neutre : \esp{debido a}.
\item La cause est-elle une \textbf{phrase complète} (verbe conjugué) ? → au milieu : \esp{porque} ; en tête : \esp{como} ; justification faible après la principale : \esp{pues}.
\item S'agit-il d'une \textbf{question} ? → \esp{¿por qué?} (deux mots, accent).
\item S'agit-il d'un \textbf{nom} « la raison » ? → \esp{el porqué} (un mot, accent).
\end{enumerate}
\end{methode}

\begin{exoresolu}[Complète avec la forme correcte]
Énoncé. Complète chaque phrase avec \esp{porque}, \esp{por qué}, \esp{como}, \esp{gracias a}, \esp{a causa de} ou \esp{pues}.\par
1. \esp{¿......... no compraste el regalo?}\par
2. \esp{No lo compré ......... era demasiado caro.}\par
3. \esp{......... no tengo tiempo, iré al mercado mañana.}\par
4. \esp{......... la crisis, muchos comercios cerraron.}\par
5. \esp{......... sus precios bajos, esta tienda tiene muchos clientes.}\par
6. \esp{Debe de estar enfermo, ......... no ha venido a trabajar.}
\tcblower
Solution détaillée.\par
1. \esp{¿Por qué no compraste el regalo?} — Question : deux mots, accent. « Pourquoi n'as-tu pas acheté le cadeau ? »\par
2. \esp{No lo compré porque era demasiado caro.} — Réponse, cause certaine, milieu de phrase. « Je ne l'ai pas acheté parce qu'il était trop cher. »\par
3. \esp{Como no tengo tiempo, iré al mercado mañana.} — Cause en tête de phrase : \esp{como}. « Comme je n'ai pas le temps, j'irai au marché demain. »\par
4. \esp{A causa de la crisis, muchos comercios cerraron.} — Cause nominale négative. « À cause de la crise, beaucoup de commerces ont fermé. »\par
5. \esp{Gracias a sus precios bajos, esta tienda tiene muchos clientes.} — Cause nominale positive. « Grâce à ses prix bas, ce magasin a beaucoup de clients. »\par
6. \esp{Debe de estar enfermo, pues no ha venido a trabajar.} — Justification faible après la principale : \esp{pues}. « Il doit être malade, car il n'est pas venu travailler. »
\end{exoresolu}

\begin{exercice}[À ton tour]
Traduis en espagnol :\par
1. « J'achète ce produit parce qu'il est de bonne qualité. »\par
2. « Comme il n'y a pas de réduction, je n'achète rien. »\par
3. « Pourquoi ce magasin est-il fermé ? — À cause de la fête. »\par
4. « Grâce au commerce équitable, les artisans vivent mieux. »
\end{exercice}

\begin{corrige}[Exercice « À ton tour »]
1. \esp{Compro este producto porque es de buena calidad.}\par
2. \esp{Como no hay rebaja, no compro nada.}\par
3. \esp{¿Por qué está cerrada esta tienda? — A causa de la fiesta.}\par
4. \esp{Gracias al comercio justo, los artesanos viven mejor.}
\end{corrige}

\begin{savaistu}[Un classique des examens... même en Espagne]
La distinction entre \esp{por qué}, \esp{porque}, \esp{porqué} et \esp{por que} est l'un des points d'orthographe les plus difficiles pour les hispanophones eux-mêmes. En Espagne, les manuels scolaires et les concours y consacrent des chapitres entiers, et la Real Academia Española reçoit chaque année des milliers de questions à ce sujet. Maîtriser ces quatre formes, c'est donc écrire mieux que beaucoup de natifs. Une astuce de mémorisation : la question « s'ouvre » (deux mots : \esp{por qué}), la réponse « se ferme » (un mot : \esp{porque}).
\end{savaistu}

\begin{aretenir}[L'essentiel]
\begin{itemize}
\item \esp{Porque} + indicatif : la cause certaine, réponse à \esp{¿por qué?}.
\item \esp{Pues} : « car », toujours après la principale.
\item \esp{Como} : « comme », uniquement en tête de phrase.
\item \esp{A causa de} (négatif), \esp{gracias a} (positif), \esp{debido a} (neutre) + nom.
\item \esp{Por} + nom ou infinitif : \esp{por falta de dinero}.
\item Quatre graphies : \esp{¿por qué?} / \esp{porque} / \esp{el porqué} / \esp{por que} (rare).
\end{itemize}
\end{aretenir}

\coursec{Grammaire 2 : le but (para) et la condition (si)}

Dans la section précédente, nous avons appris à exprimer la \cle{cause} avec \esp{porque} et ses équivalents : \esp{Compro aquí porque los productos son de calidad} (« J'achète ici parce que les produits sont de qualité »). Mais dans le dialogue entre le client et le vendeur, on trouve aussi d'autres petits mots essentiels : \esp{para} et \esp{si}. Ils permettent d'exprimer le \cle{but} (« pour ») et la \cle{condition} (« si »). Ce sont deux notions fondamentales pour rédiger votre texte sur vos habitudes de consommation et pour réussir vos traductions. Observons d'abord des exemples tirés de situations de commerce.

\courssub{Observation : à quoi servent \esp{para} et \esp{si} ?}

Lisez attentivement ces phrases. Soulignez mentalement ce qui suit \esp{para} et ce qui suit \esp{si}.

\begin{exemplebox}[Phrases d'observation]
\begin{enumerate}
\item \esp{Compro aquí para apoyar a los artesanos.} « J'achète ici pour soutenir les artisans. »
\item \esp{Te doy un descuento para que vuelvas.} « Je te fais une remise pour que tu reviennes. »
\item \esp{Este regalo es para mi madre.} « Ce cadeau est pour ma mère. »
\item \esp{Si me hace un descuento, lo compro.} « Si vous me faites une remise, je l'achète. »
\item \esp{Si tuviera dinero, compraría más.} « Si j'avais de l'argent, j'achèterais davantage. »
\end{enumerate}
\end{exemplebox}

Que remarquez-vous ? Dans la phrase 1, après \esp{para}, on trouve un \textbf{infinitif} (\esp{apoyar}) et le sujet des deux verbes est le même (« j'achète » et « je soutiens »). Dans la phrase 2, les sujets sont différents (« je » donne, « tu » reviens) : l'espagnol utilise alors \esp{para que} suivi d'un verbe conjugué spécial, le \textbf{subjonctif} (\esp{vuelvas}). Dans la phrase 3, \esp{para} est suivi d'un \textbf{nom} : il indique le destinataire. Dans les phrases 4 et 5, \esp{si} introduit une condition, mais les temps employés ne sont pas les mêmes. Détaillons chaque cas.

\courssub{Le but avec \esp{para}}

\begin{propriete}[\esp{para} + infinitif : le but quand le sujet est identique]
Pour exprimer le \textbf{but} (« pour », « afin de ») lorsque le sujet des deux actions est \textbf{le même}, l'espagnol utilise :

\begin{center}
\textbf{\esp{para} + INFINITIF}
\end{center}

\begin{itemize}
\item \esp{Ahorro para comprar un teléfono.} « J'économise pour acheter un téléphone. » (c'est moi qui économise et qui achèterai)
\item \esp{Trabaja para ganarse la vida.} « Il travaille pour gagner sa vie. »
\item \esp{Voy al mercado para comprar verduras frescas.} « Je vais au marché pour acheter des légumes frais. »
\end{itemize}

L'infinitif espagnol correspond exactement à l'infinitif français après « pour ». C'est le cas le plus simple.
\end{propriete}

\begin{propriete}[\esp{para que} + subjonctif : le but quand les sujets sont différents]
Lorsque le sujet de l'action principale et celui du but sont \textbf{différents}, l'espagnol utilise :

\begin{center}
\textbf{\esp{para que} + SUBJONCTIF}
\end{center}

\begin{itemize}
\item \esp{Ahorro para que mi hijo estudie.} « J'économise pour que mon fils étudie. » (j'économise, mais c'est lui qui étudie)
\item \esp{Te doy un descuento para que vuelvas.} « Je te fais une remise pour que tu reviennes. »
\item \esp{El vendedor baja los precios para que los clientes compren más.} « Le vendeur baisse les prix pour que les clients achètent davantage. »
\end{itemize}

\textbf{Point de comparaison :} le français utilise aussi un subjonctif après « pour que », donc la structure est parallèle. Mais attention : en espagnol, le subjonctif après \esp{para que} est \textbf{obligatoire}, sans exception, et il faut savoir le former (présent de subjonctif : \esp{compre, vuelvas, estudie, sea, vaya, tenga...}).
\end{propriete}

\begin{popfigure}
\begin{tikzpicture}[
node distance=0.9cm,
startstop/.style={rectangle, rounded corners, draw=popblue, fill=popblueL, text width=4.2cm, align=center, font=\small},
decision/.style={diamond, draw=poporange, fill=poporangeL, text width=3.4cm, align=center, font=\small, aspect=2.2},
result/.style={rectangle, rounded corners, draw=popgreen, fill=popgreenL, text width=4.6cm, align=center, font=\small},
arrow/.style={-{Stealth[length=3mm]}, thick, popdark}]
\node[startstop] (start) {Je veux exprimer le but (« pour ») devant un verbe};
\node[decision, below=of start] (q) {Le sujet des deux verbes est-il le même ?};
\node[result, below left=1.1cm and -1.2cm of q, align=center] (oui){\textbf{OUI} : \esp{para} + infinitif\\ \esp{Ahorro para comprar.}};
\node[result, below right=1.1cm and -1.2cm of q, align=center] (non){\textbf{NON} : \esp{para que} + subjonctif\\ \esp{Ahorro para que estudie.}};
\draw[arrow] (start) -- (q);
\draw[arrow] (q) -| (oui) node[midway, above left, font=\small\itshape] {même sujet};
\draw[arrow] (q) -| (non) node[midway, above right, font=\small\itshape] {sujets différents};
\end{tikzpicture}
\legende{Organigramme de décision : \esp{para} + infinitif ou \esp{para que} + subjonctif ?}
\end{popfigure}

\begin{propriete}[\esp{para} + nom ou pronom : la destination, le destinataire]
Devant un \textbf{nom} ou un \textbf{pronom}, \esp{para} indique le \textbf{destinataire}, la \textbf{destination} ou l'\textbf{usage prévu} :

\begin{itemize}
\item \esp{un regalo para mi madre} « un cadeau pour ma mère »
\item \esp{Este mensaje es para ti.} « Ce message est pour toi. »
\item \esp{Salgo mañana para Duala.} « Je pars demain pour Douala. » (destination)
\item \esp{Es un bolso para llevar al mercado.} « C'est un sac pour aller au marché. » (usage prévu)
\end{itemize}
\end{propriete}

\courssub{Rappel ciblé : \esp{para} ou \esp{por} ?}

Les francophones confondent souvent \esp{por} et \esp{para} car le français « pour » traduit les deux. Retenez cette distinction fondamentale : \esp{para} regarde \textbf{vers l'avant} (but, destination, destinataire), \esp{por} regarde \textbf{vers l'arrière} (cause, moyen, échange).

\begin{center}
\begin{tabularx}{\linewidth}{|X|X|X|}
\hline
\rowcolor{popblueL} \textbf{Valeur} & \textbf{Préposition} & \textbf{Exemple} \\
\hline
But, objectif & \esp{para} & \esp{Estudio para aprender.} « J'étudie pour apprendre. » \\
\hline
Destinataire & \esp{para} & \esp{un regalo para ti} « un cadeau pour toi » \\
\hline
Destination & \esp{para} & \esp{Salgo para España.} « Je pars pour l'Espagne. » \\
\hline
Cause, motif & \esp{por} & \esp{Lo hago por ti.} « Je le fais à cause de toi / pour toi. » \\
\hline
Échange, prix & \esp{por} & \esp{Lo compré por veinte mil francos.} « Je l'ai acheté pour vingt mille francs. » \\
\hline
Moyen & \esp{por} & \esp{Pago por teléfono.} « Je paie par téléphone. » \\
\hline
Durée & \esp{por} & \esp{Trabajé por dos horas.} « J'ai travaillé pendant deux heures. » \\
\hline
\end{tabularx}
\end{center}

\begin{methode}[Comment traduire « pour » en espagnol ?]
Suivez ces trois étapes dans l'ordre :
\begin{enumerate}
\item « Pour » est-il suivi d'un \textbf{infinitif} (but) ? → \esp{para} + infinitif. Exemple : « pour acheter » → \esp{para comprar}.
\item « Pour » est-il suivi d'un \textbf{nom de destinataire} ou d'une destination ? → \esp{para}. Exemple : « pour ma sœur » → \esp{para mi hermana}.
\item « Pour » exprime-t-il un \textbf{prix}, un \textbf{échange}, une \textbf{durée} ou une \textbf{cause} ? → \esp{por}. Exemple : « Je l'ai acheté pour vingt mille francs » → \esp{Lo compré por veinte mil francos}.
\end{enumerate}
Cas particulier : si « pour » est suivi d'une \textbf{proposition complète} avec un sujet différent (« pour que... »), utilisez \esp{para que} + subjonctif.
\end{methode}

\courssub{La condition avec \esp{si}}

\begin{propriete}[\esp{si} + présent de l'indicatif : condition réelle]
Pour exprimer une condition \textbf{réelle ou probable}, l'espagnol utilise :

\begin{center}
\textbf{\esp{si} + PRÉSENT DE L'INDICATIF, puis FUTUR (ou présent) dans la principale}
\end{center}

\begin{itemize}
\item \esp{Si me hace un descuento, lo compro.} « Si vous me faites une remise, je l'achète. »
\item \esp{Si tengo dinero, compraré el bolso.} « Si j'ai de l'argent, j'achèterai le sac. »
\item \esp{Si el producto es de calidad, los clientes vuelven.} « Si le produit est de qualité, les clients reviennent. »
\end{itemize}

\textbf{Bonne nouvelle :} cette construction est \textbf{identique au français} (« si + présent, futur »). Vous pouvez traduire presque mot à mot.
\end{propriete}

\begin{propriete}[\esp{si} + imparfait du subjonctif : condition irréelle (complément Bac)]
Pour exprimer une condition \textbf{irréelle ou hypothétique} au présent, l'espagnol utilise :

\begin{center}
\textbf{\esp{si} + IMPARFAIT DU SUBJONCTIF, puis CONDITIONNEL dans la principale}
\end{center}

\begin{itemize}
\item \esp{Si tuviera dinero, compraría más.} « Si j'avais de l'argent, j'achèterais davantage. »
\item \esp{Si fuera rico, ayudaría a los artesanos.} « Si j'étais riche, j'aiderais les artisans. »
\end{itemize}

Rappel de formation de l'imparfait du subjonctif : radical de la troisième personne du pluriel du prétérit + \esp{-ra, -ras, -ra, -ramos, -rais, -ran}. Exemple : \esp{tener} → \esp{tuvieron} → \esp{tuviera, tuvieras, tuviera...}
\end{propriete}

\begin{popfigure}
\begin{tikzpicture}[font=\small]
\draw[-{Stealth[length=3mm]}, thick, popdark] (0,0) -- (12,0);
\node[rectangle, rounded corners, draw=popblue, fill=popblueL, text width=4.8cm, align=center, anchor=south] at (3,0.3) {\textbf{Condition réelle}\\ \esp{Si + présent} $\rightarrow$ \textbf{futur}\\ \esp{Si tengo dinero, compraré.}};
\node[rectangle, rounded corners, draw=poppurple, fill=poppurpleL, text width=4.8cm, align=center, anchor=south] at (9,0.3) {\textbf{Condition irréelle}\\ \esp{Si + imparfait subj.} $\rightarrow$ \textbf{conditionnel}\\ \esp{Si tuviera dinero, compraría.}};
\node[below, font=\small\itshape] at (3,-0.15) {probable : ça peut arriver};
\node[below, font=\small\itshape] at (9,-0.15) {hypothétique : ce n'est pas le cas};
\end{tikzpicture}
\legende{Frise des constructions avec \esp{si} : deux schémas à connaître}
\end{popfigure}

\begin{attention}[Interdiction absolue : jamais de futur ni de conditionnel après \esp{si} !]
Comme en français, on ne met \textbf{jamais} de futur ni de conditionnel directement après \esp{si} dans une phrase conditionnelle :
\begin{itemize}
\item Faux : *\esp{Si tendré dinero, compraré el bolso.}
\item Faux : *\esp{Si tendría dinero, compraría más.}
\item Correct : \esp{Si tengo dinero, compraré el bolso.}
\item Correct : \esp{Si tuviera dinero, compraría más.}
\end{itemize}
Erreur typique du francophone : calquer le mauvais français « si j'aurais ». En espagnol, l'hypothèse se dit avec l'\textbf{imparfait du subjonctif} : \esp{Si tuviera dinero...} (« Si j'avais de l'argent... »). Retenez aussi que \esp{si} (conjonction, « si ») ne prend \textbf{jamais d'accent}, contrairement à \esp{sí} (« oui »).
\end{attention}

\begin{exemplebox}[Exemples traduits à mémoriser]
\begin{itemize}
\item \esp{Trabaja para ganarse la vida.} « Il travaille pour gagner sa vie. »
\item \esp{Compro productos locales para apoyar la economía del barrio.} « J'achète des produits locaux pour soutenir l'économie du quartier. »
\item \esp{El comerciante sonríe para que los clientes se sientan bienvenidos.} « Le commerçant sourit pour que les clients se sentent les bienvenus. »
\item \esp{Si el producto es de calidad, los clientes vuelven.} « Si le produit est de qualité, les clients reviennent. »
\item \esp{Si los precios suben, compraremos menos.} « Si les prix augmentent, nous achèterons moins. »
\item \esp{Si tuviera más dinero, compraría productos ecológicos.} « Si j'avais plus d'argent, j'achèterais des produits écologiques. »
\end{itemize}
\end{exemplebox}

\begin{exoresolu}[Traduction guidée : thème]
Traduisez en espagnol : « J'économise pour acheter un ordinateur, et mes parents m'aident pour que je puisse étudier. Si j'ai assez d'argent, je l'achèterai au marché de Douala. »
\tcblower
\textbf{Solution détaillée, étape par étape :}
\begin{enumerate}
\item « pour acheter » : infinitif après « pour », même sujet (« j'économise » et « j'achète ») → \esp{para comprar}.
\item « pour que je puisse étudier » : sujets différents (mes parents aident, c'est moi qui étudie) → \esp{para que} + subjonctif : \esp{para que yo pueda estudiar}. Attention au subjonctif irrégulier de \esp{poder} : \esp{pueda}.
\item « Si j'ai assez d'argent » : condition réelle → \esp{si} + présent : \esp{Si tengo suficiente dinero}. Jamais de futur après \esp{si}.
\item Principale au futur : \esp{lo compraré}.
\end{enumerate}
\textbf{Traduction complète :} \esp{Ahorro para comprar un ordenador, y mis padres me ayudan para que yo pueda estudiar. Si tengo suficiente dinero, lo compraré en el mercado de Duala.}
\end{exoresolu}

\begin{exercice}[À vous de jouer]
\begin{enumerate}
\item Complétez avec \esp{por}, \esp{para} ou \esp{para que} : \esp{a) Compro fruta \ldots{} hacer un postre. b) Este regalo es \ldots{} mi hermano. c) Pagué mil francos \ldots{} este bolso. d) Te explico el precio \ldots{} entiendas.}
\item Traduisez : « Si le vendeur est honnête, je reviendrai. » / « Si j'étais riche, j'achèterais ce magasin. » / « Elle travaille pour que ses enfants mangent bien. »
\end{enumerate}
\end{exercice}

\begin{corrige}[Exercice « À vous de jouer »]
\begin{enumerate}
\item a) \esp{para} (but + infinitif) ; b) \esp{para} (destinataire) ; c) \esp{por} (prix, échange) ; d) \esp{para que} (sujets différents : « je » explique, « tu » comprends → subjonctif \esp{entiendas}).
\item \esp{Si el vendedor es honesto, volveré.} / \esp{Si fuera rico, compraría esta tienda.} / \esp{Ella trabaja para que sus hijos coman bien.}
\end{enumerate}
\end{corrige}

\begin{aretenir}[L'essentiel de la section]
\begin{itemize}
\item But, même sujet : \esp{para} + infinitif → \esp{Ahorro para comprar un teléfono.}
\item But, sujets différents : \esp{para que} + subjonctif → \esp{Ahorro para que mi hijo estudie.}
\item Destinataire, destination : \esp{para} + nom → \esp{un regalo para mi madre}.
\item Prix, échange, cause, moyen : \esp{por} → \esp{Lo compré por veinte mil francos.}
\item Condition réelle : \esp{si} + présent → futur : \esp{Si tengo dinero, compraré el bolso.}
\item Condition irréelle (Bac) : \esp{si} + imparfait du subjonctif → conditionnel : \esp{Si tuviera dinero, compraría más.}
\item Jamais de futur ni de conditionnel après \esp{si} ; \esp{si} ne prend jamais d'accent.
\end{itemize}
\end{aretenir}

\coursec{Traduction : thème et version sur le commerce}

Nous venons d'étudier les trois connecteurs logiques essentiels de cette leçon : \esp{porque} pour la cause, \esp{para} pour le but et \esp{si} pour la condition. Il est temps de les mettre à l'épreuve dans l'exercice le plus exigeant du programme : la \cle{traduction}. Au Bac, la version (espagnol → français) et le thème (français → espagnol) portent très souvent sur des phrases du commerce et de l'éthique construites autour de ces connecteurs. Cette section vous donne une méthode infaillible, des exemples guidés pas à pas, et vous met en garde contre les quatre pièges qui coûtent le plus de points aux candidats francophones.

\courssub{La stratégie générale : traduire une phrase à connecteur}

Avant de traduire quoi que ce soit, il faut \emph{analyser} la phrase. La plupart des erreurs de thème ne viennent pas du vocabulaire, mais d'un mauvais choix de structure. Voici la démarche professionnelle du traducteur, adaptée au niveau Première.

\begin{methode}[Traduire une phrase à connecteur en cinq étapes]
\begin{enumerate}
\item \textbf{Souligner le connecteur} de la phrase de départ : \esp{porque}, \esp{para}, \esp{si}, « parce que », « pour », « pour que », « si »...
\item \textbf{Identifier sa valeur logique} : cause, but ou condition ? Attention, en français « pour » et « si » sont ambigus : il faut regarder le sens, pas le mot.
\item \textbf{Compter les sujets} : la subordonnée a-t-elle le même sujet que la principale ? Cette question décide du mode (infinitif ou subjonctif).
\item \textbf{Choisir la structure espagnole} et conjuguer le verbe au bon mode et au bon temps.
\item \textbf{Relire} la phrase traduite : accents, ponctuation espagnole (¿ ¡), ordre des mots, accords en genre et en nombre.
\end{enumerate}
\end{methode}

\begin{popfigure}
\begin{center}
\begin{tikzpicture}[
  box/.style={draw=popblue, fill=popblueL, rounded corners=3pt, align=center, font=\small, inner sep=4pt},
  res/.style={draw=popgreen, fill=popgreenL, rounded corners=3pt, align=center, font=\small, inner sep=4pt},
  arr/.style={-{Stealth[length=2.5mm]}, thick, popdark}]
\node[box, align=center] (a) at (0,3.6){parce que\\(en tête de phrase)};
\node[box, align=center] (b) at (0,2.4){parce que\\(après la principale)};
\node[box, align=center] (c) at (0,1.2){pour + infinitif\\(même sujet)};
\node[box, align=center] (d) at (0,0){pour que + verbe\\(sujets différents)};
\node[box] (e) at (0,-1.2) {si + condition réelle};
\node[res] (ra) at (7.5,3.6) {\esp{Como} + indicatif};
\node[res] (rb) at (7.5,2.4) {\esp{porque} + indicatif};
\node[res] (rc) at (7.5,1.2) {\esp{para} + infinitif};
\node[res] (rd) at (7.5,0) {\esp{para que} + \textbf{subjonctif}};
\node[res, align=center] (re) at (7.5,-1.2){\esp{si} + présent ;\\présent ou futur\\dans la principale};
\draw[arr] (a) -- (ra);
\draw[arr] (b) -- (rb);
\draw[arr] (c) -- (rc);
\draw[arr] (d) -- (rd);
\draw[arr] (e) -- (re);
\end{tikzpicture}
\end{center}
\legende{Tableau de correspondances : du connecteur français à la structure espagnole.}
\end{popfigure}

\begin{attention}[Après \esp{si}, jamais de futur ni de conditionnel]
En espagnol comme en français, la proposition introduite par \esp{si} (condition) se met au \textbf{présent de l'indicatif} ; c'est la proposition \emph{principale} qui porte le présent ou le futur : \esp{Si pagas al contado, el vendedor te hará una rebaja.} « Si tu paies comptant, le vendeur te fera une réduction. » La forme \esp{si pagarás} est une faute grave.
\end{attention}

\courssub{Version guidée : trois phrases sur l'éthique professionnelle}

Traduisons ensemble trois phrases espagnoles typiques du thème « commerce et éthique ». Pour chacune, nous justifions les choix de traduction.

\textbf{Phrase 1.} \esp{Si el vendedor cumple su palabra, el cliente confía en él.}

Analyse : \esp{si} introduit une condition réelle ; les deux verbes sont au présent de l'indicatif. Le français reproduit exactement ce schéma : « si + présent, présent ». Traduction : « Si le vendeur tient sa parole, le client lui fait confiance. » Justifications : \esp{cumplir su palabra} se traduit par la locution française « tenir sa parole » (et non « accomplir sa parole », calque à éviter) ; \esp{confiar en alguien} = « faire confiance à quelqu'un », d'où « lui fait confiance » grâce au pronom « lui » qui reprend \esp{en él}.

\textbf{Phrase 2.} \esp{Como este comerciante es honrado, tiene muchos clientes fieles.}

Analyse : \esp{como} en tête de phrase exprime la cause (« puisque, comme »). Traduction : « Comme ce commerçant est honnête, il a beaucoup de clients fidèles. » Justifications : \esp{honrado} signifie « honnête, intègre » (sens moral), alors que \esp{honesto} insiste sur la sincérité ; \esp{fieles} = « fidèles », adjectif placé après le nom comme en français.

\textbf{Phrase 3.} \esp{Trabajo mucho para que mis hijos tengan un futuro mejor.}

Analyse : \esp{para que} + subjonctif (\esp{tengan}) exprime le but avec changement de sujet. Le français possède la même structure : « pour que » + subjonctif. Traduction : « Je travaille beaucoup pour que mes enfants aient un avenir meilleur. » Justification : \esp{tengan} est le subjonctif présent de \esp{tener} ; le français rend naturellement ce subjonctif par « aient ».

\begin{exemplebox}[Deux exemples de référence à mémoriser]
\esp{Este vendedor engaña a sus clientes porque no es honrado.} « Ce vendeur trompe ses clients parce qu'il n'est pas honnête. »

\esp{Ahorro para que mis hijos tengan un futuro mejor.} « J'économise pour que mes enfants aient un avenir meilleur. »
\end{exemplebox}

\courssub{Thème guidé : quatre phrases françaises à traduire}

Appliquons maintenant la méthode en cinq étapes à quatre phrases de thème mobilisant « parce que », « pour » et « si ».

\begin{exoresolu}[Thème guidé : quatre phrases sur le commerce]
Traduire en espagnol :
\begin{enumerate}
\item « Parce qu'il est honnête, il a beaucoup de clients. »
\item « J'économise pour acheter un vélo. »
\item « Ce commerçant baisse ses prix pour que les pauvres puissent acheter. »
\item « Si tu paies comptant, le vendeur te fait une réduction. »
\end{enumerate}
\tcblower
\textbf{1.} Connecteur : « parce que » \emph{en tête de phrase} → règle : en espagnol, on emploie \esp{como} + indicatif, jamais \esp{porque} en tête. → \esp{Como es honrado, tiene muchos clientes.}

\textbf{2.} « pour + infinitif », même sujet (j'économise / j'achète) → \esp{para} + infinitif. → \esp{Ahorro para comprar una bicicleta.}

\textbf{3.} « pour que » + subordonnée, sujets différents (le commerçant / les pauvres) → \esp{para que} + subjonctif. Subjonctif présent de \esp{poder} : \esp{puedan}. → \esp{Este comerciante baja sus precios para que los pobres puedan comprar.}

\textbf{4.} « si » + présent, condition réelle → \esp{si} + présent en espagnol, présent dans la principale. → \esp{Si pagas al contado, el vendedor te hace una rebaja.}
\end{exoresolu}

\courssub{Les quatre pièges du thème sur le commerce}

\begin{attention}[Piège 1 : « parce que » en tête de phrase]
En français, on peut commencer une phrase par « parce que ». En espagnol, c'est \textbf{impossible} avec \esp{porque} : il faut \esp{como} + indicatif.

« Parce qu'il est honnête, il a beaucoup de clients » → \esp{Como es honrado, tiene muchos clientes.}

\esp{Porque es honrado, tiene muchos clientes} est une faute de construction. Retenez : \esp{porque} répond toujours à la question \esp{¿por qué?} et se place \emph{après} la proposition principale.
\end{attention}

\begin{attention}[Piège 2 : « pour que » exige le subjonctif]
Quand les deux verbes ont des sujets différents, « pour que » se traduit par \esp{para que} + \textbf{subjonctif}, jamais par \esp{para} + infinitif.

« J'économise pour que mes enfants aient un avenir meilleur » → \esp{Ahorro para que mis hijos tengan un futuro mejor.}

La phrase \esp{Ahorro para tener un futuro mejor} est correcte, mais elle signifie autre chose : c'est \emph{moi} qui aurai un avenir meilleur. Vérifiez toujours le nombre de sujets (étape 3 de la méthode).
\end{attention}

\begin{attention}[Piège 3 : « si » dans l'interrogation indirecte]
Dans « je voudrais savoir si... », le « si » n'exprime pas une condition mais une question indirecte : le verbe qui suit est à l'\textbf{indicatif} en espagnol. Le conditionnel français de politesse « je voudrais » se rend par \esp{quisiera} (imparfait du subjonctif de \esp{querer}, employé comme conditionnel de politesse).

« Je voudrais savoir si vous avez ce produit » → \esp{Quisiera saber si tiene este producto.}

Ne dites jamais \esp{quisiera saber si tendría...} : après \esp{si} interrogatif, on garde l'indicatif (présent ou passé selon le sens).
\end{attention}

\begin{attention}[Piège 4 : les faux amis du commerce et de l'éthique]
\begin{itemize}
\item \esp{asistir (a)} = « assister à » au sens d'« être présent » : \esp{Asisto a la reunión.} « J'assiste à la réunion. » Pour « aider quelqu'un », dites \esp{ayudar a alguien}.
\item \esp{el éxito} = « le succès » : \esp{tener éxito} = « réussir, avoir du succès ». Ne pas confondre avec « exister » (\esp{existir}).
\item \esp{actual} = « actuel, présent » (qui existe en ce moment). Pour « réel, véritable », dites \esp{real} ou \esp{verdadero} : \esp{el precio real} « le prix réel ».
\end{itemize}
\end{attention}

\begin{attention}[Bonus : « à cause de » exige un nom]
La locution \esp{a causa de} doit être suivie d'un \textbf{nom}, pas d'une phrase avec verbe conjugué. \esp{A causa de llueve} est une faute grave. Dites : \esp{Como llueve, la clienta no viene.} (« Parce qu'il pleut, la cliente ne vient pas ») ou \esp{A causa de la lluvia, la clienta no viene.} (« À cause de la pluie, la cliente ne vient pas »).
\end{attention}

\courssub{Texte d'entraînement : une plainte au marché}

Lisons ce court texte original : il rassemble tous les connecteurs de la leçon dans une situation camerounaise authentique. Repérez chaque connecteur et justifiez son emploi.

\begin{textebox}[En el mercado de Mokolo]
Ayer por la mañana, doña Clarisse fue al mercado de Mokolo, en Yaoundé, para comprar verduras frescas. Como llegó temprano, encontró los mejores productos. En un puesto, un vendedor le dijo: « Si compra tres kilos de tomates, le hago una rebaja ». Doña Clarisse aceptó porque el precio era justo y porque el vendedor parecía honrado.

Pero en otro puesto, tuvo una mala sorpresa: el comerciante pesó mal el arroz para ganar más dinero. La clienta protestó: « ¡No es ético engañar a los clientes! Yo trabajo duro para que mi familia coma bien ». El vendedor, avergonzado, devolvió el dinero y pidió perdón.

Un vecino que asistió a la escena comentó: « Este mercado tiene éxito porque la mayoría de los vendedores son serios. Si todos respetan a los clientes, el comercio mejora para todos ». Doña Clarisse volvió a casa contenta, pero decidió vigilar siempre la balanza. Como dice su madre, la confianza es la base del comercio.
\end{textebox}

\textbf{Glossaire.} \esp{el puesto} : l'étal, le stand ; \esp{la rebaja} : la réduction ; \esp{pesar} : peser ; \esp{engañar} : tromper ; \esp{avergonzado} : honteux ; \esp{devolver} : rendre, rembourser ; \esp{la balanza} : la balance ; \esp{vigilar} : surveiller.

\textbf{Questions de compréhension.}
\begin{enumerate}
\item ¿Por qué doña Clarisse encontró los mejores productos ?
\item ¿Qué hizo el segundo comerciante para ganar más dinero ?
\item Según el vecino, ¿por qué tiene éxito el mercado ?
\end{enumerate}

\begin{corrige}[Questions de compréhension]
1. \esp{Encontró los mejores productos porque llegó temprano al mercado.} (cause exprimée par \esp{como} en tête dans le texte, reprise ici avec \esp{porque} après la principale.)

2. \esp{Pesó mal el arroz para ganar más dinero.} (but : \esp{para} + infinitif, même sujet.)

3. \esp{Según el vecino, el mercado tiene éxito porque la mayoría de los vendedores son serios.} (faux ami : \esp{tener éxito} = « avoir du succès ».)
\end{corrige}

\courssub{Tableau récapitulatif et entraînement final}

\begin{center}
\begin{tabularx}{\linewidth}{|X|X|X|}
\hline
\rowcolor{popblueL} \textbf{Français} & \textbf{Español} & \textbf{Remarque} \\
\hline
parce que (après la principale) & \esp{porque} + indicatif & répond à « pourquoi ? » \\
\hline
parce que (en tête de phrase) & \esp{como} + indicatif & jamais \esp{porque} en tête \\
\hline
pour + infinitif (même sujet) & \esp{para} + infinitif & \esp{Ahorro para comprar...} \\
\hline
pour que + subordonnée & \esp{para que} + subjonctif & sujets différents obligatoires \\
\hline
si + condition réelle & \esp{si} + présent & présent ou futur dans la principale \\
\hline
si (question indirecte) & \esp{si} + indicatif & \esp{Quisiera saber si...} \\
\hline
à cause de + nom & \esp{a causa de} + nom & sinon \esp{porque} ou \esp{como} + phrase \\
\hline
\end{tabularx}
\end{center}

\begin{exoresolu}[Entraînement final : thème et version]
\textbf{A. Version.} Traduire : \esp{Como los precios actuales son altos, muchas familias compran menos para ahorrar.}

\textbf{B. Thème.} Traduire : « Si le vendeur est honnête, les clients reviennent, parce que la confiance est importante. »
\tcblower
\textbf{A.} « Comme les prix actuels sont élevés, beaucoup de familles achètent moins pour économiser. » Justifications : \esp{como} en tête = « comme / parce que » ; \esp{actuales} = « actuels » (faux ami : pas « réels ») ; \esp{para ahorrar} = « pour économiser » (but, même sujet, infinitif).

\textbf{B.} \esp{Si el vendedor es honrado, los clientes vuelven, porque la confianza es importante.} Justifications : condition réelle → \esp{si} + présent, présent dans la principale ; « parce que » placé après la principale → \esp{porque} (et non \esp{como}) ; « revenir » = \esp{volver} (verbe à changement radical o → ue : \esp{vuelven}).
\end{exoresolu}

\begin{savaistu}[Les connecteurs logiques au Bac]
Au baccalauréat A, les phrases de thème portent très souvent sur les connecteurs logiques : cause, but, condition. Un candidat qui maîtrise parfaitement le trio \esp{porque / para / si} — avec leurs pièges : \esp{como} en tête, \esp{para que} + subjonctif, \esp{si} sans futur ni conditionnel — sécurise plusieurs points sur l'épreuve sans effort supplémentaire. C'est un investissement rentable : trois règles à retenir, des points garantis à chaque session.
\end{savaistu}

Vous disposez désormais de tous les outils pour traduire avec précision les phrases du commerce et de l'éthique. La dernière section vous conduira à mobiliser ces mêmes connecteurs dans une production personnelle : la rédaction d'un court texte sur vos propres habitudes de consommation.

\coursec{Production écrite guidée : mes habitudes de consommation}

Dans la section précédente, nous avons traduit des phrases sur le commerce et l'éthique, du français vers l'espagnol et inversement. Traduire, c'est déjà produire : mais les mots étaient imposés. Nous franchissons maintenant la dernière étape du parcours : \cle{rédiger librement} un court texte personnel, en réutilisant tout ce que cette leçon nous a appris — le lexique du commerce et de l'éthique, les connecteurs \esp{porque}, \esp{para} et \esp{si}. C'est exactement l'exercice qui vous attendra au baccalauréat : une production écrite brève, claire et correcte. Apprenons à viser juste.

\courssub{La consigne type Bac}

\begin{definition}[La consigne]
Au baccalauréat, la production écrite se présente souvent ainsi : \esp{¿Redacta un texto de unas 80 palabras titulado « Mis hábitos de consumo ». Habla de lo que compras, dónde, por qué, y de tu compromiso como consumidor responsable!} — « Rédige un texte d'environ 80 mots intitulé "Mes habitudes de consommation". Parle de ce que tu achètes, où, pourquoi, et de ton engagement de consommateur responsable. »
\end{definition}

Trois exigences sont cachées dans cette consigne : un \cle{titre} imposé (il faut le recopier tel quel), un \cle{nombre de mots} à respecter (environ 80), et un \cle{contenu} précis (quoi, où, pourquoi, engagement). Le correcteur vérifie les trois.

\courssub{Le plan en entonnoir : du général à l'engagement}

Un bon texte de 80 mots n'est pas une liste d'idées : c'est un \cle{entonnoir}. On part du général (ses habitudes), on précise (les causes et les buts), et on termine par le plus personnel : son engagement, exprimé avec la condition.

\begin{popfigure}
\begin{center}
\begin{tikzpicture}[node distance=0.6cm]
\node[fill=popblue!25, rounded corners=4pt, align=center, text=popdark, minimum width=9cm, minimum height=1.1cm] (n1) {\textbf{1. Général : mes habitudes}\\ \esp{Normalmente compro... en el mercado...}};
\node[fill=poporange!30, rounded corners=4pt, align=center, text=popdark, minimum width=7cm, minimum height=1.1cm, below=of n1] (n2) {\textbf{2. Précis : causes et buts}\\ \esp{...porque son frescos... para apoyar a los campesinos}};
\node[fill=popgreen!30, rounded corners=4pt, align=center, text=popdark, minimum width=5cm, minimum height=1.1cm, below=of n2] (n3) {\textbf{3. Engagement : la condition}\\ \esp{Si tengo dinero, compraré artesanía.}};
\draw[-{Stealth}, very thick, popmuted] (n1.south) -- (n2.north);
\draw[-{Stealth}, very thick, popmuted] (n2.south) -- (n3.north);
\end{tikzpicture}
\end{center}
\legende{Le plan en entonnoir : général $\rightarrow$ précis $\rightarrow$ engagement}
\end{popfigure}

\begin{center}
\begin{tabularx}{\linewidth}{|l|X|X|}
\hline
\rowcolor{popblueL} \textbf{Paragraphe} & \textbf{Contenu} & \textbf{Connecteur obligatoire} \\
\hline
1 & Ce que j'achète et où & \esp{normalmente}, \esp{siempre} \\
\hline
2 & Pourquoi (cause) et pour quoi (but) & \esp{porque} + indicatif ; \esp{para} + infinitif \\
\hline
3 & Mon engagement responsable & \esp{si} + présent $\rightarrow$ présent / futur \\
\hline
\end{tabularx}
\end{center}

\courssub{La banque de phrases-modèles}

Voici des phrases de départ que vous pouvez adapter à votre propre vie. Ne les copiez pas telles quelles : changez les lieux, les produits, les raisons.

\begin{itemize}
\item \esp{Normalmente compro frutas y verduras en el mercado de mi barrio porque los productos son frescos.} « Normalement j'achète des fruits et légumes au marché de mon quartier parce que les produits sont frais. »
\item \esp{Prefiero los productos locales para apoyar a los campesinos de mi región.} « Je préfère les produits locaux pour soutenir les paysans de ma région. »
\item \esp{Si tengo suficiente dinero, compraré artesanía hecha a mano.} « Si j'ai assez d'argent, j'achèterai de l'artisanat fait main. »
\item \esp{Nunca compro productos falsificados porque es una estafa.} « Je n'achète jamais de produits contrefaits parce que c'est une arnaque. »
\item \esp{Compro ropa de segunda mano para gastar menos y proteger el medio ambiente.} « J'achète des vêtements d'occasion pour dépenser moins et protéger l'environnement. »
\item \esp{Si el vendedor es honesto, volveré a su tienda.} « Si le vendeur est honnête, je retournerai dans sa boutique. »
\end{itemize}

\begin{exemplebox}[Exemples traduits : cause, but, condition]
\esp{No compro en esa tienda porque los precios son muy altos.} « Je n'achète pas dans cette boutique parce que les prix sont très élevés. » (cause : \esp{porque} + verbe conjugué à l'indicatif)

\esp{Ahorro dinero para comprar un teléfono nuevo.} « J'économise de l'argent pour acheter un téléphone neuf. » (but : \esp{para} + infinitif)

\esp{Si los productos son de buena calidad, pagaré un poco más.} « Si les produits sont de bonne qualité, je paierai un peu plus. » (condition : \esp{si} + présent de l'indicatif, futur dans la principale)
\end{exemplebox}

\begin{attention}[Les trois fautes qui coûtent cher]
\begin{itemize}
\item Après \esp{para} (but), toujours l'\cle{infinitif} : \esp{para apoyar}, jamais \esp{para apoyo} ni \esp{para que apoyo} (ce dernier exige le subjonctif : \esp{para que apoye}, niveau supérieur).
\item Après \esp{si} (condition), \cle{jamais de futur} dans la proposition subordonnée : on dit \esp{Si tengo dinero, compraré}, jamais \esp{Si tendré dinero}. C'est la même logique qu'en français : « si j'ai », pas « si j'aurai ».
\item Attention aux faux amis : \esp{una estafa} = « une arnaque » (féminin), pas « un étage » ; \esp{los campesinos} = « les paysans », pas « les champs » (\esp{los campos}) ; \esp{la ropa} = « les vêtements » (singulier), pas « la corde ».
\end{itemize}
\end{attention}

\courssub{Le modèle de rédaction commenté}

Lisons d'abord le texte rédigé par l'enseignant, environ 80 mots (dans la fourchette 70--90), puis observons comment il est construit.

\begin{textebox}[Modèle : Mis hábitos de consumo]
Normalmente compro frutas y verduras en el mercado de Mokolo porque los productos son frescos, baratos y de temporada. Prefiero los productos locales para apoyar a los campesinos de mi región y reducir la contaminación. Si tengo suficiente dinero, compraré artesanía hecha a mano en lugar de objetos industriales. Nunca compro productos falsificados porque es una estafa que perjudica a los trabajadores. Creo firmemente que un consumidor responsable debe respetar el medio ambiente y el trabajo digno.
\end{textebox}

\textbf{Traduction.} « Normalement j'achète fruits et légumes au marché de Mokolo parce que les produits sont frais, bon marché et de saison. Je préfère les produits locaux pour soutenir les paysans de ma région et réduire la pollution. Si j'ai assez d'argent, j'achèterai de l'artisanat fait main au lieu d'objets industriels. Je n'achète jamais de produits contrefaits parce que c'est une arnaque qui nuit aux travailleurs. Je crois fermement qu'un consommateur responsable doit respecter l'environnement et le travail décent. »

\textbf{Glossaire.} \esp{de temporada} = de saison ; \esp{reducir la contaminación} = réduire la pollution ; \esp{en lugar de} = au lieu de ; \esp{perjudicar} = nuire à, porter préjudice ; \esp{trabajo digno} = travail décent ; \esp{firmemente} = fermement.

\textbf{Analyse du modèle.} Le texte contient 5 phrases (73 mots). Phrase 1 : habitude + lieu + \esp{porque} (cause multiple). Phrase 2 : préférence + \esp{para} + infinitif (but double). Phrase 3 : \esp{si} + présent, \cle{futur} (condition réaliste). Phrase 4 : refus éthique + \esp{porque} + conséquence. Phrase 5 : conclusion avec \esp{creo que} + lexique éthique (\esp{consumidor responsable}, \esp{respetar}, \esp{medio ambiente}, \esp{trabajo digno}). Toutes les contraintes sont remplies, phrases simples et correctes.

\begin{methode}[Rédiger 80 mots efficaces en 5 étapes]
\begin{enumerate}
\item \textbf{Recopie le titre} exactement : \esp{Mis hábitos de consumo}.
\item \textbf{Écris 6 à 8 phrases courtes} (10 à 12 mots chacune). Une idée principale par phrase.
\item \textbf{Mets un connecteur logique par phrase} : \esp{porque}, \esp{para}, \esp{si}, \esp{también}, \esp{pero}, \esp{creo que}, \esp{en cambio}.
\item \textbf{Vérifie les contraintes obligatoires} : au moins 1 \esp{porque} + indicatif, 1 \esp{para} + infinitif, 1 \esp{si} + présent (suivi de présent ou futur), et 3 mots du lexique de l'éthique (\esp{responsable}, \esp{honesto}, \esp{estafa}, \esp{falsificado}, \esp{respetar}, \esp{local}, \esp{digno}, \esp{justo}).
\item \textbf{Compte les mots} et relis l'orthographe (accents !) et les accords (genre, nombre). Ne copie jamais une phrase entière du texte support de la leçon : le correcteur le repérerait.
\end{enumerate}
\end{methode}

\begin{attention}[Compter les mots : 80 $\pm$ 10 \%]
Au baccalauréat, un texte trop court (moins de 70 mots) ou trop long (plus de 90 mots) est \cle{pénalisé}. Astuce : compte les mots ligne par ligne, en sachant qu'une ligne d'écriture normale contient environ 8 à 10 mots. Huit à dix lignes suffisent donc. En espagnol, pas d'apostrophe élidée : \esp{del} = 1 mot, \esp{al} = 1 mot, \esp{por qué} = 2 mots, \esp{porque} = 1 mot.
\end{attention}

\begin{savaistu}[Pourquoi 80 mots ?]
La production écrite au baccalauréat camerounais exige généralement entre 60 et 100 mots. S'entraîner à rédiger exactement 80 mots, c'est viser juste : assez long pour développer trois idées structurées, assez court pour rester correct du début à la fin. Les meilleurs candidats ne sont pas ceux qui écrivent le plus, mais ceux qui écrivent \emph{juste} et \emph{sans faute}.
\end{savaistu}

\courssub{Exercice résolu : corriger un texte d'élève}

\begin{exoresolu}[Trouve et corrige les 4 fautes]
Énoncé. Voici le début du texte d'un élève. Repère les quatre fautes et corrige-les : \esp{Normalmente compro en el mercado porque los productos es frescos. Prefiero los productos locales para apoyo a los campesinos. Si tendré dinero, compraré artesanía. Nunca compro productos falsificados porque es un estafa.}
\tcblower
Solution détaillée.
\begin{enumerate}
\item \esp{los productos es frescos} $\rightarrow$ \esp{los productos \textbf{son} frescos} : accord du verbe \esp{ser} avec un sujet pluriel (\esp{los productos}).
\item \esp{para apoyo} $\rightarrow$ \esp{para \textbf{apoyar}} : après \esp{para} (but), il faut l'infinitif, jamais le nom.
\item \esp{Si tendré dinero} $\rightarrow$ \esp{Si \textbf{tengo} dinero} : après \esp{si} (condition), jamais de futur ; le futur va dans la proposition principale (\esp{compraré}, correct ici).
\item \esp{un estafa} $\rightarrow$ \esp{\textbf{una} estafa} : \esp{estafa} est féminin (comme « une arnaque »), donc article \esp{una}.
\end{enumerate}
Texte corrigé : \esp{Normalmente compro en el mercado porque los productos son frescos. Prefiero los productos locales para apoyar a los campesinos. Si tengo dinero, compraré artesanía. Nunca compro productos falsificados porque es una estafa.}
\end{exoresolu}

\courssub{La grille d'auto-évaluation}

Avant de rendre ta copie, évalue-toi comme le ferait l'examinateur, sur 20 points.

\begin{center}
\begin{tabularx}{\linewidth}{|X|c|X|}
\hline
\rowcolor{popblueL} \textbf{Critère} & \textbf{Points} & \textbf{Questions à se poser} \\
\hline
Contenu & 4 & Ai-je parlé de quoi, où, pourquoi, et de mon engagement ? \\
\hline
Connecteurs & 4 & Ai-je au moins 1 \esp{porque}, 1 \esp{para} + infinitif, 1 \esp{si} + présent ? \\
\hline
Lexique & 4 & Ai-je 3 mots du lexique de l'éthique et du commerce ? \\
\hline
Orthographe et morphosyntaxe & 4 & Accents, accords, conjugaisons corrects ? \\
\hline
Respect des 80 mots & 4 & Entre 70 et 90 mots ? Titre recopié exactement ? \\
\hline
\textbf{Total} & \textbf{20} & \\
\hline
\end{tabularx}
\end{center}

\courssub{La phase de réécriture après correction par les pairs}

\begin{experience}[Correction croisée et réécriture]
\begin{enumerate}
\item \textbf{Rédaction} (15 min) : chaque élève rédige ses 80 mots en suivant le plan en entonnoir.
\item \textbf{Échange} (5 min) : tu donnes ta copie à ton voisin et tu reçois la sienne.
\item \textbf{Correction par les pairs} (10 min) : avec la grille ci-dessus, note le texte de ton camarade sur 20. Entoure en douceur les fautes, souligne les connecteurs trouvés, compte les mots. Écris un commentaire gentil et précis : \esp{« Buen uso de "para", pero falta un "si" »}.
\item \textbf{Réécriture} (10 min) : reprends ta copie, corrige les fautes signalées et rédige la \cle{version finale} au propre. C'est cette deuxième version, améliorée, qui sera évaluée par le professeur.
\end{enumerate}
\end{experience}

\begin{exercice}[À toi d'écrire]
Rédige ton texte \esp{Mis hábitos de consumo} (environ 80 mots). Contraintes obligatoires : 1 \esp{porque} + indicatif, 1 \esp{para} + infinitif, 1 \esp{si} + présent (suivi de présent ou futur), 3 mots du lexique de l'éthique. Contexte libre : le marché de Mokolo, une boutique de Douala, les achats de fournitures scolaires, l'artisanat camerounais, le commerce équitable\ldots{}
\end{exercice}

\begin{corrige}[Éléments de correction]
Il n'existe pas une seule bonne réponse, mais toute copie réussie doit contenir : le titre recopié exactement ; trois moments (habitudes, causes/buts, engagement) ; les trois connecteurs grammaticalement corrects ; un lexique éthique varié (\esp{responsable}, \esp{honesto}, \esp{justo}, \esp{estafa}, \esp{falsificado}, \esp{respetar}, \esp{local}, \esp{digno}) ; entre 70 et 90 mots. Exemple de phrase d'engagement attendue : \esp{Si todos compramos productos locales, nuestra economía será más fuerte y justa.} « Si nous achetons tous des produits locaux, notre économie sera plus forte et plus juste. »
\end{corrige}

\begin{aretenir}[L'essentiel de la production écrite]
Un texte de 80 mots réussi = un titre recopié + 6 à 8 phrases courtes + un connecteur par phrase (\esp{porque} pour la cause, \esp{para} + infinitif pour le but, \esp{si} + présent pour la condition) + 3 mots du lexique de l'éthique + une relecture systématique des accents et des accords. Écris simple, écris juste : c'est la clé de la réussite au baccalauréat.
\end{aretenir}

\newpage

\coursec{Activité d'intégration}

\courssub{Objectifs de l'activité}

Cette activité d'intégration mobilise, en une seule tâche de groupe, tous les savoir-faire de la leçon :

\begin{itemize}
\item jouer un \cle{dialogue client--vendeur} dans une situation de marché réaliste ;
\item employer correctement les connecteurs de \cle{cause} (\esp{porque}), de \cle{but} (\esp{para}) et de \cle{condition} (\esp{si}) ;
\item \cle{traduire} dans les deux sens : version (espagnol vers français) et thème (français vers espagnol) ;
\item \cle{rédiger} individuellement un court texte de 80 mots sur ses habitudes de consommation ;
\item développer l'esprit d'équipe, l'expression orale en public et le sens de l'\cle{éthique professionnelle}.
\end{itemize}

\courssub{Situation}

C'est samedi matin au marché Mokolo de Yaoundé. Ta classe de Première prépare la « Semaine de la consommation responsable » organisée par le lycée. Le professeur d'espagnol propose un concours : chaque groupe de trois élèves doit présenter une scène de marché en espagnol, avec des slogans éthiques traduits, et la meilleure équipe verra son affiche exposée à la bibliothèque.

Dans chaque groupe, il y a trois rôles : un \esp{cliente} (client), un \esp{vendedor} (vendeur) et un \esp{observador de la ética} (observateur de l'éthique), chargé des traductions et du commentaire.

Voici le document officiel du concours, affiché au tableau :

\begin{textebox}[Concurso « En el mercado responsable » — Reglamento]
CONCURSO DE LA SEMANA DEL CONSUMO RESPONSABLE

Cada grupo prepara una escena en el mercado con tres personajes: un cliente, un vendedor y un observador de la ética.

El diálogo debe tener diez réplicas y contener, como mínimo, dos veces \esp{porque}, dos veces \esp{para} y una vez \esp{si}.

El observador traduce al francés tres réplicas importantes y traduce al español dos eslóganes éticos.

Después, cada alumno escribe un texto de ochenta palabras: « Mis hábitos de consumo ».

¡El mejor diálogo y el mejor eslogan ganan un premio!
\end{textebox}

\begin{definition}[Glossaire du règlement]
\esp{la réplica} : la réplique ; \esp{como mínimo} : au minimum ; \esp{el eslogan} (pluriel \esp{eslóganes}) : le slogan ; \esp{los hábitos de consumo} : les habitudes de consommation ; \esp{el premio} : le prix, la récompense ; \esp{ganar} : gagner ; \esp{el personaje} : le personnage.
\end{definition}

\begin{attention}[Orthographe : \esp{eslogan} ou \esp{eslógan} ?]
Le mot \esp{eslogan} est un mot \emph{grave} (accentué sur l'avant-dernière syllabe) terminé par \emph{-n} : il ne prend donc \textbf{pas d'accent écrit} au singulier. Au pluriel, l'accent tonique tombe sur la troisième syllabe en partant de la fin : il faut alors écrire \esp{eslóganes} avec un accent. Même phénomène que \esp{joven} / \esp{jóvenes}.
\end{attention}

\courssub{Consignes}

\begin{experience}[Étape 1 — Écrire le dialogue (20 min)]
Par groupes de trois, rédigez un dialogue de \cle{dix répliques} au marché de votre ville (Mokolo, marché central de Douala, marché de Bafoussam...). Le client demande des informations, négocie le prix et pose des questions sur la qualité ; le vendeur répond avec honnêteté. Contraintes obligatoires :
\begin{enumerate}
\item employer au moins \textbf{deux fois} \esp{porque} (cause) ;
\item employer au moins \textbf{deux fois} \esp{para} + infinitif (but) ;
\item employer au moins \textbf{une fois} \esp{si} + présent (condition) ;
\item utiliser au moins cinq mots du lexique du commerce : \esp{el precio, la calidad, el descuento, la factura, el cliente, honesto, caro, barato}.
\end{enumerate}
\end{experience}

\begin{experience}[Étape 2 — Traduire (15 min)]
L'observateur de l'éthique :
\begin{enumerate}
\item choisit \cle{trois répliques clés} du dialogue et les traduit en français (version) ;
\item traduit en espagnol \cle{deux slogans éthiques} (thème), par exemple : « Un vendeur honnête garde ses clients » et « J'achète local pour soutenir les producteurs de ma région ».
\end{enumerate}
Attention aux pièges : \esp{honesto} se traduit par « honnête », et « soutenir » se dit \esp{apoyar} ou \esp{sostener}, jamais \esp{soportar}.
\end{experience}

\begin{experience}[Étape 3 — Jouer la scène et afficher (15 min)]
Chaque groupe joue sa scène devant la classe (deux minutes maximum), puis colle ses slogans traduits au tableau sous forme d'affiche. La classe écoute et note sur une grille : prononciation, correction des connecteurs, naturel du jeu.
\end{experience}

\begin{experience}[Étape 4 — Rédiger le bilan individuel (15 min)]
Chacun rédige seul un texte de \cle{80 mots} intitulé \esp{Mis hábitos de consumo}, en réutilisant au moins un \esp{porque}, un \esp{para} et un \esp{si}. Terminez par un vote de la classe : le dialogue le plus convaincant et le slogan le mieux traduit reçoivent le prix de la « Semaine de la consommation responsable ».
\end{experience}

\courssub{Ressources à mobiliser}

\begin{aretenir}[Connecteurs à employer]
\begin{itemize}
\item \textbf{Cause} : \esp{porque} + verbe conjugué. \esp{No compro esta camisa porque es muy cara.} « Je n'achète pas cette chemise parce qu'elle est très chère. »
\item \textbf{But} : \esp{para} + infinitif. \esp{Vengo al mercado para comprar fruta fresca.} « Je viens au marché pour acheter des fruits frais. »
\item \textbf{Condition} : \esp{si} + présent, jamais de futur après \esp{si}. \esp{Si bajas el precio, compro dos kilos.} « Si tu baisses le prix, j'achète deux kilos. »
\end{itemize}
\end{aretenir}

\begin{propriete}[Rappel : les trois structures en un coup d'œil]
\begin{center}
\begin{tabularx}{\linewidth}{|X|X|X|}
\hline
\rowcolor{popblueL} Relation logique & Structure & Exemple \\
\hline
Cause & \esp{porque} + verbe conjugué & \esp{...porque son frescos.} \\
\hline
But & \esp{para} + infinitif & \esp{...para ahorrar dinero.} \\
\hline
Condition & \esp{si} + présent (+ présent ou futur dans l'autre proposition) & \esp{Si tengo dinero, lo compro.} \\
\hline
\end{tabularx}
\end{center}
Attention à ne pas confondre \esp{si} (si, conjonction de condition, sans accent) et \esp{sí} (oui, adverbe, avec accent).
\end{propriete}

\begin{aretenir}[Lexique utile du marché]
\esp{el precio} (le prix), \esp{la calidad} (la qualité), \esp{el descuento} (la remise), \esp{la factura/el recibo} (la facture/le reçu), \esp{caro/barato} (cher/bon marché), \esp{regatear} (marchander), \esp{honesto} (honnête), \esp{el productor local} (le producteur local), \esp{gastar} (dépenser), \esp{ahorrar} (économiser), \esp{¿Cuánto cuesta?} (Combien ça coûte ?), \esp{Es demasiado caro.} (C'est trop cher.)
\end{aretenir}

\begin{attention}[Pièges de traduction]
\begin{itemize}
\item « Un vendeur honnête \textbf{garde} ses clients » : \esp{Un vendedor honesto \textbf{conserva/mantiene} a sus clientes.} N'oubliez pas la \cle{préposition a} devant le complément de personne (\esp{a sus clientes}) : c'est l'\emph{a} personnelle, obligatoire en espagnol.
\item « Pour soutenir » : \esp{para apoyar}, et non \esp{soportar} (qui signifie « supporter, tolérer ») : faux ami classique.
\item Après \esp{si}, pas de futur : on dit \esp{si tengo dinero}, jamais \esp{si tendré dinero}. Le futur se place dans l'autre proposition : \esp{Si tengo dinero, lo compraré.} « Si j'ai de l'argent, je l'achèterai. »
\item « Combien coûtent les avocats ? » : \esp{¿Cuánto cuestan los aguacates?} — le verbe \esp{costar} s'accorde avec la chose qui coûte, pas avec la personne.
\end{itemize}
\end{attention}

\courssub{Proposition de production et corrigé commenté}

\begin{exoresolu}[Modèle de dialogue, traductions et bilan]
\textbf{Énoncé.} Produire le dialogue de dix répliques, les traductions demandées et le bilan de 80 mots.
\tcblower
\textbf{Diálogo modelo (10 réplicas) :}

\esp{-- Buenos días, señor. ¿Cuánto cuestan los aguacates?}

\esp{-- Buenos días. Cuestan quinientos francos porque son de muy buena calidad.}

\esp{-- Es un poco caro. Vengo al mercado para comprar productos frescos y baratos.}

\esp{-- Le hago un descuento si compra tres kilos.}

\esp{-- ¿De dónde vienen estos aguacates?}

\esp{-- Vienen de un productor local porque quiero apoyar a los campesinos de mi región.}

\esp{-- Muy bien. Compro productos locales para ayudar a la economía de mi país.}

\esp{-- Usted es un cliente responsable. Si vuelve la semana próxima, le regalo un aguacate.}

\esp{-- Gracias, señor. Vuelvo porque usted es un vendedor honesto.}

\esp{-- Aquí tiene su recibo. ¡Hasta pronto!}

\textbf{Version (3 répliques traduites en français) :}

\esp{Cuestan quinientos francos porque son de muy buena calidad.} « Ils coûtent cinq cents francs parce qu'ils sont de très bonne qualité. »

\esp{Le hago un descuento si compra tres kilos.} « Je vous fais une remise si vous achetez trois kilos. »

\esp{Vuelvo porque usted es un vendedor honesto.} « Je reviens parce que vous êtes un vendeur honnête. »

\textbf{Thème (2 slogans traduits en espagnol) :}

« Un vendeur honnête garde ses clients. » \esp{Un vendedor honesto conserva a sus clientes.}

« J'achète local pour soutenir les producteurs de ma région. » \esp{Compro productos locales para apoyar a los productores de mi región.}

\textbf{Bilan individuel (80 mots) :}

\esp{Mis hábitos de consumo}

\esp{Compro fruta y verdura en el mercado porque los productos son frescos y baratos. Voy al mercado para apoyar a los productores locales de mi región. Si tengo poco dinero, compro solamente lo necesario y no gasto en cosas inútiles. Siempre pregunto el precio antes de comprar porque quiero ahorrar. Creo que un consumidor responsable elige la calidad y respeta a los vendedores honestos. Consumir bien es importante para mi familia y para mi país.}

\textbf{Commentaire des choix de langue :}
\begin{enumerate}
\item \esp{porque} est suivi d'un verbe conjugué (\esp{son}, \esp{quiero}) : c'est la structure de la cause.
\item \esp{para} est toujours suivi de l'infinitif (\esp{apoyar}, \esp{comprar}) : c'est le but, équivalent de « pour ».
\item Dans \esp{si tengo poco dinero, compro...}, le présent suit \esp{si} : jamais de futur dans la proposition conditionnelle, comme en français.
\item L'\emph{a} personnelle apparaît dans \esp{apoyar a los productores} et \esp{conserva a sus clientes} devant des compléments de personne.
\item Le vouvoiement du client se fait avec \esp{usted} et la troisième personne : \esp{le hago}, \esp{si compra}, \esp{aquí tiene}.
\item Le lexique du commerce est bien réinvesti : \esp{precio, descuento, recibo, calidad, honesto, ahorrar}.
\end{enumerate}
\end{exoresolu}

\begin{methode}[Comment réussir le texte de 80 mots à l'examen]
\begin{enumerate}
\item \textbf{Préparer un plan rapide} : quoi j'achète, où, pourquoi, avec qui, et mon opinion.
\item \textbf{Placer les trois connecteurs} : une cause avec \esp{porque}, un but avec \esp{para} + infinitif, une condition avec \esp{si} + présent.
\item \textbf{Réutiliser le lexique de la leçon} : \esp{precio, calidad, mercado, ahorrar, gastar, productor local}.
\item \textbf{Employer le présent de l'indicatif} : c'est le temps des habitudes ; éviter les temps non maîtrisés.
\item \textbf{Compter les mots et relire} : vérifier les accents, l'accord des adjectifs (\esp{productos frescos y baratos}) et les points d'exclamation et d'interrogation ouvrants (¡ ¿).
\end{enumerate}
\end{methode}

\courssub{Bilan de l'activité}

À l'issue de cette activité, tu as mobilisé l'ensemble des compétences de la leçon dans une tâche authentique et concrète :

\begin{itemize}
\item tu sais \cle{conduire un dialogue} client--vendeur : saluer, demander le prix, marchander, remercier ;
\item tu maîtrises les trois connecteurs essentiels : \esp{porque} (cause), \esp{para} + infinitif (but), \esp{si} + présent (condition) ;
\item tu as pratiqué la \cle{traduction} dans les deux sens, en évitant les faux amis (\esp{soportar}/soutenir) et en respectant l'\emph{a} personnelle ;
\item tu as rédigé un texte personnel de 80 mots sur tes \cle{habitudes de consommation}, compétence directement exigible à l'évaluation ;
\item tu as développé une attitude de \cle{consommateur responsable} : acheter local, demander un reçu, respecter les vendeurs honnêtes.
\end{itemize}

\begin{aretenir}[L'essentiel]
Une phrase de cause : \esp{porque} + verbe conjugué. Une phrase de but : \esp{para} + infinitif. Une phrase de condition : \esp{si} + présent. Avec ces trois outils et le lexique du marché, tu peux dialoguer, traduire et rédiger sur la consommation responsable : c'est exactement ce que l'on te demandera en expression écrite et en traduction.
\end{aretenir}

\newpage

\coursec{Méthodes et exercices résolus}

\coursec{Leçon EA16 : El consumo responsable y la ética profesional (continuación) : expresión escrita}

\courssub{Texte support : Dialogue au marché de Mokolo}

\begin{textebox}[Dialogue authentique : achat de tomates au marché]
\esp{-- Vendedor : ¡Buenos días, señora! ¿Qué desea llevar hoy?} \\
\esp{-- Cliente : Buenos días. Quisiera dos kilos de tomates, por favor.} \\
\esp{-- Vendedor : Aquí tiene, son muy frescos, acaban de llegar de la finca.} \\
\esp{-- Cliente : ¿Cuánto cuesta el kilo?} \\
\esp{-- Vendedor : Quinientos francos CFA. Pero si lleva tres kilos, le dejo el kilo a cuatrocientos.} \\
\esp{-- Cliente : De acuerdo, los compro porque quiero hacer una salsa para toda la familia.} \\
\esp{-- Vendedor : Muy bien, son tres kilos a cuatrocientos: mil doscientos francos en total.} \\
\esp{-- Cliente : Aquí tiene. Muchas gracias, hasta luego.} \\
\esp{-- Vendedor : Gracias a usted, que le aproveche.}
\end{textebox}

\begin{center}
\begin{tabularx}{\linewidth}{|X|X|}
\hline
\rowcolor{popblueL} \textbf{Español} & \textbf{Français} \\
\hline
\esp{la finca} & la plantation, l'exploitation agricole \\
\esp{acaban de llegar} & ils viennent d'arriver (venir de + infinitif) \\
\esp{el kilo} & le kilo \\
\esp{dejar a buen precio} & laisser à bon prix, faire un prix \\
\esp{que le aproveche} & bon appétit / profitez-en (formule de politesse) \\
\hline
\end{tabularx}
\end{center}

\courssub{Lexique : Le commerce et l'éthique}

\begin{definition}[Vocabulaire essentiel du module 4]
\begin{center}
\begin{tabularx}{\linewidth}{|X|X|X|}
\hline
\rowcolor{popblueL} \textbf{Noms / Sustantivos} & \textbf{Verbes / Verbos} & \textbf{Adjectifs / Adjetivos} \\
\hline
\esp{el comercio} (commerce) & \esp{comprar} (acheter) & \esp{barato/a} (bon marché) \\
\esp{el mercado} (marché) & \esp{vender} (vendre) & \esp{caro/a} (cher) \\
\esp{la tienda} (magasin) & \esp{costar} (coûter) & \esp{fresco/a} (frais) \\
\esp{el vendedor / la vendedora} & \esp{regatear} (marchander) & \esp{justo/a} (juste, équitable) \\
\esp{el cliente / la clienta} & \esp{ahorrar} (économiser) & \esp{honrado/a} (honnête) \\
\esp{el precio} (prix) & \esp{engañar} (tromper) & \esp{responsable} (responsable) \\
\esp{el descuento} (remise) & \esp{confiar} (faire confiance) & \esp{local/a} (local) \\
\esp{el producto} (produit) & \esp{consumir} (consommer) & \esp{ético/a} (éthique) \\
\esp{el comercio justo} (commerce équitable) & \esp{proteger} (protéger) & \esp{sostenible} (durable) \\
\esp{el productor} (producteur) & \esp{apoyar} (soutenir) & \esp{transparente} (transparent) \\
\esp{la etiqueta} (étiquette) & \esp{certificar} (certifier) & \esp{ecológico/a} (écologique) \\
\hline
\end{tabularx}
\end{center}
\end{definition}

\begin{savaistu}[Le Cameroun et la Guinée équatoriale : commerce et cacao]
Le marché de \textbf{Mokolo} (Yaoundé) et celui de \textbf{Nkoulouloun} (Douala) sont des pôles majeurs où s'applique la négociation (\esp{regatear}). La Guinée équatoriale (Malabo, Bata) partage la monnaie (FCFA) et une forte culture du \textbf{cacao} (\esp{el cacao}). Le \esp{comercio justo} garantit un \esp{precio justo} aux \esp{productores} (paysans) et encourage une agriculture \esp{sostenible}. Au Cameroun, les labels \og Made in Cameroon \fg et les coopératives de café/cacao illustrent cette éthique locale.
\end{savaistu}

\courssub{Grammaire 1 : Exprimer la cause (\textit{porque})}

\begin{propriete}[La cause : \textit{porque}, \textit{como}, \textit{puesto que}]
\begin{itemize}
\item \textbf{\esp{porque}} (un mot, sans accent) : introduit la cause réelle, réponse à \esp{¿por qué?}. Suivi de l'\textbf{indicatif} (fait certain). \esp{No compro \textbf{porque} es caro.}
\item \textbf{\esp{como}} / \textbf{\esp{puesto que}} (registre plus formel) : cause connue ou évidente, souvent en tête de phrase. \esp{\textbf{Como} llueve, no voy al mercado.}
\item \textbf{Attention :} \esp{por qué} (deux mots, accent) = « pourquoi » (question) ; \esp{por que} (deux mots, sans accent) = « pour lequel » (préposition + pronom relatif) ; \esp{el porqué} (nom masculin) = « la raison ».
\end{itemize}
\end{propriete}

\begin{exemplebox}[Cause]
\esp{Compro productos locales \textbf{porque} son más baratos.} (J'achète des produits locaux \textbf{parce qu'}ils sont moins chers.) \\
\esp{\textbf{Como} no tengo dinero, no puedo comprar.} (\textbf{Comme} je n'ai pas d'argent, je ne peux pas acheter.)
\end{exemplebox}

\courssub{Grammaire 2 : Exprimer le but (\textit{para}) et la condition (\textit{si})}

\begin{propriete}[Le but : \textit{para} + infinitif]
\esp{para} + \textbf{infinitif} exprime la finalité, l'intention (réponse à \esp{¿para qué?}). Sujet identique dans les deux propositions.
\esp{Ahorro dinero \textbf{para} comprar un teléfono.} (J'économise \textbf{pour} acheter un téléphone.)
\emph{Si sujets différents :} \esp{para que} + \textbf{subjonctif}. \esp{Te lo explico \textbf{para que} lo entiendas.}
\end{propriete}

\begin{propriete}[La condition : \textit{si} + indicatif présent]
Structure : \esp{Si} + \textbf{présent de l'indicatif} $\rightarrow$ \textbf{Présent} (vérité générale/habitude) ou \textbf{Futur simple} (condition réelle future) ou \textbf{Impératif} (ordre/conseil).
\emph{Jamais de futur ni de conditionnel après \esp{si} (si conditionnel).}
\begin{center}
\begin{tabular}{|l|l|l|}
\hline
\rowcolor{popblueL} \textbf{Type} & \textbf{Proposition \esp{si}} & \textbf{Proposition principale} \\
\hline
Vérité générale & \esp{Si llueve, \textbf{no voy}} (présent) & \textbf{Présent} \\
\hline
Condition réelle future & \esp{Si tengo tiempo, \textbf{iré}} (présent) & \textbf{Futur simple} \\
\hline
Ordre / Conseil & \esp{Si ves a Juan, \textbf{dile}} (présent) & \textbf{Impératif} \\
\hline
\end{tabular}
\end{center}
\end{propriete}

\begin{exemplebox}[But et Condition]
\esp{Trabajo \textbf{para} mantener a mi familia.} (But : même sujet) \\
\esp{\textbf{Si} los precios bajan, \textbf{compro} más.} (Habitude : présent + présent) \\
\esp{\textbf{Si} me baja el precio, \textbf{me lo llevo}.} (Futur réel : présent + futur) \\
\esp{\textbf{Si} quiere calidad, \textbf{cómprelo} aquí.} (Conseil : présent + impératif)
\end{exemplebox}

\begin{popfigure}
\begin{center}\tcbox[colback=popink,colframe=popink,coltext=white,arc=9pt,boxsep=4pt,left=6pt,right=6pt]{\bfseries\small Conectores lógicos}\end{center}
\vspace{-2pt}
\begin{tcbraster}[raster columns=2,raster equal height=rows,raster column skip=6pt,raster row skip=6pt,enhanced,blankest]
\begin{tcolorbox}[enhanced,colback=white,colframe=popgreen,boxrule=0.9pt,arc=3pt,title={CAUSA \esp{porque} + ind. \esp{como} + ind.},fonttitle=\bfseries\scriptsize,coltitle=white,colbacktitle=popgreen,left=4pt,right=4pt,top=2pt,bottom=2pt,boxsep=1pt,toptitle=1.5pt,bottomtitle=1.5pt,halign=left]
\begin{itemize}[nosep,leftmargin=*,itemsep=1pt,font=\scriptsize]
\item Ex: \esp{No voy \textbf{porque} llueve}
\end{itemize}
\end{tcolorbox}
\begin{tcolorbox}[enhanced,colback=white,colframe=poporange,boxrule=0.9pt,arc=3pt,title={BUT \esp{para} + inf. \esp{para que} + subj.},fonttitle=\bfseries\scriptsize,coltitle=white,colbacktitle=poporange,left=4pt,right=4pt,top=2pt,bottom=2pt,boxsep=1pt,toptitle=1.5pt,bottomtitle=1.5pt,halign=left]
\begin{itemize}[nosep,leftmargin=*,itemsep=1pt,font=\scriptsize]
\item Ex: \esp{Estudio \textbf{para} aprobar}
\end{itemize}
\end{tcolorbox}
\begin{tcolorbox}[enhanced,colback=white,colframe=poppink,boxrule=0.9pt,arc=3pt,title={CONDITION \esp{si} + prés. ind. $\rightarrow$ prés / futur / imp.},fonttitle=\bfseries\scriptsize,coltitle=white,colbacktitle=poppink,left=4pt,right=4pt,top=2pt,bottom=2pt,boxsep=1pt,toptitle=1.5pt,bottomtitle=1.5pt,halign=left]
\begin{itemize}[nosep,leftmargin=*,itemsep=1pt,font=\scriptsize]
\item Ex: \esp{\textbf{Si} tengo, \textbf{iré}}
\end{itemize}
\end{tcolorbox}
\end{tcbraster}

\legende{Carte mentale : Cause, But, Condition}
\end{popfigure}

\courssub{Savoir-faire 1 : Jouer et comprendre un dialogue client-vendeur}

\begin{methode}[Réussir un dialogue commercial]
\begin{enumerate}
\item \textbf{Identifier la situation} : qui parle (\esp{cliente}, \esp{vendedor}), où (\esp{tienda}, \esp{mercado}), pour quoi faire (\esp{comprar}, \esp{devolver}, \esp{regatear}).
\item \textbf{Mémoriser les formules d'ouverture} : \esp{Buenos días, ¿qué desea?} / \esp{Quisiera...} / \esp{Estoy buscando...}.
\item \textbf{Structurer l'échange en 4 temps} : saludo, petición, negociación del precio, despedida y agradecimiento.
\item \textbf{Utiliser le conditionnel de politesse} : \esp{¿Podría hacerme un descuento?} « Pourriez-vous me faire une remise ? » ; \esp{Me gustaría...} « J'aimerais... ».
\item \textbf{Justifier avec la cause, le but et la condition} : \esp{Lo compro \textbf{porque} es de buena calidad}, \esp{lo quiero \textbf{para} regalar}, \esp{\textbf{Si} me baja el precio, me lo llevo}.
\item \textbf{Conclure poliment} : \esp{Muchas gracias, hasta luego.}
\end{enumerate}
\end{methode}

\begin{exoresolu}[Dialogue au marché de Mokolo (complétion)]
\textbf{Énoncé.} Complétez ce dialogue avec les expressions convenables (\esp{porque}, \esp{para}, \esp{si}, \esp{quisiera}, \esp{¿cuánto cuesta?}), puis traduisez les répliques 1 et 4 en français.

--- Vendedor : Buenos días, señora, ¿qué desea?

--- Cliente (1) : \underline{\hspace{3cm}} dos kilos de tomates, por favor.

--- Vendedor : Aquí tiene. Son muy frescos.

--- Cliente (2) : \underline{\hspace{3cm}} el kilo ?

--- Vendedor : Quinientos francos. Pero (3) \underline{\hspace{3cm}} lleva tres kilos, le dejo el kilo a cuatrocientos.

--- Cliente (4) : De acuerdo, los compro \underline{\hspace{2cm}} quiero hacer una salsa \underline{\hspace{2cm}} toda la familia.

\tcblower
\textbf{Solution détaillée.}

(1) \esp{Quisiera} : conditionnel de politesse de \esp{querer}, indispensable pour formuler une demande sans brutalité. Traduction : « Bonjour madame, que désirez-vous ? — \textbf{Je voudrais} deux kilos de tomates, s'il vous plaît. »

(2) \esp{¿Cuánto cuesta} : question de prix ; attention à l'accent interrogatif sur \esp{cuánto} et au \esp{¿} initial.

(3) \esp{si} : conditionnel avec présent de l'indicatif : \esp{si lleva tres kilos} « si vous prenez trois kilos ». \textbf{Jamais de futur après \esp{si} conditionnelle.}

(4) \esp{porque} (cause : la raison de l'achat) puis \esp{para} (but : destination de la sauce). Traduction : « D'accord, je les achète \textbf{parce que} je veux faire une sauce \textbf{pour} toute la famille. »

\textbf{Conclusion} : un bon dialogue commercial alterne politesse (conditionnel), question de prix et justification (cause, but, condition).
\end{exoresolu}

\begin{exercice}[Entraînement : Dialogue à la tienda]
\textbf{Consigne.} Complétez le dialogue suivant en conjuguant les verbes entre parenthèses et en choisissant le connecteur logique (\esp{porque}, \esp{para}, \esp{si}).

--- Vendedor : Buenas tardes. ¿En qué le puedo ayudar?

--- Cliente : Busco una camisa de algodón. (1) \underline{\hspace{2cm}} (\esp{querer} - cond.) una talla M.

--- Vendedor : Tenemos esta azul. Cuesta veinte mil francos.

--- Cliente : ¿(2) \underline{\hspace{2cm}} (\esp{poder} - cond.) hacerme un descuento? (3) \underline{\hspace{2cm}} (\esp{ser}) un cliente habitual.

--- Vendedor : (4) \underline{\hspace{2cm}} (\esp{llevar} - prés. 2sg) dos, le hago un 10\%.

--- Cliente : Vale, (5) \underline{\hspace{2cm}} (\esp{llevar} - fut. 1sg) las dos (6) \underline{\hspace{2cm}} (\esp{necesitar} - inf.) para un viaje.
\end{exercice}

\begin{corrige}[Exercice 1]
(1) \esp{Quisiera} (politesse). (2) \esp{¿Podría} (politesse). (3) \esp{porque} (cause : je suis client habituel). (4) \esp{Si lleva} (condition réelle : présent). (5) \esp{llevaré} (futur dans la principale). (6) \esp{para} (but : même sujet).
\textbf{Traduction :} « Bonsoir. Je voudrais une taille M. Pourriez-vous me faire une remise car je suis client habituel ? Si vous en prenez deux, je fais 10\%. D'accord, j'en prendrai deux pour un voyage. »
\end{corrige}

\courssub{Savoir-faire 2 : Rédiger un texte de 80 mots sur ses habitudes de consommation}

\begin{methode}[Rédiger « Mis hábitos de consumo »]
\begin{enumerate}
\item \textbf{Compter les mots} : 80 mots = environ 8 à 10 phrases simples. Ni moins (hors sujet de quantité), ni beaucoup plus.
\item \textbf{Plan en 3 parties} : présentation (\esp{dónde y cuándo compro}), habitudes (\esp{qué compro y por qué}), opinion éthique (\esp{comercio justo, productos locales}).
\item \textbf{Varier les connecteurs} : \esp{normalmente}, \esp{además}, \esp{sin embargo}, \esp{por eso}, \esp{por tanto}.
\item \textbf{Utiliser les trois structures de la leçon} : au moins une cause (\esp{porque}), un but (\esp{para} + infinitif), une condition (\esp{si} + présent).
\item \textbf{Employer le présent de l'indicatif} (habitudes) et la première personne : \esp{suelo comprar} « j'ai l'habitude d'acheter » (\esp{soler} + infinitif).
\item \textbf{Relire} : accents, accords genre/nombre, pas de calque du français (\esp{barato} pas \esp{bon marché}).
\end{enumerate}
\end{methode}

\begin{exoresolu}[Rédaction guidée : mes habitudes de consommation]
\textbf{Énoncé.} Rédigez un texte d'environ 80 mots intitulé \esp{Mis hábitos de consumo} : vous direz où et quand vous faites vos achats, ce que vous achetez et pourquoi, et vous donnerez votre avis sur la consommation responsable.

\tcblower
\textbf{Raisonnement} : plan en 3 parties ; présent de l'indicatif ; intégrer \esp{porque}, \esp{para}, \esp{si} ; viser 80 mots.

\textbf{Texte rédigé (82 mots) :}

\esp{Normalmente compro en el mercado de mi barrio, en Yaoundé, los sábados por la mañana. Compro frutas y verduras locales porque son frescas y baratas. A veces voy a la tienda para comprar arroz y aceite. Si tengo bastante dinero, prefiero los productos del comercio justo, porque ayudan a los campesinos. Creo que debemos consumir de manera responsable para proteger el medio ambiente y apoyar la economía de nuestro país.}

« Normalement, j'achète au marché de mon quartier, à Yaoundé, le samedi matin. J'achète des fruits et légumes locaux parce qu'ils sont frais et bon marché. Parfois je vais au magasin pour acheter du riz et de l'huile. Si j'ai assez d'argent, je préfère les produits du commerce équitable, parce qu'ils aident les paysans. Je pense que nous devons consommer de manière responsable pour protéger l'environnement et soutenir l'économie de notre pays. »

\textbf{Analyse linguistique} : \esp{Normalmente} (marqueur d'habitude) ; \esp{porque son frescas} (cause) ; \esp{para comprar} (but) ; \esp{Si tengo} (condition réelle) ; \esp{prefiero} (présent irrégulier \esp{preferir} e$\rightarrow$ie) ; \esp{debemos consumir} (obligation morale \esp{deber} + inf.) ; \esp{para proteger / apoyar} (buts).
\end{exoresolu}

\begin{exercice}[Production écrite : Mes habitudes]
\textbf{Consigne.} Rédigez votre propre texte \esp{Mis hábitos de consumo} (80 mots $\pm$ 10). Vous êtes élève à Douala ou Yaoundé. Utilisez \textbf{au moins} : \esp{porque}, \esp{para} + inf., \esp{si} + prés., \esp{suelo} + inf., \esp{comercio justo} ou \esp{productos locales}.
\end{exercice}

\begin{corrige}[Exercice 2 - Proposition de correction type]
\esp{Suelo hacer la compra en el supermercado cerca de mi casa en Douala los viernes por la tarde. Compro productos locales porque apoyan a los agricultores de mi país. Para la ropa, voy al mercado central si hay rebajas. Si encuentro ropa de segunda mano en buen estado, la compro para reducir residuos. Pienso que el consumo responsable es importante para el futuro de nuestro planeta.}
(81 mots. Structures : \esp{suelo hacer}, \esp{porque}, \esp{para} + inf., \esp{si} + prés. x2, vocabulaire éthique.)
\end{corrige}

\courssub{Savoir-faire 3 : Traduire en espagnol (thème) des phrases sur le commerce}

\begin{methode}[Réussir le thème sur la consommation]
\begin{enumerate}
\item \textbf{Repérer le temps français} : présent d'habitude = présent de l'indicatif espagnol.
\item \textbf{Traduire « parce que » par \esp{porque}} (un mot, sans accent). Ne pas écrire \esp{por que} (deux mots) pour la conjonction de cause.
\item \textbf{Traduire « pour » + infinitif de but par \esp{para} + infinitif} : « pour économiser » = \esp{para ahorrar}.
\item \textbf{Traduire « si » + présent par \esp{si} + présent} ; attention, « si » interrogatif indirect = \esp{si} aussi, mais sans valeur conditionnelle.
\item \textbf{Choisir le bon vocabulaire} : « bon marché » = \esp{barato} (pas \esp{bon marché} calqué) ; « faire des achats » = \esp{hacer compras} / \esp{hacer la compra} ; « un client » = \esp{un cliente} ; « honnête » = \esp{honrado} / \esp{honesto}.
\item \textbf{Vérifier les accords} : \esp{los productos locales}, \esp{una tienda pequeña}, \esp{clientes habituales}.
\item \textbf{Préposition \esp{a} personnelle} : devant personne animée objet direct (\esp{engañar a los clientes}, \esp{ayudar a los productores}).
\end{enumerate}
\end{methode}

\begin{exoresolu}[Thème : le commerce équitable]
\textbf{Énoncé.} Traduisez en espagnol :

1. J'achète des produits locaux parce qu'ils sont moins chers.

2. Si les prix baissent, les clients achètent plus.

3. Nous travaillons pour aider les petits producteurs.

4. Ce vendeur est honnête : il ne trompe jamais ses clients.

\tcblower
\textbf{Solutions justifiées.}

1. \esp{Compro productos locales \textbf{porque} son más baratos.} — \esp{porque} + indicatif (cause réelle) ; comparatif \esp{más baratos} (accord masc. pluriel).

2. \esp{\textbf{Si} los precios bajan, los clientes compran más.} — \esp{si} + présent, présent dans la principale (vérité générale/habitude) ; \textbf{jamais de futur après \esp{si}}.

3. \esp{Trabajamos \textbf{para} ayudar \textbf{a} los pequeños productores.} — \esp{para} + infinitif (but) ; \esp{a} personnelle devant \esp{los pequeños productores} (personnes).

4. \esp{Este vendedor \textbf{es honrado}: nunca engaña \textbf{a} sus clientes.} — \esp{ser} (qualité permanente) et non \esp{estar} ; \esp{engañar a} avec \esp{a} personnelle ; « honnête » = \esp{honrado/honesto} (h muet, pas d'accent).

\textbf{Conclusion} : le thème exige la maîtrise conjointe du vocabulaire commercial et des trois connecteurs de la leçon.
\end{exoresolu}

\begin{exercice}[Thème : Entraînement au Bac]
\textbf{Consigne.} Traduisez en espagnol.

1. Nous achetons du café équitable pour soutenir les producteurs.

2. Si tu vas au marché, achète des bananes plantains.

3. Ce produit n'est pas cher parce qu'il est de saison.

4. Les clients respectent ce commerçant car il ne ment jamais.
\end{exercice}

\begin{corrige}[Exercice 3]
1. \esp{Compramos café de comercio justo \textbf{para} apoyar a los productores.} (\esp{para} + inf. but ; \esp{a} pers.)
2. \esp{\textbf{Si} vas al mercado, \textbf{cómprate} plátanos.} (Condition + impératif affirmatif \esp{cómprate} pronominal).
3. \esp{Este producto no es caro \textbf{porque} es de temporada.} (Cause \esp{porque} + ind. ; \esp{de temporada} = de saison).
4. \esp{Los clientes respetan a este comerciante \textbf{porque} nunca miente.} (\esp{respetar a} pers. ; \esp{mentir} e$\rightarrow$ie : \esp{miente}).
\end{corrige}

\courssub{Savoir-faire 4 : Traduire en français (version) un texte sur l'éthique}

\begin{methode}[Réussir la version sur l'éthique professionnelle]
\begin{enumerate}
\item \textbf{Lire tout le texte} avant de traduire pour saisir le sens global.
\item \textbf{Traduire les connecteurs avec précision} : \esp{porque} = « parce que » ; \esp{para} + infinitif = « pour » ; \esp{si} = « si » ; \esp{sin embargo} = « cependant » ; \esp{por tanto} = « par conséquent ».
\item \textbf{Éviter les faux amis} : \esp{decepción} = « déception » (pas « tromperie» = \esp{engaño}) ; \esp{actualmente} = « actuellement » (pas « actuellement » au sens de « en ce moment » seulement, mais c'est le sens principal) ; \esp{largo} = « long » (pas « large » = \esp{ancho}) ; \esp{éxito} = « succès » (pas « sortie ») ; \esp{asistir} = « assister à » (pas « aider » = \esp{ayudar}) ; \esp{una tienda} = « un magasin » (pas « une tente »).
\item \textbf{Restituer le \esp{se} passif / impersonnel} : \esp{se venden productos} = « on vend des produits » / « des produits sont vendus ».
\item \textbf{Rédiger un français correct} : pas de mot à mot, accords et ponctuation française (pas de \esp{¿} \esp{¡}).
\end{enumerate}
\end{methode}

\begin{exoresolu}[Version : un comerciante honesto]
\textbf{Énoncé.} Traduisez en français :

\esp{Don Miguel es un comerciante de Malabo que vende cacao desde hace veinte años. Sus clientes lo respetan porque nunca engaña a nadie y siempre da el precio justo. Si un producto es de mala calidad, no lo vende, porque quiere conservar la confianza de la gente. Trabaja mucho para mantener a su familia y para pagar los estudios de sus hijos. Para él, la honestidad es más importante que el dinero.}

\tcblower
\textbf{Analyse lexicale et grammaticale} :
\begin{itemize}
\item \esp{desde hace veinte años} = « depuis vingt ans » (marqueur de durée avec présent).
\item \esp{dar el precio justo} = « donner le juste prix / pratiquer un prix juste ».
\item \esp{mantener a su familia} = « faire vivre sa famille / subvenir aux besoins de sa famille ».
\item \esp{conservar la confianza} = « conserver la confiance / garder la confiance ».
\item \esp{más importante que} = comparatif de supériorité « plus important que ».
\end{itemize}

\textbf{Traduction} :
« Don Miguel est un commerçant de Malabo qui vend du cacao depuis vingt ans. Ses clients le respectent parce qu'il ne trompe jamais personne et qu'il donne toujours le juste prix. Si un produit est de mauvaise qualité, il ne le vend pas, parce qu'il veut conserver la confiance des gens. Il travaille beaucoup pour faire vivre sa famille et pour payer les études de ses enfants. Pour lui, l'honnêteté est plus importante que l'argent. »

\textbf{Conclusion} : une bonne version restitue le sens en français naturel, sans calquer la construction espagnole.
\end{exoresolu}

\begin{exercice}[Version : Entraînement]
\textbf{Consigne.} Traduisez en français.

\esp{En mi barrio hay una tienda de comercio justo. Los productos son un poco más caros, pero la gente los compra porque saben que el dinero ayuda a las comunidades rurales. Si todos consumimos de forma responsable, podemos cambiar el mundo. Para la dueña de la tienda, la ética no es una moda, es una forma de vida.}
\end{exercice}

\begin{corrige}[Exercice 4]
« Dans mon quartier, il y a un magasin de commerce équitable. Les produits sont un peu plus chers, mais les gens les achètent parce qu'ils savent que l'argent aide les communautés rurales. Si nous consommons tous de manière responsable, nous pouvons changer le monde. Pour la propriétaire du magasin, l'éthique n'est pas une mode, c'est un mode de vie. »
\emph{Notes :} \esp{un poco más caros} (un peu plus chers) ; \esp{la gente} (singulier en espagnol, pluriel en français « les gens ») ; \esp{saben} (savent, \esp{saber}) ; \esp{la dueña} (la propriétaire/patronne) ; \esp{una forma de vida} (un mode de vie).
\end{corrige}

\courssub{Savoir-faire 5 : Transformer des phrases (cause, but, condition)}

\begin{methode}[Transformer avec \textit{porque}, \textit{para}, \textit{si}]
\begin{enumerate}
\item \textbf{Identifier la relation logique} entre les deux idées : raison (cause), intention (but) ou hypothèse (condition).
\item \textbf{Cause} : joindre avec \esp{porque} + verbe conjugué : \esp{No compro carne. Es cara.} $\rightarrow$ \esp{No compro carne \textbf{porque} es cara.}
\item \textbf{But} : joindre avec \esp{para} + infinitif (même sujet) : \esp{Ahorro dinero. Quiero comprar un teléfono.} $\rightarrow$ \esp{Ahorro dinero \textbf{para} comprar un teléfono.}
\item \textbf{Condition} : \esp{si} + présent, puis présent/futur/impératif dans la principale : \esp{Si tengo dinero, compraré un regalo.}
\item \textbf{Vérifier} que la phrase transformée garde le sens des deux phrases de départ.
\end{enumerate}
\end{methode}

\begin{exoresolu}[Transformations logiques]
\textbf{Énoncé.} Reliez chaque paire de phrases avec \esp{porque}, \esp{para} ou \esp{si} selon le sens.

1. Voy al mercado. / Necesito verduras frescas.

2. Estudio español. / Quiero trabajar en el comercio internacional.

3. Tienes tiempo mañana. / Ven conmigo a la tienda.

\tcblower
\textbf{Solutions justifiées.}

1. \esp{Voy al mercado \textbf{porque} necesito verduras frescas.} — relation de cause : la nécessité explique le déplacement ; verbe conjugué après \esp{porque}.

2. \esp{Estudio español \textbf{para} trabajar en el comercio internacional.} — relation de but : même sujet, donc \esp{para} + infinitif (et non \esp{para que} + subjonctif, réservé aux sujets différents).

3. \esp{\textbf{Si} tienes tiempo mañana, \textbf{ven} conmigo a la tienda.} — condition : \esp{si} + présent, impératif \esp{ven} dans la principale (ordre sous condition).

\textbf{Conclusion} : le choix du connecteur dépend uniquement de la relation logique : raison $\rightarrow$ \esp{porque}, intention $\rightarrow$ \esp{para}, hypothèse $\rightarrow$ \esp{si}.
\end{exoresolu}

\begin{exercice}[Transformations : Entraînement]
\textbf{Consigne.} Reliez les phrases en utilisant le connecteur logique approprié (\esp{porque}, \esp{para}, \esp{si}).

1. No como carne. / Soy vegetariano.

2. Ahorro dinero. / Quiero comprar una moto.

3. Llueve mañana. / No voy al mercado.

4. Estudio mucho. / Quiero sacar buenas notas.

5. Tienes hambre. / Come algo.
\end{exercice}

\begin{corrige}[Exercice 5]
1. \esp{No como carne \textbf{porque} soy vegetariano.} (Cause)
2. \esp{Ahorro dinero \textbf{para} comprar una moto.} (But, même sujet)
3. \esp{\textbf{Si} llueve mañana, no voy al mercado.} (Condition, présent + présent/futur)
4. \esp{Estudio mucho \textbf{para} sacar buenas notas.} (But, même sujet)
5. \esp{\textbf{Si} tienes hambre, \textbf{come} algo.} (Condition + impératif)
\end{corrige}

\begin{attention}[Les 5 erreurs qui coûtent des points au Bac]
\begin{enumerate}
\item \textbf{Oublier l'accent interrogatif et le signe d'ouverture} : écrire \esp{cuanto cuesta} au lieu de \esp{¿Cuánto cuesta?} ; sans accent sur \esp{cuánto}, \esp{dónde}, \esp{cómo}, \esp{qué}, \esp{quién}... la phrase est fautive. N'oubliez jamais le \esp{¿} initial et le \esp{?} final.
\item \textbf{Mettre un futur ou un conditionnel après \esp{si}} : \esp{si tendré dinero} ou \esp{si tendría dinero} sont \textbf{faux} ; on dit \esp{Si tengo dinero, compraré...} « Si j'ai de l'argent, j'achèterai... ».
\item \textbf{Calquer « pour » par \esp{por}} : « je travaille pour gagner ma vie » = \esp{trabajo \textbf{para} ganar la vida} ; \esp{por} exprime la cause (\esp{por eso}), le moyen (\esp{por teléfono}) ou la durée, \esp{para} exprime le but, la destination, l'échéance.
\item \textbf{Confondre les faux amis fréquents} :
\begin{itemize}
\item \esp{asistir} = « assister à » (pas « aider » = \esp{ayudar})
\item \esp{una tienda} = « un magasin » (pas « une tente » = \esp{una tienda de campaña})
\item \esp{engañar} = « tromper » (pas « engager » = \esp{contratar})
\item \esp{el éxito} = « le succès » (pas « la sortie » = \esp{la salida})
\item \esp{actualmente} = « actuellement » (pas « actuellement » $\neq$ \esp{actualmente} est correct, mais \esp{realmente} = « vraiment »)
\item \esp{la decepción} = « la déception » (pas « la tromperie » = \esp{el engaño})
\end{itemize}
\item \textbf{Négliger les accords} : \esp{los productos locales}, \esp{una cliente satisfecha}, \esp{precios justos} ; l'adjectif s'accorde en genre et en nombre avec son nom, y compris dans la rédaction de 80 mots. Vérifiez \esp{barato/barata}, \esp{caro/cara}, \esp{fresco/fresca}.
\end{enumerate}
\end{attention}

\begin{experience}[Jeu de rôle : La négociation au marché]
\textbf{Objectif} : Réinvestir le lexique commercial, les connecteurs (\esp{porque}, \esp{para}, \esp{si}) et le conditionnel de politesse à l'oral.
\textbf{Déroulement} (15 min) :
\begin{enumerate}
\item Par binômes : un élève \textbf{Vendedor} (a une liste de prix : tomates 500, oignons 300, bananes 400 FCFA/kg), un élève \textbf{Cliente} (a une liste de courses et un budget max 2000 FCFA).
\item La \textbf{Cliente} doit acheter 3 produits différents, négocier le prix (\esp{¿Me hace descuento?}, \esp{Si llevo más, ¿me baja el precio?}), justifier ses choix (\esp{porque son frescos}, \esp{para cocinar hoy}).
\item Le \textbf{Vendedor} argumente la qualité (\esp{son de la finca}, \esp{es precio justo}) et propose des lots (\esp{Si lleva 3 kilos...}).
\item Changement de rôles. Quelques binômes jouent devant la classe.
\end{enumerate}
\textbf{Évaluation} : Fluidité, correction du \esp{si} + présent, usage de \esp{porque/para}, politesse (\esp{quisiera/podría}).
\end{experience}

\begin{savaistu}[Éthique et labels : du local au global]
Au Cameroun, le label \og \textbf{Produit du Cameroun} \fg et les initiatives \og \textbf{Commerce Équitable Cameroun} \fg certifient le cacao, le café, le coton. En Espagne, le label \esp{Comercio Justo} (ONG \esp{IDEAS}, \esp{Intermón Oxfam}) garantit : prix minimum vital, préfinancement, pas de travail d'enfants, respect de l'environnement. En Amérique latine, les coopératives de \esp{café}, \esp{banane}, \esp{quinoa} (Bolivie, Pérou, Colombie) exportent sous label \esp{Fairtrade / Max Havelaar}. Acheter \esp{local} et \esp{justo} réduit l'empreinte carbone (\esp{huella de carbono}) et soutient l'\esp{economía solidaria}.
\end{savaistu}

\newpage

\coursec{Exercices}

\courssub{Je vérifie mes connaissances}

\begin{exercice}[QCM : cause, but, condition]
Pour chaque phrase, choisis la bonne réponse (a, b ou c).
\begin{enumerate}
\item \esp{Compro en este mercado \cle{porque} los precios son justos.} — \esp{porque} exprime : a) le but ; b) la cause ; c) la condition.
\item \esp{Ahorro dinero \cle{para} comprar una bicicleta.} — \esp{para} exprime : a) la cause ; b) la conséquence ; c) le but.
\item \esp{\cle{Si} bajas el precio, compro dos camisas.} — \esp{si} exprime : a) la condition ; b) la cause ; c) le but.
\item \esp{Trabaja mucho \cle{para que} sus hijos estudien.} — Après \esp{para que}, on emploie : a) l'infinitif ; b) l'indicatif ; c) le subjonctif.
\item \esp{No fui al mercado \cle{a causa de} la lluvia.} — \esp{a causa de} est suivi de : a) un nom ; b) un verbe conjugué ; c) un infinitif.
\end{enumerate}
\end{exercice}

\begin{exercice}[Vrai ou faux justifié]
Dis si les affirmations suivantes sont vraies (V) ou fausses (F), puis justifie ta réponse en une phrase en français.
\begin{enumerate}
\item On peut dire \esp{¿Por qué no compras nada?} pour demander la cause.
\item \esp{Porque} (un seul mot, sans accent) sert à poser une question.
\item Après \esp{para} + infinitif, le sujet de l'infinitif est le même que celui du verbe principal.
\item \esp{Si} (condition) prend toujours un accent : \esp{sí}.
\item \esp{Gracias a} exprime une cause positive, tandis que \esp{a causa de} exprime souvent une cause négative.
\end{enumerate}
\end{exercice}

\begin{exercice}[Vocabulaire : le commerce et l'éthique]
Trouve l'équivalent espagnol des mots français suivants, puis utilise chacun d'eux dans une courte phrase espagnole.
\begin{enumerate}
\item un commerçant honnête
\item un prix juste
\item une réduction
\item économiser (de l'argent)
\item le commerce équitable
\item un client / une cliente
\item la qualité
\item tromper le client
\end{enumerate}
\end{exercice}

\begin{attention}[Faux amis et pièges de traduction]
\begin{itemize}
\item « honnête » se dit \esp{honrado} ou \esp{honesto} ; attention, \esp{honesto} peut aussi signifier « décent, pudique » selon le contexte.
\item « tromper » (un client) se dit \esp{engañar} ; \esp{traicionar} signifie « trahir ».
\item « économiser » se dit \esp{ahorrar} ; ne dis jamais \esp{economizar} au sens courant, et surtout pas \esp{salvar} (« sauver » une personne).
\item « une réduction » se dit \esp{un descuento} ou \esp{una rebaja} ; \esp{una reducción} est un faux ami qui signifie « une diminution » (abstraite).
\end{itemize}
\end{attention}

\begin{exercice}[Conjugaison à compléter]
Complète les phrases avec le verbe entre parenthèses au subjonctif présent ou à l'infinitif, selon le cas.
\begin{enumerate}
\item Ahorro dinero para \dots{} (comprar) un teléfono.
\item Mi padre trabaja para que nosotros \dots{} (estudiar) en buenas escuelas.
\item Te doy un descuento para que \dots{} (volver) pronto.
\item Vengo al mercado para \dots{} (vender) mi cacao.
\item El vendedor baja el precio para que los clientes \dots{} (estar) contentos.
\item Si \dots{} (tener, tú) dinero, compra productos de calidad.
\end{enumerate}
\end{exercice}

\courssub{Je m'entraîne}

\begin{exercice}[Complétion : les connecteurs]
Complète les dix phrases suivantes avec : \esp{porque}, \esp{¿por qué?}, \esp{para}, \esp{para que}, \esp{si}, \esp{a causa de} ou \esp{gracias a}.
\begin{enumerate}
\item \dots{} no compras en esta tienda? — \dots{} es muy cara.
\item Trabajo los sábados \dots{} ganar más dinero.
\item El mercado está cerrado \dots{} la fiesta nacional.
\item Te presto mi moto \dots{} llegues rápido al mercado.
\item \dots{} haces una rebaja, llevo tres sacos de arroz.
\item \dots{} su honestidad, este comerciante tiene muchos clientes.
\item Compro productos locales \dots{} ayudar a los campesinos.
\item La vendedora habla despacio \dots{} el cliente entienda bien.
\item \dots{} llueve mañana, no voy a Douala.
\item No pudo abrir su tienda \dots{} una enfermedad.
\end{enumerate}
\end{exercice}

\begin{exercice}[Transformation : para / para que]
Réécris les phrases selon la consigne : transforme \esp{para} + infinitif en \esp{para que} + subjonctif avec le sujet imposé, ou l'inverse.
\begin{enumerate}
\item Ahorro para comprar un ordenador. (sujet imposé : \esp{mi hermano})
\item El vendedor baja el precio para que yo compre más. (même sujet)
\item Trabajo duro para que mis hijos vayan a la universidad. (même sujet)
\item Vengo temprano para elegir los mejores productos. (sujet imposé : \esp{los clientes})
\item Ella estudia español para que sus padres estén orgullosos. (même sujet)
\end{enumerate}
\end{exercice}

\begin{textebox}[En el mercado Mokolo]
— Buenos días, señora. ¿Cuánto cuestan estos aguacates?\\
— Mil francos el montón, señorita.\\
— ¡Es caro! \dots{} usted me hace un \dots{}, compro dos montones.\\
— Está bien, mil ochocientos, \dots{} quiero que mis \dots{} vuelvan.\\
— Gracias. Compro aquí \dots{} usted es una vendedora \dots{} y la \dots{} de sus productos es buena.\\
— Trabajo así \dots{} tener una buena reputación en el mercado.
\end{textebox}

\begin{exercice}[Texte à trous : au marché de Yaoundé]
Complète le dialogue ci-dessus avec les mots de la liste : \esp{descuento}, \esp{calidad}, \esp{porque}, \esp{para}, \esp{si}, \esp{precio}, \esp{honesta}, \esp{clientes}. Attention : un mot de la liste est un intrus.
\end{exercice}

\begin{exercice}[Appariements : phrases à traduire]
Associe chaque phrase française (1 à 6) à sa traduction espagnole (a à f), puis recopie les paires correctes.
\begin{enumerate}
\item J'achète chez ce vendeur parce qu'il est honnête.
\item J'économise pour acheter un ordinateur.
\item Si tu fais une réduction, je prends deux sacs.
\item Il travaille dur pour que ses enfants réussissent.
\item Grâce à la pluie, la récolte de cacao est bonne.
\item Pourquoi ne viens-tu pas au marché ?
\end{enumerate}
\begin{itemize}
\item[a.] \esp{Si haces una rebaja, llevo dos sacos.}
\item[b.] \esp{¿Por qué no vienes al mercado?}
\item[c.] \esp{Compro a este vendedor porque es honrado.}
\item[d.] \esp{Trabaja duro para que sus hijos tengan éxito.}
\item[e.] \esp{Ahorro para comprar un ordenador.}
\item[f.] \esp{Gracias a la lluvia, la cosecha de cacao es buena.}
\end{itemize}
\end{exercice}

\courssub{Je me prépare au Bac}

\begin{textebox}[Un comerciante ejemplar de Yaoundé]
Don Samuel tiene una pequeña tienda en el barrio de Briqueterie, en Yaoundé. Desde hace veinte años, vende arroz, aceite y jabón a precios justos. Nunca engaña a sus clientes porque cree que la honestidad es la base del comercio. « Si un cliente no tiene dinero hoy, le fío y paga mañana », explica con una sonrisa.\\
Cada mañana, don Samuel abre su tienda a las seis para que los trabajadores puedan comprar antes de ir al trabajo. Ahorra una parte de sus ganancias para enviar a sus tres hijos a la universidad. Gracias a su buena reputación, mucha gente del barrio compra solamente en su tienda.\\
« El dinero es importante, dice, pero la confianza de la gente vale más que el dinero. » Los jóvenes comerciantes del mercado vienen a verlo para aprender su manera de trabajar. Don Samuel les repite siempre el mismo consejo : sé honesto, respeta al cliente y trabaja con paciencia.
\end{textebox}

\textbf{Glossaire} : \esp{engañar} = tromper ; \esp{fiar} = faire crédit ; \esp{las ganancias} = les bénéfices ; \esp{la confianza} = la confiance ; \esp{valer más que} = valoir plus que ; \esp{el consejo} = le conseil.

\begin{exercice}[Compréhension et étude de langue]
Lis le texte \emph{Un comerciante ejemplar de Yaoundé}, puis réponds aux questions.

\textbf{Questions de compréhension} (réponds en espagnol, avec des phrases complètes) :
\begin{enumerate}
\item ¿Qué vende don Samuel y dónde está su tienda ?
\item ¿Por qué no engaña nunca a sus clientes ?
\item ¿Para qué abre la tienda a las seis de la mañana ?
\item ¿Qué hace don Samuel con una parte de sus ganancias ?
\item ¿Qué consejo da a los jóvenes comerciantes ?
\end{enumerate}

\textbf{Questions de langue} :
\begin{enumerate}
\item Relève dans le texte deux connecteurs de cause et un connecteur de but, et traduis les phrases concernées en français.
\item Dans la phrase \esp{Abre su tienda a las seis para que los trabajadores puedan comprar}, pourquoi emploie-t-on le subjonctif ? Justifie la règle.
\item Conjugue les verbes \esp{vender} et \esp{abrir} au présent de l'indicatif (toutes les personnes).
\end{enumerate}
\end{exercice}

\begin{exercice}[Thème et version type Bac]
\textbf{Thème.} Traduis en espagnol les phrases suivantes :
\begin{enumerate}
\item Parce que ce vendeur est honnête, j'achète toujours chez lui.
\item J'économise pour acheter un ordinateur.
\item Si tu fais une réduction, je prends deux sacs.
\item Il travaille dur pour que ses enfants réussissent.
\item Grâce à ce commerçant, les habitants du quartier paient des prix justes.
\end{enumerate}

\textbf{Version.} Traduis en français le paragraphe \emph{El comercio justo} ci-dessous.
\end{exercice}

\begin{textebox}[El comercio justo]
Mama Béa vende café y cacao en el mercado central. Practica el comercio justo porque quiere ayudar a los campesinos de su región. Compra sus productos directamente a los productores para que ellos ganen más dinero. Si los clientes pagan un poco más, ella explica que es por la calidad. Gracias a su trabajo, muchas familias del campo viven mejor.
\end{textebox}

\begin{exercice}[Production écrite et jeu de rôles type Bac]
\textbf{A. Production écrite.} Rédige un texte en espagnol d'environ \textbf{80 mots} sur le sujet suivant :

\begin{itemize}
\item Sujet : \esp{Describe tus hábitos de consumo y explica por qué es importante ser un consumidor responsable.}
\item Contraintes : utilise au moins deux connecteurs de cause (\esp{porque}, \esp{gracias a}, \esp{a causa de}...), un connecteur de finalité (\esp{para}, \esp{para que}) et une phrase conditionnelle avec \esp{si}.
\end{itemize}

\textbf{B. Jeu de rôles noté.} Par groupes de deux, jouez un dialogue client-vendeur au marché (8 à 10 répliques) : le client négocie le prix d'un produit et le vendeur donne un argument éthique (commerce équitable, qualité, honnêteté). Tu seras évalué sur :
\begin{itemize}
\item la prononciation et la fluidité ;
\item l'emploi correct des connecteurs \esp{porque}, \esp{para}, \esp{para que} et \esp{si} ;
\item la richesse du lexique du commerce et de l'éthique.
\end{itemize}
\end{exercice}

\begin{methode}[Réussir la production écrite en 5 étapes]
\begin{enumerate}
\item \textbf{Analyse le sujet} : deux parties sont demandées (décrire tes habitudes + expliquer l'importance d'être responsable). Traite les deux.
\item \textbf{Liste ton vocabulaire} avant de rédiger : \esp{comprar, gastar, ahorrar, precio justo, calidad, productos locales, comercio justo}.
\item \textbf{Place les connecteurs imposés} : une cause avec \esp{porque}, une cause avec \esp{gracias a} ou \esp{a causa de}, un but avec \esp{para} ou \esp{para que}, une condition avec \esp{si} + présent.
\item \textbf{Rédige des phrases courtes} (sujet + verbe + complément) : à ton niveau, une phrase simple correcte vaut mieux qu'une phrase longue fautive.
\item \textbf{Relis} en vérifiant : accents, ¿ ¡, accord de l'adjectif (\esp{una vendedora honesta}), subjonctif après \esp{para que}, présent après \esp{si}.
\end{enumerate}
\end{methode}

\begin{exemplebox}[Modèle de production écrite (environ 80 mots)]
\esp{Normalmente compro frutas y verduras en el mercado de mi barrio porque los productos son frescos y los precios son justos. Prefiero comprar a los campesinos locales para ayudar a la economía de mi región. Si un vendedor es honrado, vuelvo siempre a su tienda. Ahorro una parte de mi dinero para comprar productos de buena calidad. Creo que es importante ser un consumidor responsable porque nuestras compras pueden cambiar la vida de los productores. Gracias al comercio justo, muchas familias viven mejor.}\\
« Normalement, j'achète des fruits et légumes au marché de mon quartier parce que les produits sont frais et les prix justes. Je préfère acheter aux paysans locaux pour aider l'économie de ma région. Si un vendeur est honnête, je retourne toujours dans sa boutique. J'économise une partie de mon argent pour acheter des produits de bonne qualité. Je crois qu'il est important d'être un consommateur responsable parce que nos achats peuvent changer la vie des producteurs. Grâce au commerce équitable, beaucoup de familles vivent mieux. »
\end{exemplebox}

\begin{aretenir}[L'essentiel de la leçon]
\begin{itemize}
\item \textbf{Cause} : \esp{porque} + verbe conjugué ; \esp{gracias a} + nom (cause positive) ; \esp{a causa de} + nom (cause négative) ; question : \esp{¿por qué?} (deux mots, accent).
\item \textbf{But} : \esp{para} + infinitif si le sujet ne change pas ; \esp{para que} + subjonctif si le sujet change.
\item \textbf{Condition} : \esp{si} + présent de l'indicatif, futur ou présent dans l'autre proposition ; jamais de subjonctif après \esp{si} ; pas d'accent sur le \esp{si} conditionnel.
\item \textbf{Lexique} : \esp{comerciante honrado, precio justo, descuento/rebaja, ahorrar, comercio justo, calidad, engañar al cliente}.
\end{itemize}
\end{aretenir}

\newpage

\coursec{Corrigés des exercices}

\begin{corrige}[Exercice 1 — QCM : cause, but, condition]
\begin{enumerate}
\item \textbf{b) la \cle{cause}.} \esp{Porque} répond à la question \esp{¿por qué?} et introduit la raison de l'achat : \esp{Compro en este mercado porque los precios son justos.} « J'achète dans ce marché parce que les prix sont justes. »
\item \textbf{c) le \cle{but}.} \esp{Para} + infinitif exprime la finalité : \esp{Ahorro dinero para comprar una bicicleta.} « J'économise de l'argent pour acheter une bicyclette. »
\item \textbf{a) la \cle{condition}.} \esp{Si} introduit l'hypothèse dont dépend l'achat : \esp{Si me haces un descuento, compro dos camisas.} « Si tu me fais une réduction, j'achète deux chemises. »
\item \textbf{c) le \cle{subjonctif}.} Après \esp{para que}, le sujet change (\esp{él} $\rightarrow$ \esp{sus hijos}), donc le subjonctif est obligatoire : \esp{estudien}.
\item \textbf{a) un \cle{nom}.} \esp{A causa de} + nom (\esp{la lluvia}) ; devant un verbe conjugué, il faut \esp{porque}.
\end{enumerate}
\textbf{Points de vigilance} : ne pas confondre \esp{porque} (réponse, un seul mot) et \esp{¿por qué?} (question, deux mots, accent sur \esp{qué}) ; retenir que \esp{para que} appelle toujours le \cle{subjonctif}.
\end{corrige}

\begin{corrige}[Exercice 2 — Vrai ou faux justifié]
\begin{enumerate}
\item \textbf{V.} \esp{¿Por qué no compras nada?} « Pourquoi n'achètes-tu rien ? » : la tournure interrogative \esp{¿por qué?} (deux mots, accent sur \esp{qué}) demande bien la \cle{cause}.
\item \textbf{F.} \esp{Porque} (un seul mot, sans accent) sert à \emph{répondre} et à donner la cause ; c'est \esp{¿por qué?} qui pose la question.
\item \textbf{V.} \esp{Para} + infinitif n'est possible que si le \cle{sujet ne change pas} : \esp{Ahorro para comprar} (c'est moi qui économise et qui achète). Si le sujet change, on emploie \esp{para que} + \cle{subjonctif}.
\item \textbf{F.} Le \esp{si} conditionnel ne prend jamais d'accent ; \esp{sí} avec accent signifie « oui » (ou « soi-même » dans \esp{él sí}).
\item \textbf{V.} \esp{Gracias a} introduit une \cle{cause favorable} (« grâce à »), \esp{a causa de} une \cle{cause défavorable} (« à cause de »).
\end{enumerate}
\end{corrige}

\begin{corrige}[Exercice 3 — Vocabulaire : le commerce et l'éthique]
\begin{enumerate}
\item \esp{un comerciante honrado} — \esp{Don Samuel es un comerciante honrado.} « Don Samuel est un commerçant honnête. »
\item \esp{un precio justo} — \esp{En esta tienda los precios son justos.} « Dans ce magasin, les prix sont justes. »
\item \esp{un descuento / una rebaja} — \esp{El vendedor me hace un descuento.} « Le vendeur me fait une réduction. »
\item \esp{ahorrar (dinero)} — \esp{Ahorro dinero para comprar una moto.} « J'économise de l'argent pour acheter une moto. »
\item \esp{el comercio justo} — \esp{El comercio justo ayuda a los campesinos.} « Le commerce équitable aide les paysans. »
\item \esp{un cliente / una clienta} — \esp{La clienta paga y se va contenta.} « La cliente paie et repart contente. »
\item \esp{la calidad} — \esp{La calidad de este cacao es excelente.} « La qualité de ce cacao est excellente. »
\item \esp{engañar al cliente} — \esp{Un comerciante honrado nunca engaña al cliente.} « Un commerçant honnête ne trompe jamais le client. »
\end{enumerate}
\textbf{Points de vigilance} : « tromper » = \esp{engañar} (pas \esp{traicionar}, « trahir ») ; « économiser » = \esp{ahorrar} (pas \esp{salvar}, « sauver ») ; \esp{una reducción} est un \cle{faux ami} (« diminution » abstraite, pas « réduction commerciale »).
\end{corrige}

\begin{corrige}[Exercice 4 — Conjugaison à compléter]
\begin{enumerate}
\item \esp{Ahorro dinero para \textbf{comprar} un teléfono.} Même sujet (\esp{yo}), donc \esp{para} + \cle{infinitif}.
\item \esp{Mi padre trabaja para que nosotros \textbf{estudiemos} en buenas escuelas.} Le sujet change (\esp{mi padre} $\rightarrow$ \esp{nosotros}), donc \esp{para que} + \cle{subjonctif}.
\item \esp{Te doy un descuento para que \textbf{vuelvas} pronto.} Le sujet change (\esp{yo} $\rightarrow$ \esp{tú}) ; \esp{volver} a un changement de radical \esp{o} $\rightarrow$ \esp{ue} au subjonctif.
\item \esp{Vengo al mercado para \textbf{vender} mi cacao.} Même sujet (\esp{yo}), donc \cle{infinitif}.
\item \esp{El vendedor baja el precio para que los clientes \textbf{estén} contentos.} \cle{Subjonctif présent} d'\esp{estar}, irrégulier et accentué : \esp{esté, estés, esté, estemos, estéis, estén}.
\item \esp{Si \textbf{tienes} dinero, compra productos de calidad.} Après \esp{si} conditionnel, \cle{présent de l'indicatif}, jamais de subjonctif ; \esp{tener} : \esp{e} $\rightarrow$ \esp{ie} (\esp{tienes}).
\end{enumerate}
\end{corrige}

\begin{corrige}[Exercice 5 — Complétion : les connecteurs]
\begin{enumerate}
\item \esp{\textbf{¿Por qué} no compras en esta tienda? — \textbf{Porque} es muy cara.} (question : deux mots + accent ; réponse : un seul mot).
\item \esp{Trabajo los sábados \textbf{para} ganar más dinero.} (\cle{but} + infinitif, même sujet).
\item \esp{El mercado está cerrado \textbf{a causa de} la fiesta nacional.} (\cle{cause négative} + nom).
\item \esp{Te presto mi moto \textbf{para que} llegues rápido al mercado.} (le sujet change : \esp{yo} $\rightarrow$ \esp{tú}, d'où le \cle{subjonctif} \esp{llegues}).
\item \esp{\textbf{Si} haces una rebaja, llevo tres sacos de arroz.} (\cle{condition} + présent de l'indicatif).
\item \esp{\textbf{Gracias a} su honestidad, este comerciante tiene muchos clientes.} (\cle{cause positive}).
\item \esp{Compro productos locales \textbf{para} ayudar a los campesinos.} (\cle{but} + infinitif, même sujet).
\item \esp{La vendedora habla despacio \textbf{para que} el cliente entienda bien.} (changement de sujet, \cle{subjonctif} \esp{entienda}, radical \esp{e} $\rightarrow$ \esp{ie}).
\item \esp{\textbf{Si} llueve mañana, no voy a Douala.} (\cle{condition}).
\item \esp{No pudo abrir su tienda \textbf{a causa de} una enfermedad.} (\cle{cause négative} + nom).
\end{enumerate}
\end{corrige}

\begin{corrige}[Exercice 6 — Transformation : para / para que]
\begin{enumerate}
\item \esp{Ahorro \textbf{para que mi hermano compre} un ordenador.} (\cle{subjonctif} de \esp{comprar} : \esp{compre}).
\item \esp{El vendedor baja el precio \textbf{para comprar} más.} (même sujet : \cle{infinitif}).
\item \esp{Trabajo duro \textbf{para ir} a la universidad.} (même sujet : \cle{infinitif}).
\item \esp{Vengo temprano \textbf{para que los clientes elijan} los mejores productos.} (\cle{subjonctif} d'\esp{elegir} : \esp{elijan}, avec changement \esp{g} $\rightarrow$ \esp{j} devant \esp{a}).
\item \esp{Ella estudia español \textbf{para estar orgullosa}.} (même sujet : \cle{infinitif} ; accord féminin \esp{orgullosa}).
\end{enumerate}
\textbf{Point de vigilance} : la transformation impose toujours le même raisonnement — sujet identique : \esp{para} + \cle{infinitif} ; sujet différent : \esp{para que} + \cle{subjonctif}.
\end{corrige}

\begin{corrige}[Exercice 7 — Texte à trous : au marché de Yaoundé]
\esp{— ¡Es caro! \textbf{Si} usted me hace un \textbf{descuento}, compro dos montones.}\par
\esp{— Está bien, mil ochocientos, \textbf{porque} quiero que mis \textbf{clientes} vuelvan.}\par
\esp{— Gracias. Compro aquí \textbf{porque} usted es una vendedora \textbf{honesta} y la \textbf{calidad} de sus productos es buena.}\par
\esp{— Trabajo así \textbf{para} tener una buena reputación en el mercado.}\par
L'intrus est \esp{precio}.\\
\textbf{Justifications} : \esp{si} + présent introduit la \cle{condition} de l'achat ; \esp{porque} donne la \cle{cause} (deux fois) ; \esp{para} + infinitif exprime le \cle{but} (même sujet) ; \esp{honesta} s'accorde au féminin avec \esp{vendedora} — on ne peut pas écrire \esp{una vendedora honesto}.
\end{corrige}

\begin{corrige}[Exercice 8 — Appariements : phrases à traduire]
1 $\rightarrow$ c ; 2 $\rightarrow$ e ; 3 $\rightarrow$ a ; 4 $\rightarrow$ d ; 5 $\rightarrow$ f ; 6 $\rightarrow$ b.
\begin{itemize}
\item \esp{Compro a este vendedor porque es honrado.} « J'achète chez ce vendeur parce qu'il est honnête. »
\item \esp{Ahorro para comprar un ordenador.} « J'économise pour acheter un ordinateur. »
\item \esp{Si haces una rebaja, llevo dos sacos.} « Si tu fais une réduction, je prends deux sacs. »
\item \esp{Trabaja duro para que sus hijos tengan éxito.} « Il travaille dur pour que ses enfants réussissent. »
\item \esp{Gracias a la lluvia, la cosecha de cacao es buena.} « Grâce à la pluie, la récolte de cacao est bonne. »
\item \esp{¿Por qué no vienes al mercado?} « Pourquoi ne viens-tu pas au marché ? »
\end{itemize}
\textbf{Point de vigilance} : « acheter chez quelqu'un » se traduit par \esp{comprar a alguien} (ou \esp{comprar en la tienda de alguien}), jamais par \esp{comprar en este vendedor}.
\end{corrige}

\begin{corrige}[Exercice 9 — Compréhension et étude de langue (type Bac)]
\textbf{Barème indicatif (sur 10)} : compréhension 5 pts (1 pt par question), langue 5 pts (2 + 1,5 + 1,5).

\textbf{Compréhension} :
\begin{enumerate}
\item \esp{Don Samuel vende arroz, aceite y jabón. Su tienda está en el barrio de Briqueterie, en Yaoundé.}
\item \esp{Nunca engaña a sus clientes porque cree que la honestidad es la base del comercio.}
\item \esp{Abre la tienda a las seis para que los trabajadores puedan comprar antes de ir al trabajo.}
\item \esp{Ahorra una parte de sus ganancias para enviar a sus tres hijos a la universidad.}
\item \esp{Les repite siempre el mismo consejo : sé honesto, respeta al cliente y trabaja con paciencia.}
\end{enumerate}
\textbf{Langue} :
\begin{enumerate}
\item \cle{Cause} : \esp{porque} — \esp{Nunca engaña a sus clientes porque cree que la honestidad es la base del comercio.} « Il ne trompe jamais ses clients parce qu'il croit que l'honnêteté est la base du commerce. » ; \esp{gracias a} — \esp{Gracias a su buena reputación, mucha gente del barrio compra solamente en su tienda.} « Grâce à sa bonne réputation, beaucoup de gens du quartier n'achètent que dans sa boutique. » \cle{But} : \esp{para que} — \esp{Abre su tienda a las seis para que los trabajadores puedan comprar.} « Il ouvre sa boutique à six heures pour que les travailleurs puissent acheter. »
\item On emploie le \cle{subjonctif} \esp{puedan} parce qu'après \esp{para que} (but), le sujet du second verbe est différent (\esp{don Samuel} $\rightarrow$ \esp{los trabajadores}) ; quand le sujet change, \esp{para que} exige toujours le subjonctif.
\item \esp{vender} : \esp{vendo, vendes, vende, vendemos, vendéis, venden} (régulier en \esp{-er}). \esp{abrir} : \esp{abro, abres, abre, abrimos, abrís, abren} (régulier en \esp{-ir} au présent ; participe passé irrégulier : \esp{abierto}).
\end{enumerate}
\textbf{Points de vigilance} : répondre par des phrases complètes en reprenant les mots de la question ; ne pas oublier l'accent de \esp{está} ni le \esp{¿} ouvrant dans une citation interrogative.
\end{corrige}

\begin{corrige}[Exercice 10 — Thème et version type Bac]
\textbf{Barème indicatif (sur 10)} : thème 5 pts (1 pt par phrase), version 5 pts (fidélité 3 pts, qualité du français 2 pts).

\textbf{Thème} :
\begin{enumerate}
\item \esp{Porque este vendedor es honrado, compro siempre en su tienda.} (ou \esp{le compro siempre a él})
\item \esp{Ahorro para comprar un ordenador.}
\item \esp{Si haces una rebaja, llevo dos sacos.}
\item \esp{Trabaja duro para que sus hijos tengan \textbf{éxito}.} (accent sur \esp{éxito})
\item \esp{Gracias a este comerciante, los habitantes del barrio pagan precios justos.}
\end{enumerate}
\textbf{Version} : « Mama Béa vend du café et du cacao au marché central. Elle pratique le commerce équitable parce qu'elle veut aider les paysans de sa région. Elle achète ses produits directement aux producteurs pour qu'ils gagnent plus d'argent. Si les clients paient un peu plus, elle explique que c'est à cause de la qualité. Grâce à son travail, beaucoup de familles de la campagne vivent mieux. »

\textbf{Points de vigilance} : « réussir » = \esp{tener \textbf{éxito}} (avec accent) ; « chez lui » = \esp{en su tienda} ou \esp{a él}, pas de calque \esp{en él} ; dans la version, \esp{por la calidad} se rend naturellement par « à cause de la qualité » ; \esp{el campo} = « la campagne », pas « le champ ».
\end{corrige}

\begin{corrige}[Exercice 11 — Production écrite et jeu de rôles type Bac]
\textbf{Barème indicatif (sur 10)} : respect du sujet et des deux parties 2 pts ; connecteurs imposés correctement employés 3 pts (\cle{cause} $\times$2, \cle{but}, \cle{condition}) ; correction grammaticale (accords, subjonctif, accents) 3 pts ; richesse du lexique 2 pts. Pour le jeu de rôles : prononciation et fluidité 4 pts, connecteurs 3 pts, lexique du commerce et de l'éthique 3 pts.

\textbf{Proposition de rédaction modèle (environ 80 mots)} :

\esp{Normalmente compro frutas y verduras en el mercado de mi barrio porque los productos son frescos y los precios son justos. Prefiero comprar a los campesinos locales para ayudar a la economía de mi región. Si un vendedor es honrado, vuelvo siempre a su tienda. Ahorro una parte de mi dinero para comprar productos de buena calidad. Creo que es importante ser un consumidor responsable porque nuestras compras pueden cambiar la vida de los productores. Gracias al comercio justo, muchas familias viven mejor.}

« Normalement, j'achète des fruits et légumes au marché de mon quartier parce que les produits sont frais et les prix justes. Je préfère acheter aux paysans locaux pour aider l'économie de ma région. Si un vendeur est honnête, je retourne toujours dans sa boutique. J'économise une partie de mon argent pour acheter des produits de bonne qualité. Je crois qu'il est important d'être un consommateur responsable parce que nos achats peuvent changer la vie des producteurs. Grâce au commerce équitable, beaucoup de familles vivent mieux. »

\textbf{Exemple de dialogue pour le jeu de rôles} :

\esp{— Buenos días, ¿cuánto cuesta este saco de arroz?}\par
\esp{— Diez mil francos, señora.}\par
\esp{— ¡Es caro! Si me hace un descuento, llevo dos sacos.}\par
\esp{— Le dejo los dos a dieciocho mil porque usted es una buena clienta.}\par
\esp{— ¿Por qué es más caro que en la otra tienda?}\par
\esp{— Porque es arroz de comercio justo : compro directamente a los campesinos para que ganen más dinero.}\par
\esp{— Entiendo. Gracias a su honestidad, tiene usted muchos clientes.}\par
\esp{— Trabajo así para tener una buena reputación.}\par
\esp{— Está bien, llevo los dos sacos. ¡Hasta pronto!}\par
\esp{— ¡Gracias, vuelva pronto!}

\textbf{Points de vigilance} :
\begin{itemize}
\item Vérifier que les deux parties du sujet sont traitées (habitudes + importance d'être responsable).
\item Après \esp{si}, \cle{présent de l'indicatif} uniquement : \esp{si un vendedor es honrado} (jamais de subjonctif, jamais de futur).
\item Après \esp{para que}, \cle{subjonctif} obligatoire : \esp{para que ganen más dinero}.
\item Accords de l'adjectif : \esp{una buena clienta}, \esp{precios justos}, \esp{productos frescos}.
\item Ne pas oublier les signes \esp{¿ ¡} et les accents : \esp{está}, \esp{región}, \esp{economía}, \esp{reputación}.
\end{itemize}
\end{corrige}

\newpage

\coursec{Fiche bilan}

\begin{popfigure}
\begin{center}\tcbox[colback=popblue,colframe=popblue,coltext=white,arc=9pt,boxsep=4pt,left=6pt,right=6pt]{\bfseries\small Consumo responsable y ética}\end{center}
\vspace{-2pt}
\begin{tcbraster}[raster columns=3,raster equal height=rows,raster column skip=6pt,raster row skip=6pt,enhanced,blankest]
\begin{tcolorbox}[enhanced,colback=poporange,colframe=poporange,coltext=white,boxrule=0.9pt,arc=3pt,left=4pt,right=4pt,top=3pt,bottom=3pt,boxsep=1pt,halign=left]\bfseries\scriptsize Diálogo cliente--vendedor\end{tcolorbox}
\begin{tcolorbox}[enhanced,colback=popgreen,colframe=popgreen,coltext=white,boxrule=0.9pt,arc=3pt,left=4pt,right=4pt,top=3pt,bottom=3pt,boxsep=1pt,halign=left]\bfseries\scriptsize Léxico del comercio y la ética\end{tcolorbox}
\begin{tcolorbox}[enhanced,colback=poppurple,colframe=poppurple,coltext=white,boxrule=0.9pt,arc=3pt,left=4pt,right=4pt,top=3pt,bottom=3pt,boxsep=1pt,halign=left]\bfseries\scriptsize Causa porque, pues, ya que, como\end{tcolorbox}
\begin{tcolorbox}[enhanced,colback=poppink,colframe=poppink,coltext=white,boxrule=0.9pt,arc=3pt,left=4pt,right=4pt,top=3pt,bottom=3pt,boxsep=1pt,halign=left]\bfseries\scriptsize Finalidad para + inf. para + nombre\end{tcolorbox}
\begin{tcolorbox}[enhanced,colback=popgold,colframe=popgold,coltext=white,boxrule=0.9pt,arc=3pt,left=4pt,right=4pt,top=3pt,bottom=3pt,boxsep=1pt,halign=left]\bfseries\scriptsize Condición si + presente, futuro/imperativo\end{tcolorbox}
\begin{tcolorbox}[enhanced,colback=popblue!70,colframe=popblue!70,coltext=white,boxrule=0.9pt,arc=3pt,left=4pt,right=4pt,top=3pt,bottom=3pt,boxsep=1pt,halign=left]\bfseries\scriptsize Traducción thème / version\end{tcolorbox}
\begin{tcolorbox}[enhanced,colback=popgreen!70!popdark,colframe=popgreen!70!popdark,coltext=white,boxrule=0.9pt,arc=3pt,left=4pt,right=4pt,top=3pt,bottom=3pt,boxsep=1pt,halign=left]\bfseries\scriptsize Producción escrita : 80 palabras\end{tcolorbox}
\end{tcbraster}

\legende{Carte mentale de la leçon : consommation responsable et éthique professionnelle}
\end{popfigure}

\begin{bilan}
\textbf{1. Lexique clé du commerce et de l'éthique}
\begin{itemize}
  \item \esp{el/la cliente} : le/la client(e) ; \esp{el/la vendedor(a)} : le/la vendeur(-euse) ; \esp{la tienda} : le magasin ; \esp{el mercado} : le marché.
  \item \esp{el precio} : le prix ; \esp{barato / caro} : bon marché / cher ; \esp{el descuento} : la remise ; \esp{la oferta} : l'offre, la promotion.
  \item \esp{comprar / vender / pagar} : acheter / vendre / payer ; \esp{en efectivo} : en espèces ; \esp{con tarjeta} : par carte.
  \item \esp{la calidad} : la qualité ; \esp{la garantía} : la garantie ; \esp{devolver} : rendre, rembourser ; \esp{el recibo} : le reçu.
  \item \esp{el consumo responsable} : la consommation responsable ; \esp{honesto/a} : honnête ; \esp{engañar} : tromper ; \esp{la publicidad engañosa} : la publicité mensongère ; \esp{respetar} : respecter.
\end{itemize}

\textbf{2. Grammaire à connaître par cœur}
\begin{center}
\begin{tabularx}{\linewidth}{|l|X|X|}
\hline
\rowcolor{popblueL} Notion & Forme & Exemple \\
\hline
Cause & \esp{porque} (répond à ¿por qué ?) & \esp{Compro aquí porque es barato.} « J'achète ici parce que c'est bon marché. » \\
\hline
Cause & \esp{como} + verbe (en tête de phrase) & \esp{Como no tengo dinero, no compro nada.} « Comme je n'ai pas d'argent, je n'achète rien. » \\
\hline
Cause & \esp{pues / ya que / puesto que} & \esp{Ya que estás aquí, mira los precios.} « Puisque tu es là, regarde les prix. » \\
\hline
But & \esp{para} + infinitif & \esp{Ahorro para comprar un teléfono.} « J'économise pour acheter un téléphone. » \\
\hline
But & \esp{para} + nom (destination) & \esp{Este regalo es para mi madre.} « Ce cadeau est pour ma mère. » \\
\hline
Condition & \esp{si} + présent, futur ou impératif & \esp{Si hay descuento, compraré más.} « S'il y a une remise, j'achèterai plus. » \\
\hline
Condition & jamais de futur après \esp{si} & \esp{Si tengo dinero, te lo doy.} « Si j'ai de l'argent, je te le donne. » \\
\hline
\end{tabularx}
\end{center}

\textbf{3. Méthodes express}
\begin{itemize}
  \item \textbf{Dialogue client--vendeur} : saluer (\esp{Buenos días}) $\rightarrow$ demander (\esp{¿Cuánto cuesta...?}) $\rightarrow$ négocier (\esp{¿Me hace un descuento?}) $\rightarrow$ conclure (\esp{Me lo llevo. Gracias.}).
  \item \textbf{Texte de 80 mots} : 3 parties (introduction : mes habitudes ; développement : où, quoi, pourquoi, avec \esp{porque}, \esp{para}, \esp{si} ; conclusion : mon opinion sur le consumo responsable).
  \item \textbf{Traduction} : repérer d'abord la logique de la phrase (cause ? but ? condition ?), choisir le bon connecteur, puis vérifier le temps du verbe.
\end{itemize}
\end{bilan}

\begin{aretenir}[Checklist avant le Bac]
\begin{itemize}
  \item $\square$ Je sais conjuguer les verbes du dialogue : \esp{querer} (quiero), \esp{poder} (puedo), \esp{costar} (cuesta), \esp{llevarse} (me lo llevo).
  \item $\square$ Je connais le vocabulaire du magasin : \esp{precio, descuento, oferta, recibo, garantía, en efectivo}.
  \item $\square$ J'emploie \esp{porque} pour répondre à \esp{¿por qué?} et je sais placer \esp{como} en tête de phrase.
  \item $\square$ Je n'oublie pas que \esp{para} + infinitif exprime le but (« pour » + infinitif).
  \item $\square$ Je ne mets jamais de futur après \esp{si} : \esp{si tengo}, pas \esp{si tendré}.
  \item $\square$ Je distingue \esp{porque} (cause) et \esp{para qué} (but) dans une question.
  \item $\square$ Je sais écrire les accents : \esp{está, compré, teléfono, ética, garantía}.
  \item $\square$ Je fais attention aux faux amis : \esp{una oferta} = une promotion (pas « une offre d'emploi ») ; \esp{asistir} = être présent.
  \item $\square$ Je structure mon texte de 80 mots en trois paragraphes avec au moins un connecteur de cause, un de but et un de condition.
  \item $\square$ Je relis ma traduction : sujet, temps, connecteur logique, accents.
  \item $\square$ Je sais saluer et remercier dans un dialogue : \esp{Buenos días, ¿en qué puedo ayudarle? / Muchas gracias, adiós.}
  \item $\square$ Je peux donner mon opinion sur la consommation responsable : \esp{Pienso que... / Me parece importante...}
\end{itemize}
\end{aretenir}

\findecours

\end{document}
